\documentclass[10pt,letterpaper]{article}
\usepackage[margin=0.88in]{geometry}
\usepackage[T1]{fontenc}
\usepackage{lmodern,microtype,amsmath,amssymb,amsthm,mathtools}
\usepackage{booktabs,longtable,array,enumitem,xurl}
\usepackage[dvipsnames]{xcolor}
\usepackage[colorlinks=true,linkcolor=MidnightBlue,urlcolor=MidnightBlue,citecolor=MidnightBlue]{hyperref}
\hypersetup{pdftitle={Residue certificates and nonnegative gates for the consistency of ZFC},
 pdfauthor={Prepared for the ProveIt project},
 pdfsubject={Explicit degree-four Diophantine consistency equations and reproducible exact algebra}}
\usepackage{fancyhdr}
\setlength{\headheight}{14pt}
\setlength{\emergencystretch}{3em}
\pagestyle{fancy}\fancyhf{}
\fancyhead[L]{\small Explicit quartics for $\operatorname{Con}(\mathrm{ZFC})$}
\fancyhead[R]{\small ProveIt research}
\fancyfoot[C]{\thepage}
\setlist{itemsep=3pt,topsep=5pt}
\newtheorem{theorem}{Theorem}[section]
\newtheorem{lemma}[theorem]{Lemma}
\newtheorem{proposition}[theorem]{Proposition}
\newtheorem{corollary}[theorem]{Corollary}
\theoremstyle{definition}\newtheorem{definition}[theorem]{Definition}
\theoremstyle{remark}\newtheorem{remark}[theorem]{Remark}
\newcommand{\N}{\mathbb N}\newcommand{\Z}{\mathbb Z}\newcommand{\Q}{\mathbb Q}
\newcommand{\ZFC}{\mathrm{ZFC}}\newcommand{\ZF}{\mathrm{ZF}}
\newcommand{\Con}{\operatorname{Con}}\newcommand{\Prf}{\operatorname{Prf}}
\newcommand{\file}[1]{\texttt{\nolinkurl{#1}}}
\newcommand{\sqd}[2]{(#1-#2)^2}
\newcommand{\Pn}{P_{407}}\newcommand{\Pz}{Q_{888}}\newcommand{\Ps}{Q_{917}}
\title{\vspace{-1em}\textbf{Residue certificates and nonnegative gates\\for the consistency of ZFC}\\[0.6em]
\large Explicit quartics in 407 nonnegative integers,\\888 unrestricted integers, and a small-coefficient 917-integer form}
\author{Prepared for Vladimir Reshetnikov's ProveIt project\\\small AI-assisted mathematical research with reproducible exact algebra}
\date{24 October 2026}
\begin{document}
\maketitle
\begin{abstract}
We supply a completely specified polynomial $\Pn\in\Z[x_0,\ldots,x_{406}]$
of exact degree four such that
\[
 \ZFC\vdash
 \left(\nexists x\in\N^{407}:\Pn(x)=0\right)\ \longleftrightarrow\ \Con(\ZFC),
 \qquad\N=\{0,1,2,\ldots\}.
\]
It is the sum of squares of 323 listed quadratic residuals. Its full
expansion has 8,420 monomials and coefficient height 83,843,774.
Two accompanying integer-domain quartics have 888 and 917 variables;
the latter has coefficient height only 63,906,036.
The construction directly encodes the calibrated numerical calculus
$C_{222}$. It does not simulate K215 and does not use Jones's universal
polynomial or Sun's variable reduction.

The main arithmetic improvement is to encode a binomial coefficient only
modulo the modulus required by the proof. A radix divisible by that modulus
removes the unnecessary exact coefficient witnesses. Only its lower digits
need to be carry-free, allowing an already required factorial numerator
to supply the entire radix bound. Nonnegative difference-square gates,
additive reuse of Pell indices, a projected eight-column proof certificate,
and certified sign inference yield the final quartics.
Against the strongest inspected incoming direct constructions, the natural
variable count falls from 522 to 407, and the integer count from 1,690 to
888. These are improvements in a stated degree/coefficient tradeoff, not
a claim of globally minimal variables or worldwide priority.

The proof is ordinary mathematics formalizable in ZFC, using the explicitly
identified inherited $C_{222}$ calibration and standard Pell theory. Exact
finite checks certify the literal polynomials, inverse substitutions, and
domain conversions. No new end-to-end Lean/Rocq proof term or object-language
ZFC derivation is claimed, and no consistency assumption is used.
\end{abstract}

\clearpage
\tableofcontents
\newpage
\section{The result and its exact scope}
\label{sec:result}

\begin{theorem}[Delivered ordinary equations]\label{thm:main}
Let $\Pn$ be the polynomial whose 323 residuals are printed in
Appendix~\ref{app:rows}, in the variable order of
\file{equations/natural407/equation.json}. Let $\Pz$ and $\Ps$ be the
polynomials in the corresponding \file{integer888} and \file{integer917}
directories. Each is the sum of the squares of its listed residuals.
Then
\begin{align}
 \ZFC&\vdash (\exists x\in\N^{407}\;\Pn(x)=0)
       \longleftrightarrow (\exists p\;\Prf_{\ZFC}(p,\bot)),\label{eq:mainN}\\
 \ZFC&\vdash (\exists z\in\Z^{888}\;\Pz(z)=0)
       \longleftrightarrow (\exists p\;\Prf_{\ZFC}(p,\bot)),\label{eq:mainZ}\\
 \ZFC&\vdash (\exists z\in\Z^{917}\;\Ps(z)=0)
       \longleftrightarrow (\exists p\;\Prf_{\ZFC}(p,\bot)).\label{eq:mainS}
\end{align}
All three polynomials have exact degree four and fixed, literal integer
coefficients. Negation gives the requested consistency equivalences.
\end{theorem}

Here $\Prf_{\ZFC}$ denotes the fixed syntactic refutation predicate used
by the repository's sealed ZFC calculus. Standard effective translations
identify its consistency sentence with the usual consistency sentence for
ZFC. Neither the coefficients nor their generation depend on deciding
whether a refutation exists.

\begin{table}[ht]
\centering\small
\begin{tabular}{lrrrrr}
\toprule
Equation/domain & Variables & Residuals & Degree & Monomials & Height\\
\midrule
$\Pn$, $\N$ & \textbf{407} & 323 & 4 & 8,420 & 83,843,774\\
Sparse companion, $\N$ & 410 & 326 & 4 & 7,029 & 211,184,624\\
$\Pz$, $\Z$ & \textbf{888} & 470 & 4 & 10,156 & 2,445,790,974,704\\
$\Ps$, $\Z$ & 917 & 499 & 4 & 6,986 & \textbf{63,906,036}\\
\bottomrule
\end{tabular}
\caption{Exact coefficient collection. Height is the maximum absolute
coefficient of the fully expanded single polynomial, not of a residual.}
\end{table}

Every term is supplied as a literal coefficient and a sorted multiset of
variable indices. There are no free program parameters, unspecified
coefficient ports, variable exponents, factorials, binomial coefficients,
or bounded quantifiers in the delivered polynomial files. Those operations
appear only in the explanatory construction and are eliminated by the
finite quadratic graph described below.

\subsection{Comparison with the actual incoming frontier}

The comparison snapshot is commit
\file{85ae4070019efd66ea3e58e9a67d287528094e8b}. The incoming ZIPs at this
commit are materially stronger than the opening summary of the numbered
manuscript collection. We inspected all 23 incoming archives' inventories
and relevant overview files, and read the arithmetic dependencies used
here. Their immutable Git blob identifiers are recorded in
\file{provenance/incoming_index.json}.

The synchronized incoming package gives a direct quartic in 522 naturals,
with 20,225 terms and height 83,900,628. The batched incoming package gives
a direct quartic in 1,690 integers, with 21,381 terms and height
2,649,643,805,758 \cite{synchronized,batched}. Relative to those respective
baselines, $\Pn$ uses 22.03\% fewer variables and 58.37\% fewer terms;
$\Pz$ uses 47.46\% fewer variables and 52.50\% fewer terms.
The 917-integer alternative lowers the coefficient height to less than
$2^{26}$ while still nearly halving the baseline variable count.

The repository also has equations with nine positive variables, ten
unrestricted variables, and machine-derived quartics with far fewer than
407 variables. Those use different degrees, arithmetic circuits, or large
fixed coefficient payloads. The incoming interpolation package additionally
contains twelve-integer forms of degree 4,832 and 16,000
\cite{repo,interpolation}. We do not assert that the present quartics
dominate those constructions. The achievement is a substantially smaller
\emph{direct, explicitly expanded quartic with moderate literal coefficients}.
There is no established scalar measure that makes one of all these forms
unconditionally the simplest.

\subsection{Attribution of the mathematical changes}

The numerical calculus, its calibration, proof-flow interpretation, beta
coding, basic arithmetic graph identities, synchronized Pell frame, and
shared auxiliary index witnesses are inherited \cite{synchronized,calibration}.
The eight-column projection was already present in the direct-projection
archive \cite{projection}. This contribution combines that projection
with the synchronized arithmetic and introduces or completes the following
specific transformations in this construction:

\begin{enumerate}
\item A range test represents divisibility of a binomial coefficient,
      rather than its complete value.
\item A lower-digit bound allows an existing factorial numerator to
      replace the independently generated common radix bound.
\item A square of a difference is one total nonnegative quadratic gate,
      avoiding the separate positive and negative expression components.
\item Requested Pell pairs at $H+2,t+2,t+1$ are derived by addition
      from pairs at $H,t$, leaving only 15 independent index tests.
\item Exact certificates permit nonnegative coordinate restorations
      with signed expanded coefficients.
\item Only 167 of the 690 source coordinates need independent sign
      guards; the other 523 signs follow from the existing equations.
\end{enumerate}
These are precise algebraic changes and measured consequences. We do not
claim that every general identity used in them is new to the literature.

\section{The calibrated numerical proof calculus}
\label{sec:calculus}

The source is the numerical calculus $C_{222}$, whose false target is
the integer 1. Label 0 is the freely available padding label. Its
machine-independent calibration is
\begin{equation}\label{eq:calibration}
 C_{222}(1)\ \longleftrightarrow\ \Prf_{\ZF}(\bot)
       \ \longleftrightarrow\ \Prf_{\ZFC}(\bot),
\end{equation}
where each proof notation abbreviates existence of a finite refutation.
The named inherited theorems are
\file{KEndpoint.c222_iff_zf_refutation} and
\file{ConstructibleChoice.zf_refutation_iff_zfc_refutation}.
Their immediate source files and pins are included under
\file{upstream/provenance/upstream/calibration/}. The second theorem uses
the constructible-universe equiconsistency of ZF and ZFC. The first
theorem occurs before the machine-specific section of its module and has
no machine-table hypothesis. Its surrounding development is also used by
machine results; that shared provenance does not make this equation a
simulation of a Turing machine.

For clarity, we give the complete arithmetic rule interface. Put
\begin{align*}
 I(a,b)&=2(a+b)(a+b+1)+4b,\\
 \mathsf M(a,b)&=I(a,b)+1,&\mathsf A(a,b)&=I(a,b)+2,&\mathsf E(a,b)&=I(a,b)+3,\\
 \mathsf N(a)&=I(a,1),&\mathsf X(a,b)&=\mathsf N(\mathsf A(a,\mathsf N(b))).
\end{align*}
The eight unconditional axiom labels are
\begin{align*}
 c_1&=\mathsf X(a,\mathsf E(a,b)),\\
 c_2&=I(\mathsf A(a,I(b,c)),I(\mathsf A(a,b),\mathsf A(a,c))),\\
 c_3&=I(\mathsf E(a,b),I(\mathsf M(c,a),\mathsf M(c,b))),\\
 c_4&=I(\mathsf A(a,\mathsf A(b,c)),\mathsf A(b,\mathsf A(a,c))),\\
 c_5&=I(\mathsf E(c,a),I(\mathsf M(c,b),\mathsf M(a,b))),\\
 c_6&=I(\mathsf X(a,\mathsf A(a,b)),b),\\
 c_7&=I(\mathsf E(c,a),I(\mathsf E(c,b),\mathsf A(\mathsf M(a,b),\mathsf E(a,b)))),\\
 c_8&=I(I(I(c,d),a),I(I(a,c),I(b,c))).
\end{align*}
The guarded bundle label is specified by the following fixed polynomial
compositions:
\begin{align*}
 \mathsf R(b,k,h)&=\mathsf A(b,I(h,\mathsf N(\mathsf A(b,I(h,\mathsf N(\mathsf M(b,k))))))),\\
 \mathsf S(a,b,k)&=\mathsf A(k,I(\mathsf M(k,a),\mathsf M(k,b))),\\
 C_0(a,b,k,h)&=I(\mathsf S(a,b,k),\mathsf R(b,k,h)),\\
 C_1(a,b,k,h)&=\mathsf N(\mathsf A(a,C_0(a,b,k,h))),\\
 C_2(a,b,k,h)&=\mathsf N(\mathsf A(a,I(h,C_1(a,b,k,h)))),\\
 \mathsf C(a,b,k,h)&=\mathsf A(b,I(\mathsf M(b,k),C_2(a,b,k,h))),\\
 \mathsf V(a,b,c)&=I(\mathsf M(a,b),\mathsf A(c,\mathsf M(a,b))),\\
 \mathsf K(a,b,c)&=\mathsf A(c,I(\mathsf C(a,b,c,c),\mathsf N(\mathsf R(b,c,c)))),\\
 c_*&=\mathsf N(I(\mathsf V(a,b,c),\mathsf K(a,b,c))).
\end{align*}
Its guard is the polynomial equation
\begin{equation}\label{eq:guard}
 g=(c-a)^2(c-b)^2-(d+1)^2=0.
\end{equation}
For natural parameters this permits some $d$ exactly when $c\ne a,b$:
the positive integer $|(c-a)(c-b)|$ must equal $d+1$.

All rules have two padded premises and one conclusion. Their canonical
parameterization is the following.
\begin{center}\small
\begin{tabular}{llll}
\toprule
Rule & First premise & Second premise & Conclusion / condition\\
\midrule
Truth padding & 0 & 0 & 0\\
Eight unconditional axioms & 0 & 0 & $c_j(a,b,c,d)$\\
Bundle & 0 & 0 & $c_*(a,b,c)$, $g=0$\\
Modus ponens & $a$ & $I(a,b)$ & $b$\\
Generalization & $b$ & 0 & $\mathsf A(a,b)$\\
\bottomrule
\end{tabular}
\end{center}
Projecting the premise and conclusion labels in this way preserves both
directions of rule-instance existence. The four schema parameters remain
free, so no rule instance is lost.

\section{From finite refutations to bounded arithmetic paths}
\label{sec:flow}

For a two-premise rule $(p,q)\to v$, its flow vector is
$e_v-e_p-e_q$ in the free abelian group on positive labels, with $e_0=0$.
A finite proof tree of target 1 supplies a multiset of rule instances
whose total flow is $e_1$. Conversely, let $D$ be the least rule-closed set
of derivable labels. If 1 were outside $D$, summing coordinates outside
$D$ would give a contradiction: every valid rule has nonpositive total
contribution there, whereas $e_1$ has contribution one. The argument needs
only the finite support of the proposed certificate. Cycles and repeated
premises cannot create a false positive.

\subsection{Padding and the carrier}

Expand a finite derivation into a tree. Choose an integer $t\ge1$ at least
its number of rule occurrences and every value of the four parameters in
the canonical rule table. Pad to exactly $t$ occurrences using truth rows.
Thus the horizon itself bounds all four local schema coordinates.

Every label is a polynomial with nonnegative coefficients in these
parameters. A literal, 266-gate exponent certificate in
\file{equations/source690/label_majorant.json} proves
\begin{equation}\label{eq:labelbound}
 \text{every label}+1\le (t+1)^{C_0},\qquad C_0=2^{22}=4,194,304.
\end{equation}
Its computed exponent is 3,445,908; choosing the next power of two
permits 22 repeated squarings. The certificate is obtained by assigning
exponent one to an input, $\lceil\log_2c\rceil$ to a constant $c>1$,
$\max(e_1,e_2)+1$ to addition, and $e_1+e_2$ to multiplication. Since
the base $t+1$ is at least two, induction proves every recorded bound.
The bounded expression is the sum of all possible conclusion and premise
labels, plus one, so it bounds each required label plus one.

Choose
\begin{equation}\label{eq:carrier}
 M=((t+1)^t)^{C_0},\qquad M=M_0+1.
\end{equation}
Only $(t+1)^t$ is a variable-exponent relation; it is subsequently
represented by the shared Pell graph. The exponent $C_0$ is fixed and
uses only quadratic squaring gates.

\begin{lemma}[One-sided degree bootstrap]\label{lem:carrier}
Suppose $t\ge1$, all labels are at most $h\ge1$, and $M>1$ satisfies
$M\ge(h+1)^t$. Put
\[
 A(X)=\prod_{i<t}(1+v_iX),\qquad
 D(X)=(1+X)\prod_{i<t}(1+p_iX)(1+q_iX).
\]
If $A(M)=D(M)$, then $A=D$ as integer polynomials, and the rule instances
have total flow $e_1$.
\end{lemma}
\begin{proof}
We have
$A(M)\le(1+hM)^t<((h+1)M)^t\le M^{t+1}$.
If $\deg D>t$, its nonnegative integer leading coefficient would imply
$D(M)\ge M^{t+1}$, a contradiction. Thus $D$ has at most $t$ nonconstant
linear factors, including its initial factor $1+X$. Consequently
$A(1),D(1)\le(h+1)^t\le M$. Their constant coefficients are one, so
every coefficient is strictly smaller than $M$. Uniqueness of base-$M$
expansion gives $A=D$. Unique factorization in $\Q[X]$, or comparison
of the roots $-1/j$ with multiplicity for each positive label $j$,
gives the asserted flow identity.
\end{proof}

The preliminary degree argument is essential. Applying a naive bound to
all $2t$ premise factors would demand an unnecessarily larger carrier.
Equation~\eqref{eq:labelbound} makes \eqref{eq:carrier} sufficient for
Lemma~\ref{lem:carrier}, even though the axiom labels can be much larger
than $t$.

\subsection{The accumulator path and eight local columns}

Let $w(j)=1+Mj$. A finite path consists of
\begin{equation}\label{eq:path}
 S_0=M+1,\qquad S_t=1,\qquad
 S_{i+1}w(v_i)=S_iw(p_i)w(q_i)\quad(i<t).
\end{equation}
Every state is positive, by induction from the first endpoint.
Multiplication and cancellation give $A(M)=D(M)$, hence a refutation by
Lemma~\ref{lem:carrier}. Conversely, expand a proof tree from its root
to its leaves. At each step, let $S_i$ be the product of the weights of
the outstanding nonzero obligations. Replacing one conclusion by its
premises gives exactly \eqref{eq:path}; the final empty product is one.
Truth padding leaves the state unchanged.

Encode the finite state sequence by beta digits
\[
 \beta(x,y,i)=x\bmod (1+(i+1)y).
\]
For each row use the eight coordinates
\[
 (a,b,c,d,S,S',u,u'),\qquad
 (a,b,c,d)\le t,\quad S,S'\le y,\quad u,u'\le x.
\]
The two beta equations are
\begin{equation}\label{eq:beta}
 x=S+(1+(i+1)y)u,\qquad
 x=S'+(1+(i+2)y)u'.
\end{equation}
Because $S,S'\le y$ and each modulus exceeds $y$, these are exact
remainder equations. No separate remainder-bound columns are needed.

The outer endpoint rows are
\begin{align}
 x&=M_0+2+(y+1)u_0,& y&=M_0+2+r_0,\label{eq:initial}\\
 x&=1+(1+(t+1)y)u_t.\label{eq:final}
\end{align}
All variables here are natural. These rows enforce the desired endpoints.
For completeness choose $y$ above every state and divisible by $t!$.
The moduli $1+(i+1)y$, $0\le i\le t$, are pairwise coprime: a common
divisor is coprime to $y$ and divides the nonzero index difference, which
divides $t!$ and hence $y$. CRT supplies $x$ with the prescribed digits.
The resulting natural quotients satisfy $u,u'\le x$ automatically.
The order of choices is proof, then $t$, then $M$, then the finite states,
then $y$ and $x$. No trace code is required to bound itself.

\section{A nonnegative local rule polynomial}
\label{sec:local}

Write $J_h=S'(1+Mh)-S$ and use the axiom labels of
Section~\ref{sec:calculus}. The projected local Horn polynomial is
\begin{align}
 \mathcal H={}&
 \left(J_0^2\prod_{j=1}^8J_{c_j}^2\right)
 (J_{c_*}^2+g^2)\,U_{\rm mp}^2\,U_{\rm gen}^2,\label{eq:horn}\\
 U_{\rm mp}={}&S'(1+Mb)-S(1+Ma)(1+MI(a,b)),\nonumber\\
 U_{\rm gen}={}&S'(1+M\mathsf A(a,b))-S(1+Mb).\nonumber
\end{align}
Let $\mathcal L(i,z,w)$ be $\mathcal H$ plus the squares of the two
residuals in \eqref{eq:beta}; here $z=(t,x,y,M_0)$ and $w$ is the
eight-coordinate local tuple. Every summand and factor is nonnegative,
so $\mathcal L=0$ means precisely that one valid rule and its accumulator
update hold, together with both beta equations.

The degree of this intermediate polynomial is large. It is never expanded
to form the final quartic. Its finite arithmetic graph is represented by
quadratic defining equations.

\subsection{Difference squares as primitive gates}

For natural stored values $a,b$, create a new natural value $z$ using the
single quadratic residual
\begin{equation}\label{eq:squaregate}
 z-a^2+2ab-b^2=0.
\end{equation}
This graph is total and functional: its unique output is $(a-b)^2\ge0$.
It remains nonnegative even when the inputs are unrestricted integers.
Thus a circuit assembled from addition, multiplication, and difference
squares uses one gate value per operation, without maintaining separate
positive and negative evaluations of signed expressions.

For example, the bundle guard uses
\[
 (c-a)^2(c-b)^2
   =\bigl(c^2+ab-c(a+b)\bigr)^2.
\]
Two difference-square gates and one multiplication give its inner value.
Another difference square against $(d+1)^2$ gives $g^2$.
For the unprojected precursor, the independent source identity audit
checks the literal constructor DAGs and the ring identity relating its
nonnegative circuit to the inherited factored source. It does not infer
equality from numerical samples. The eight-column projection is a separate
equivalence of rule-instance existence, justified by the canonical
substitutions in Section~\ref{sec:calculus}; it is not asserted to be an
identity of the two differently parameterized source polynomials.

\subsection{A bound for a nonmonotone circuit}

Difference squares are not monotone in their inputs, so evaluating them
at coordinate bounds is not a valid majorant. Instead, at each such gate
replace $(a-b)^2$ by $(A+B)^2$ when $0\le a\le A$ and $0\le b\le B$.
Addition and multiplication preserve the resulting bounds. This gives a
monotone nonnegative circuit $\mathcal U$ that dominates $\mathcal L$
on the whole bounded box.

Put
\begin{equation}\label{eq:H}
 R=2+t+x+y+M_0,\qquad H=R^{2^{23}}.
\end{equation}
A literal 338-gate exponent certificate for
$\mathcal U+t+x+y+1$ has exponent 6,893,253, below $2^{23}=8,388,608$.
The same elementary gate induction used for the labels proves
\begin{equation}\label{eq:strictH}
 0\le\mathcal L(i,z,w)<H,
 \quad t,x,y<H,
 \quad\text{all local bounds}+1\le H.
\end{equation}
The input index is bounded by $t$ in this induction. The extra one
ensures the strict bound. The endpoint rows give $R\ge4$, more than
enough for all estimates below.

\section{Eliminating the bounded family of local rows}
\label{sec:CRT}

Set
\[
 T=H!,\qquad d_i=1+(i+1)T,\qquad B=\prod_{i<t}d_i.
\]
The $d_i$ are pairwise coprime. Every prime divisor of any $d_i$ exceeds
$H$, since every prime at most $H$ divides $H!$.
Let $(b_1,\ldots,b_8)=(t,t,t,t,y,y,x,x)$ be the local bounds.
The entire bounded family is encoded by positive integers
$A_1,\ldots,A_8$ and a natural index $C$ satisfying
\begin{align}
 B&\mid\binom{A_j}{b_j+1}\quad(1\le j\le8),\label{eq:ranges}\\
 1+(C+1)T&=B,\label{eq:canonical}\\
 \mathcal L(C,z,A)&=Bq_0\quad\text{for some }q_0\in\N.\label{eq:localmod}
\end{align}
All $A_j$ are represented as a fresh natural variable plus one.

\begin{lemma}[Exact bounded-row decoding]\label{lem:CRT}
Equations \eqref{eq:ranges}--\eqref{eq:localmod} are equivalent to
$\forall i<t\,\exists w\le b:\mathcal L(i,z,w)=0$.
Moreover their witnesses can be chosen to satisfy
\begin{equation}\label{eq:aggregate}
 A_1+\cdots+A_8+g_0=8B,\qquad g_0\in\N.
\end{equation}
\end{lemma}
\begin{proof}
Fix $i<t$ and a prime $\ell\mid d_i$. The falling-factorial identity
\[
 (b_j+1)!\binom{A_j}{b_j+1}=\prod_{r=0}^{b_j}(A_j-r)
\]
gives a residue $w_j\in[0,b_j]$ with $A_j\equiv w_j\pmod\ell$.
Equation~\eqref{eq:canonical} gives $C\equiv i\pmod\ell$, since $T$
is invertible modulo $\ell$. Polynomial congruence then yields
$\ell\mid\mathcal L(i,z,w)$. By \eqref{eq:strictH}, this nonnegative
number is smaller than $\ell$, so it is zero.

Conversely select a valid row for each $i<t$. CRT chooses $A_j$
congruent to its $j$th coordinate modulo every $d_i$. Because
$(b_j+1)!\mid T$ and $\gcd(B,T)=1$, cancellation in the falling-factorial
identity gives \eqref{eq:ranges}. The local congruence follows from
polynomial evaluation, and $\mathcal L\ge0$ gives a natural $q_0$.
Since $B\equiv1\pmod T$ and $B\ge1+T$, the canonical value
$C=(B-1)/T-1$ is natural and has every required index residue.
Choose each positive representative $A_j$ in $[1,B]$; then
$g_0=8B-\sum A_j\ge0$ gives \eqref{eq:aggregate}.
\end{proof}

The canonical index removes an unnecessary divisibility quotient. The
nonnegative local circuit similarly permits \eqref{eq:localmod} with
one natural quotient. The aggregate representative bound will supply a
radix estimate without representing eight separate upper-bound slacks.

\section{Represent only the binomial residue that is used}
\label{sec:radix}

The following elementary identity is the main arithmetic compression.

\begin{lemma}[Lower digits suffice]\label{lem:lowdigits}
Let $n,k\in\N$ and $q\ge2$. Suppose
$\binom nj<q$ for every $0\le j\le k$, with $\binom nj=0$ for $j>n$.
Then
\[
 \left\lfloor\frac{(q+1)^n}{q^k}\right\rfloor
       =\binom nk+q h\quad\text{for some }h\in\N.
\]
If $B\ge1$ and $B\mid q$, the three natural witnesses $u,r,s$ in
\begin{equation}\label{eq:fused}
 (q+1)^n=q^k B u+r,\qquad q^k=r+s+1
\end{equation}
exist exactly when $B\mid\binom nk$.
\end{lemma}
\begin{proof}
Expand $(q+1)^n=\sum_j\binom njq^j$. The terms below $k$ have total
less than $q^k$; the coefficient at $k$ is smaller than $q$; and the
higher part is divisible by $q^{k+1}$. This proves the quotient identity.
Reduction modulo $B$ erases its higher part. Equations~\eqref{eq:fused}
are precisely Euclidean division by $q^k$ with a quotient divisible by
$B$. Both directions follow. No bound on coefficients above $k$ is used.
\end{proof}

Apply \eqref{eq:fused} to the eight range requests in
\eqref{eq:ranges}. There is no variable for the exact binomial value,
no separate high-digit variable, and no exact digit-bound slack for
these requests. The three other binomial calls still require their
complete values and retain ordinary digit-extraction graphs.

\subsection{An existing factorial numerator supplies the radix}

The factorial graph for a positive $n$ uses
\begin{equation}\label{eq:factorialsemantic}
 \rho_n=(n+2)^{n+2},\qquad V_n=(\rho_n+1)^n,
 \qquad c_n=\binom{\rho_n+n}{n}.
\end{equation}
It encodes $n!=\lceil V_n/c_n\rceil$. For $n=1$ this is immediate.
For $n\ge2$,
\[
 \frac{V_n}{c_n}=n!\prod_{j=1}^n\frac{\rho_n+1}{\rho_n+j}.
\]
The product is at most one and its loss from one is less than
$n(n-1)/(2\rho_n)$. Since $n!n(n-1)/2<\rho_n$, the ratio lies strictly
above $n!-1$ and at most $n!$, proving the ceiling formula.
Two natural remainder rows represent this ceiling exactly.

Use the already requested value $V=V_H$ to define
\begin{equation}\label{eq:recycled}
 \boxed{q=B(V_H+1).}
\end{equation}
The only exact binomial values needed are
\[
 \binom{\rho_H+H}{H},\qquad
 \binom{\rho_t+t}{t},\qquad
 \binom{Q+t}{t},\qquad
 Q=(1+(t+1)T)^t+2.
\]
The first two occur in $H!$ and $t!$. The third occurs in the CRT product
formula. The remaining eight calls are only divisibility requests.

\begin{proposition}[The recycled radix is sufficient]\label{prop:radix}
For the released construction, the radix \eqref{eq:recycled}, the
aggregate bound \eqref{eq:aggregate}, and the existing arithmetic rows
make every exact and residue-only binomial graph correct. This does not
presuppose correctness of the CRT product formula.
\end{proposition}
\begin{proof}
First recover all requested powers by the simultaneous Pell theorem of
Section~\ref{sec:pell}. Its modulus row excludes a zero base. In
particular $q>0$, so \eqref{eq:recycled} gives $B\ge1$ and $q\ge V_H+1$.
This step does not require a binomial or factorial to be correct.

For any positive $n$ and $0\le j\le n$, binomial unimodality and
the product formula give
\begin{equation}\label{eq:factlow}
 \binom{\rho_n+n}{j}\le\binom{\rho_n+n}{n}
       =\prod_{i=1}^n\left(1+\frac{\rho_n}{i}\right)
       \le(\rho_n+1)^n=V_n.
\end{equation}
Thus the first factorial's lower digits are bounded by $V_H<q$.
Its exact graph is correct, and the ceiling identity now proves $T=H!$.
Since $t\le H$ and $V_n$ increases with $n$, the same argument validates
$t!$.

The raw remainder rows for the proposed CRT product already imply
$0\le B<TQ$, whether or not its numerator has yet been validated.
Let $R$ be as in \eqref{eq:H}. We have $R\ge4$, $H\ge R^4$,
$1\le t\le R$, and every range lower index is at most $R$.
The number $r=8$ of columns is at most $H$. From $T=H!\le H^H$
and $t+2\le H$,
\[
 1+(t+1)T\le H^{H+1},\qquad
 Q+t\le H^{t(H+1)+1}.
\]
The aggregate inequality gives, for every encoded coordinate,
\[
 A_j\le rB<rTQ\le H^{(t+2)(H+1)}.
\]
The same final bound holds for $Q+t$. For any of these upper arguments
$n$ and its lower argument $1\le k\le R$,
\begin{align*}
 \binom nj&\le n^j\le n^k
       \le H^{k(t+2)(H+1)}\le H^{H(H+1)}\quad(0\le j\le k),\\
 V_H&>(H+2)^{H(H+2)}>H^{H(H+1)}.
\end{align*}
Here $k(t+2)\le R(R+2)\le H$. All required lower coefficients are
therefore smaller than $V_H<q$. Lemma~\ref{lem:lowdigits} now validates
the third exact binomial and every fused range request.

Only now identify $B$ with its intended product. Indeed
\[
 T^t t!\binom{Q+t}{t}=\prod_{j=1}^t(TQ+jT)
       \equiv\prod_{j=1}^t(1+jT)\pmod{TQ-1}.
\]
The product on the right is less than $Q-2<TQ-1$, so its ordinary
remainder is exactly the intended $B$.
\end{proof}

The order in this proof resolves an otherwise serious apparent cycle:
$B$ helps define the radix, while a radix-dependent binomial helps
compute $B$. We first recover $H!$, then use the \emph{raw} bound $B<TQ$,
then validate the remaining binomials, and only finally use the product
identity. Completeness chooses the true finite arithmetic values, the
positive canonical CRT representatives, and their ordinary quotients
before selecting the auxiliary Pell witnesses.

\section{Shared Pell certificates and derived indices}
\label{sec:pell}

For $a\ge2$ let $(X_j(a),Y_j(a))$ be defined by
\[
 (X_0,Y_0)=(1,0),\qquad
 X_{j+1}=aX_j+(a^2-1)Y_j,\quad Y_{j+1}=X_j+aY_j.
\]
Then $X_j^2-(a^2-1)Y_j^2=1$. We use the standard classification
of natural solutions and the finite Matiyasevich indexed-Pell criterion.
The precise inherited interface is
\file{PellRepresentation.criterion_iff_pell}; the corresponding Mathlib
results include \file{Pell.eq_pell}, \file{Pell.matiyasevic}, and the
addition identities \cite{mathlib,pellinterface}.

\subsection{One frame recovers all positive powers}

There are 23 distinct requested power relations $m=n^k$. Every exponent
is positive by the literal source construction. Set
\[
 W=2+\sum\{\text{distinct stored atoms used as }n,k,m\},\qquad Z=z+1,
\]
and impose
\begin{equation}\label{eq:frame}
 A^2-W^3(W+2)Z^2=1.
\end{equation}
For every distinct base $n$, use a natural $T_n$ with
\begin{equation}\label{eq:basemod}
 2An=T_n+n^2+1.
\end{equation}
The latter excludes $n=0$. Classification in \eqref{eq:frame} gives
$A=X_j(W+1)$ and $WZ=Y_j(W+1)$. Modulo $W$, the recurrence yields
$Y_j(W+1)\equiv j$. Thus $W\mid j$ and $j\ge W$, because $WZ>0$.
Consequently
\[
 A\ge(W+1)^W>W^W,
 \qquad T_n\ge2A-2>W^W\ge n^k,
 \qquad m<T_n.
\]
The inequality for $T_n=2An-n^2-1$ follows by monotonicity on
$1\le n\le W$, where its minimum is $2A-2$.
For existence, take $j=W$; the same congruence makes
$Z=Y_W(W+1)/W$ a positive integer.

\subsection{Fifteen independent index tests}

For each of 15 seed exponents $k_i$, use $X_i,Y_i,S_i,R_i$ and natural
slacks so that
\begin{align*}
 Y_i&=k_i+e_i,&X_i^2-(A^2-1)Y_i^2&=1,\\
 S_i^2-(b^2-1)R_i^2&=1,&S_i&=X_i+Us_i,&R_i&=k_i+4Y_ir_i.
\end{align*}
The auxiliary values $U,V,b$ are shared, with
\begin{align*}
 U^2-(A^2-1)V^2&=1,\\
 V&=\left(\prod_{i=1}^{15}Y_i^2\right)(v_0+1),\\
 b&=1+4V(b_0+1),&b&=A+Ub_1.
\end{align*}
The finite indexed-Pell criterion proves
$(X_i,Y_i)=(X_{k_i}(A),Y_{k_i}(A))$.
The fixed 15-factor product becomes ordinary multiplication gates.

For simultaneous completeness, first choose the true seed pairs. In the
finite ring $(\Z/L\Z)[s]/(s^2-(A^2-1))$, where
$L=4\prod_iY_i^2$, the unit $A+s$ has a positive power equal to one.
Its Pell coordinates give $U\equiv1\pmod L$, $V\equiv0\pmod L$,
and $V>0$. The Pell norm gives $\gcd(U,V)=1$; $U$ is odd, so
$\gcd(U,4V)=1$. CRT chooses an arbitrarily large $b$ satisfying
$b\equiv A\pmod U$ and $b\equiv1\pmod{4V}$.
Taking the true $k_i$th Pell pair at base $b$ gives every remaining
congruence, and monotonicity in the base gives natural quotients.
This is a simultaneous witness construction, rather than an unjustified
identification of separately existential witnesses.

\subsection{Three index tests replaced by exact addition}

Introduce $D=A^2-1$ using $A^2=D+1$. The identities
\begin{align}
 X_{u+v}&=X_uX_v+DY_uY_v,\label{eq:addition}\\
 Y_{u+v}&=X_uY_v+Y_uX_v\nonumber
\end{align}
follow from multiplication in $\Z[s]/(s^2-D)$. The fixed pairs for
indices one and two are $(A,1)$ and $(2D+1,2A)$.
The actual source requests $H+2,t+2,t+1$ as well as $H,t$, so it
derives those three pairs by \eqref{eq:addition}. Every dependency points
to an earlier arithmetic node. Induction therefore fixes the derived
pairs uniquely, and they need not appear in the shared auxiliary product.

Finally a power request with index $k$ imposes
\begin{equation}\label{eq:powlast}
 X_k+nY_k=AY_k+m+T_nq_k,\qquad q_k\in\N.
\end{equation}
Set $H_j=X_j-(A-n)Y_j$. The recurrence gives
$H_{j+1}=nH_j+T_nY_j$ and $H_0=1$, hence
\[
 H_k-n^k=T_n\sum_{j=0}^{k-1}n^{k-1-j}Y_j.
\]
The quotient is nonnegative. Both $m$ and $n^k$ are below $T_n$, so
\eqref{eq:powlast} is equivalent to $m=n^k$.
This proof holds simultaneously for arbitrary proposed inputs and outputs
of all power requests. A dependency between their semantic values does
not create a circular proof of the power relation itself.

\section{From the arithmetic graph to literal quartics}
\label{sec:compile}

The source generator introduces a variable for each required arithmetic
value and a quadratic row for each gate. Addition, multiplication, and
difference-square rows are respectively
\[
 z-a-b=0,\qquad z-ab=0,\qquad z-a^2+2ab-b^2=0.
\]
Every higher product and each fixed power is decomposed into these rows.
After initial affine inlining, the complete source has 690 natural
variables and 606 quadratic residuals. The uncontracted sum of squares
has 3,762 monomials and coefficient height 63,905,979.
The generator emits the full residual and coefficient lists, not just
their counts.

\subsection{Certified natural-coordinate elimination}

\begin{lemma}[Nonnegative polynomial restoration]\label{lem:restore}
Suppose a residual is $\varepsilon(z-P(x))$, $\varepsilon\in\{-1,1\}$,
where the coordinates $x$ exclude $z$, and $P\in\Z[x]$ has an exact certificate
\[
 P=\sum_j q_jL_j^2+R,
\]
where $q_j>0$ are rational, $L_j$ are rational affine forms, and all
coefficients of $R$ are nonnegative. Deleting this residual and
substituting $z=P(x)$ in all others preserves natural solutions, with
inverse restoration $z=P(x)$. If $P$ is quadratic and every occurrence
of $z$ in a remaining row is a constant multiple of $z$ (so neither
$z^2$ nor $zx_i$ occurs), all residuals remain quadratic.
\end{lemma}
\begin{proof}
The certificate makes $P(x)\ge0$ for natural $x$; its integer
coefficients make its value integral. The monic defining row forces
exactly that value. Substitution and restoration are mutual inverses on
the zero sets. The stated occurrence condition prevents quadratic
substitution from increasing any residual degree beyond two.
\end{proof}

The producer uses exact rational algebra to find such certificates.
The independent checker expands each supplied affine square and verifies
the identity; it does not trust the certificate-search algorithm.
The final 407-variable equation removes 283 coordinates from the source,
including two restorations whose expanded polynomials have negative
coefficients. It has 323 rows. A second deterministic support limit
gives the 410-variable sparse companion.

For $\Pn$, each residual is quadratic or linear. Its sum of squares has
degree at most four. Literal expansion finds nonzero degree-four terms,
establishing exact degree four and all reported term counts and heights.
The appendix is generated directly from the checked final residual data.

\section{Integer solutions without guarding every coordinate}
\label{sec:integer}

The integer construction starts from the same 690-coordinate natural
source, before natural-coordinate contraction. Guarding every coordinate
would waste many witnesses: most signs already follow from the equations.

\begin{lemma}[Finite sign inference]\label{lem:sign}
Suppose some coordinates of an integer solution are already known
nonnegative. A residual of the form $c z=F$, $c>0$, proves $z\ge0$
in either of the following cases:
\begin{enumerate}
\item $F$ has nonnegative coefficients, and every variable occurring to
      an odd power in one of its monomials is already known nonnegative;
\item $F=L^2$ for an explicitly given integer affine form $L$.
\end{enumerate}
An ordered finite sequence of these inferences is sound.
\end{lemma}
\begin{proof}
Even powers are nonnegative for arbitrary integers, and all remaining
factors have known nonnegative sign. In the second case the right side
is a square. Division of the inequality $cz\ge0$ by $c>0$ gives the
claim. Induction over the supplied order prevents circular sign reasoning.
\end{proof}

The actual certificate requires only 167 seed coordinates. It derives the
remaining 523 signs: 506 by the first rule and 17 by exact affine-square
identities. Every inference is checked against a literal original residual.

For each seed coordinate append
\begin{equation}\label{eq:integerguard}
 z=a^2+b^2+c(c+1),\qquad a,b,c\in\Z.
\end{equation}
This forces $z\ge0$. Conversely every natural $z$ has such a representation.
Indeed $4z+1$ is not of the excluded form $4^r(8s+7)$ in Legendre's
three-square theorem. In a representation of $4z+1$ by three squares,
exactly one summand is odd. Write it as $(2c+1)^2$ and the other two
as $(2a)^2,(2b)^2$; dividing by four gives \eqref{eq:integerguard}.

Thus the guarded system has $690+3\cdot167=1,191$ integer variables and
$606+167=773$ rows. Every integer zero recovers the entire natural
source by Lemma~\ref{lem:sign}, and every natural source zero extends
using \eqref{eq:integerguard}.

Only after this equivalence is established do we perform signed monic
coordinate eliminations. Over integers, a restoration polynomial need not
be nonnegative. For a quadratic restoration, the eliminated coordinate
again occurs only as a constant scalar multiple of itself in every
remaining row, never multiplied by a variable.
The primary projection removes 303 coordinates, leaving $\Pz$ with
888 variables and 470 residuals. A smaller support threshold removes
274 coordinates and leaves $\Ps$ with 917 variables and 499 residuals.
The latter retains substantially smaller coefficients and fewer expanded
monomials. Their exact polynomial inverse maps are supplied.

\section{The ZFC theorem, proof boundaries, and reproducibility}
\label{sec:validation}

\begin{proof}[Proof of Theorem~\ref{thm:main}]
The natural arithmetic graph is equivalent to the finite bounded-row
certificate by the shared-power theorem, Proposition~\ref{prop:radix},
the factorial and product identities, and exact gate functionality.
Lemma~\ref{lem:CRT} identifies its bounded rows with a valid accumulator
path. The source polynomial implements exactly the projected $C_{222}$
rules. The path/flow equivalence of Section~\ref{sec:flow}, followed by
the inherited calibration \eqref{eq:calibration}, identifies existence
with a ZFC refutation. Lemma~\ref{lem:restore} transfers this equivalence
to $\Pn$. The sign inference, Legendre guards, and signed monic
restorations transfer it to $\Pz$ and $\Ps$.

Every new step is finite integer arithmetic, polynomial identity,
induction, CRT, or the stated classical Pell/three-square theorem.
The finite-proof and model-theoretic arguments in the inherited
calibration have their ordinary set-theoretic interpretations in ZFC.
The entire argument can therefore be carried out in ZFC over its internal
natural numbers. This proves the biconditionals, without assuming
consistency or existence of a refutation. Taking negations proves their
no-solution forms.
\end{proof}

The exact boundary is important. The mathematical theorem uses the
identified existing calculus calibration; this package is not an
independent reconstruction of all its imported logical files. The immediate
calibration files and source pins are retained, but their entire Lean
dependency closure was not rebuilt here. No new full Lean/Rocq theorem
or object-language ZFC proof term is emitted. The finite algebra audits
certify the supplied construction and its transformations; they do not
replace the unbounded mathematical arguments above.

\subsection{What is checked exactly}

The release verifier checks the following logically distinct layers:
\begin{itemize}
\item Regeneration of the complete 690-variable source from its fixed
      constructor code, and agreement with the supplied source bytes.
\item The code-constructor identities, the unprojected precursor's
      square-gate rewrite,
      the two majorant exponent ledgers, and the three actual Pell
      addition dependencies.
\item The actual recycled-radix and aggregate-bound rows, the factorial
      numerator binding, positivity of all 23 exponent atoms, and the
      15-seed auxiliary product.
\item All 283 natural inverse substitutions, including rational
      nonnegativity certificates; every sign inference and every integer
      guard; all 303 or 274 integer inverse substitutions.
\item Independent full expansion and collection of each delivered quartic,
      with actual variable occurrence, degree, support, and coefficient
      height checks.
\end{itemize}
Deliberately corrupted radix, nonnegativity, and sign-inference records
are rejected. Small Pell and lower-digit fixtures also check transcription;
their finite ranges are not used as proofs of the universal lemmas.
The complete verifier is run normally and with Python optimization enabled,
so correctness checks do not disappear when Python assertions are disabled.

\subsection{Archive contents and commands}

The archive contains this article and its editable LaTeX sources; the
complete residual and expanded coefficient lists; the uncontracted source;
the inverse and sign certificates; independent checking code; pinned
dependencies and attribution; and verification receipts.
Python 3.10 or later and its standard library suffice for the algebra.
Run from the package root:
\begin{verbatim}
python3 verify.py
python3 -O verify.py
python3 reproduce.py --output /path/to/new-empty-directory
pdflatex -interaction=nonstopmode -halt-on-error report.tex
pdflatex -interaction=nonstopmode -halt-on-error report.tex
\end{verbatim}
The reproduction command constructs coefficient data. It never searches
for a refutation or a polynomial zero. The verifier authenticates each
certificate before using it, and does not overwrite the delivered
equations. File digests identify exact bytes, while the polynomial lists
and proved transformations give the mathematical meaning.

\subsection{What remains open}

This construction improves a specific direct quartic frontier. It neither
settles the globally smallest variable count nor supplies a lower bound
against other architectures. A useful next mathematical target is to
reduce the independent Pell-index tests for the encoded coordinates, or
to replace their lower-digit range certificates by a smaller exact
relation. A useful formal target is the end-to-end internalization of this
already fixed coefficient list. Neither open task is an assumption in
the ordinary proof supplied here.

\begin{thebibliography}{20}
\bibitem{repo} Vladimir Reshetnikov and ProveIt contributors,
\emph{DiophantineConsistency research}, snapshot
\file{85ae4070019efd66ea3e58e9a67d287528094e8b}.
\url{https://github.com/VladimirReshetnikov/ProveIt/tree/85ae4070019efd66ea3e58e9a67d287528094e8b/Computability/DiophantineConsistency}.
\bibitem{synchronized} ProveIt incoming contribution,
\emph{Synchronized arithmetic certificates for ZFC refutability},
\file{ConZFC_Synchronized_Certificates_2026-10-24.zip}, at the comparison
snapshot. Its immediate source code, proofs, calibration files, notices,
and hashes are included under \file{upstream/}.
\url{https://github.com/VladimirReshetnikov/ProveIt/blob/85ae4070019efd66ea3e58e9a67d287528094e8b/docs/incoming/ConZFC_Synchronized_Certificates_2026-10-24.zip}.
\bibitem{batched} ProveIt incoming contribution,
\emph{Batched certificates},
\file{con_zfc_batched_certificates_2026-10-23.zip}, same snapshot.
The signed elimination method is inherited with attribution.
\url{https://github.com/VladimirReshetnikov/ProveIt/blob/85ae4070019efd66ea3e58e9a67d287528094e8b/docs/incoming/con_zfc_batched_certificates_2026-10-23.zip}.
\bibitem{projection} ProveIt incoming contribution,
\emph{Project first, pad second, and flatten last}, in
\file{ConZFC_Direct_Projection_7617.zip}, same snapshot. The prior proof
is copied to \file{provenance/prior_projection_proof.md}.
\url{https://github.com/VladimirReshetnikov/ProveIt/blob/85ae4070019efd66ea3e58e9a67d287528094e8b/docs/incoming/ConZFC_Direct_Projection_7617.zip}.
\bibitem{interpolation} ProveIt incoming contribution,
\file{ConZFC_Twelve_Integer_Interpolation_2026-10-24.zip}, same snapshot.
\url{https://github.com/VladimirReshetnikov/ProveIt/blob/85ae4070019efd66ea3e58e9a67d287528094e8b/docs/incoming/ConZFC_Twelve_Integer_Interpolation_2026-10-24.zip}.
\bibitem{calibration} ProveIt contributors,
\file{Code254.lean}, \file{KCode222.lean}, \file{KEndpoint.lean}, and
\file{ConstructibleChoice.lean}, calibration snapshot
\file{08200eb10fde2db408c69dffd840d6920a569b1f}; immediate files and source
pins included in the archive.
\bibitem{pellinterface} ProveIt contributors,
\file{PellRepresentation.lean}, especially
\file{criterion_iff_pell}; immediate source included under
\file{upstream/provenance/upstream/calibration/}.
\bibitem{mathlib} Mathlib contributors, \emph{Pell's equation and
Matiyasevic's theorem}, including the Lean development by Mario Carneiro.
\url{https://leanprover-community.github.io/mathlib4_docs/Mathlib/NumberTheory/PellMatiyasevic.html}.
\end{thebibliography}

\clearpage\appendix
\section{The complete 407-variable equation}\label{app:rows}

Every $x_i$ in this appendix ranges independently over $\N$, including
zero, and the indices run from 0 to 406. The single equation is
\begin{equation}\label{eq:literal}
 \boxed{\displaystyle \Pn(x_0,\ldots,x_{406})=
                  \sum_{j=0}^{322}R_j(x_0,\ldots,x_{406})^2=0.}
\end{equation}
The following is the complete residual list, with no omitted coefficients
or auxiliary definitions. The arithmetic variable names in the construction
are mapped to these indices in \file{equation.json}. In particular the
large fixed exponents used to choose witnesses do not appear as large
degrees in these residuals. The entire expanded polynomial is also
supplied as \file{expanded_polynomial.json}.

% Generated from the checked literal residual data.
\begingroup\small
\allowdisplaybreaks[4]
\setlength{\jot}{2pt}
\begin{align*}
R_{0}={}&-x_{3}^{2}+x_{4}\\[3pt]
R_{1}={}&-x_{4}^{2}+x_{5}\\[3pt]
R_{2}={}&-x_{5}^{2}+x_{6}\\[3pt]
R_{3}={}&-x_{6}^{2}+x_{7}\\[3pt]
R_{4}={}&-x_{7}^{2}+x_{8}\\[3pt]
R_{5}={}&-x_{8}^{2}+x_{9}\\[3pt]
R_{6}={}&-x_{9}^{2}+x_{10}\\[3pt]
R_{7}={}&-x_{10}^{2}+x_{11}\\[3pt]
R_{8}={}&-x_{11}^{2}+x_{12}\\[3pt]
R_{9}={}&-x_{12}^{2}+x_{13}\\[3pt]
R_{10}={}&-x_{13}^{2}+x_{14}\\[3pt]
R_{11}={}&-x_{14}^{2}+x_{15}\\[3pt]
R_{12}={}&-x_{15}^{2}+x_{16}\\[3pt]
R_{13}={}&-x_{16}^{2}+x_{17}\\[3pt]
R_{14}={}&-x_{17}^{2}+x_{18}\\[3pt]
R_{15}={}&-x_{18}^{2}+x_{19}\\[3pt]
R_{16}={}&-x_{19}^{2}+x_{20}\\[3pt]
R_{17}={}&-x_{20}^{2}+x_{21}\\[3pt]
R_{18}={}&-x_{21}^{2}+x_{22}\\[3pt]
R_{19}={}&-x_{22}^{2}+x_{23}\\[3pt]
R_{20}={}&-x_{23}^{2}+x_{24}\\[3pt]
R_{21}={}&-1-x_{2}+x_{24}^{2}\\[3pt]
R_{22}={}&-4-2x_{0}-x_{0}x_{2}-x_{0}x_{26}-2x_{2}-2x_{26}+x_{28}\\[3pt]
R_{23}={}&-2+x_{1}-x_{2}-x_{2}x_{25}-3x_{25}-x_{25}x_{26}\\[3pt]
R_{24}={}&-1+x_{1}-x_{27}-x_{27}x_{28}\\[3pt]
R_{25}={}&-25-10x_{0}-x_{0}^{2}-2x_{0}x_{1}-4x_{0}x_{2}-2x_{0}x_{26}-10x_{1}\\*
&{}-x_{1}^{2}-4x_{1}x_{2}-2x_{1}x_{26}-20x_{2}-4x_{2}^{2}-4x_{2}x_{26}\\*
&{}-10x_{26}-x_{26}^{2}+x_{38}\\[3pt]
R_{26}={}&-x_{38}^{2}+x_{39}\\[3pt]
R_{27}={}&-x_{39}^{2}+x_{40}\\[3pt]
R_{28}={}&-x_{40}^{2}+x_{41}\\[3pt]
R_{29}={}&-x_{41}^{2}+x_{42}\\[3pt]
R_{30}={}&-x_{42}^{2}+x_{43}\\[3pt]
R_{31}={}&-x_{43}^{2}+x_{44}\\[3pt]
R_{32}={}&-x_{44}^{2}+x_{45}\\[3pt]
R_{33}={}&-x_{45}^{2}+x_{46}\\[3pt]
R_{34}={}&-x_{46}^{2}+x_{47}\\[3pt]
R_{35}={}&-x_{47}^{2}+x_{48}\\[3pt]
R_{36}={}&-x_{48}^{2}+x_{49}\\[3pt]
R_{37}={}&-x_{49}^{2}+x_{50}\\[3pt]
R_{38}={}&-x_{50}^{2}+x_{51}\\[3pt]
R_{39}={}&-x_{51}^{2}+x_{52}\\[3pt]
R_{40}={}&-x_{52}^{2}+x_{53}\\[3pt]
R_{41}={}&-x_{53}^{2}+x_{54}\\[3pt]
R_{42}={}&-x_{54}^{2}+x_{55}\\[3pt]
R_{43}={}&-x_{55}^{2}+x_{56}\\[3pt]
R_{44}={}&-x_{56}^{2}+x_{57}\\[3pt]
R_{45}={}&-x_{57}^{2}+x_{58}\\[3pt]
R_{46}={}&-x_{58}^{2}+x_{59}\\[3pt]
R_{47}={}&-x_{59}^{2}+x_{60}\\[3pt]
R_{48}={}&-x_{61}-x_{61}x_{64}-x_{61}x_{65}+x_{63}+x_{65}\\[3pt]
R_{49}={}&-x_{0}x_{61}-2x_{61}+x_{66}\\[3pt]
R_{50}={}&-x_{69}-x_{69}x_{72}-x_{69}x_{73}+x_{71}+x_{73}\\[3pt]
R_{51}={}&-x_{68}x_{69}+x_{75}\\[3pt]
R_{52}={}&2-2x_{61}-x_{61}x_{67}+x_{76}+x_{78}\\[3pt]
R_{53}={}&x_{74}x_{75}-x_{76}-x_{76}x_{77}-x_{77}-x_{77}x_{78}\\[3pt]
R_{54}={}&1+x_{37}x_{61}+x_{61}-x_{76}\\[3pt]
R_{55}={}&-6-5x_{29}-x_{29}^{2}-2x_{29}x_{30}-5x_{30}-x_{30}^{2}+x_{79}\\[3pt]
R_{56}={}&-72-68x_{30}-16x_{30}^{2}-16x_{30}x_{79}-34x_{79}-4x_{79}^{2}+x_{80}\\[3pt]
R_{57}={}&-30-11x_{29}-x_{29}^{2}-4x_{29}x_{80}-22x_{80}-4x_{80}^{2}+x_{81}\\[3pt]
R_{58}={}&-380-312x_{80}-64x_{80}^{2}-32x_{80}x_{81}-78x_{81}-4x_{81}^{2}+x_{82}\\[3pt]
R_{59}={}&-6-5x_{30}-x_{30}^{2}-2x_{30}x_{31}-5x_{31}-x_{31}^{2}+x_{83}\\[3pt]
R_{60}={}&-6-5x_{29}-x_{29}^{2}-2x_{29}x_{31}-5x_{31}-x_{31}^{2}+x_{84}\\[3pt]
R_{61}={}&-414-22x_{29}-2x_{29}^{2}-16x_{29}x_{31}-8x_{29}x_{83}-200x_{30}-32x_{30}^{2}\\*
&{}-64x_{30}x_{31}-32x_{30}x_{79}-32x_{30}x_{84}-320x_{31}-64x_{31}^{2}\\*
&{}-32x_{31}x_{79}-32x_{31}x_{83}-32x_{31}x_{84}-100x_{79}-8x_{79}^{2}\\*
&{}-16x_{79}x_{84}-52x_{83}-8x_{83}^{2}-108x_{84}-8x_{84}^{2}+x_{85}\\[3pt]
R_{62}={}&-1344-800x_{30}-128x_{30}^{2}-256x_{30}x_{31}-128x_{30}x_{79}\\*
&{}-128x_{30}x_{84}-864x_{31}-128x_{31}^{2}-128x_{31}x_{79}-128x_{31}x_{84}\\*
&{}-400x_{79}-32x_{79}^{2}-64x_{79}x_{84}-432x_{84}-32x_{84}^{2}-2x_{85}\\*
&{}-2x_{85}^{2}+x_{86}\\[3pt]
R_{63}={}&-247-168x_{29}-32x_{29}^{2}-64x_{29}x_{30}-32x_{29}x_{83}-32x_{29}x_{84}\\*
&{}-188x_{30}-32x_{30}^{2}-32x_{30}x_{83}-32x_{30}x_{84}-2x_{79}-92x_{83}\\*
&{}-8x_{83}^{2}-16x_{83}x_{84}-84x_{84}-8x_{84}^{2}+x_{87}\\[3pt]
R_{64}={}&-960-672x_{29}-128x_{29}^{2}-256x_{29}x_{30}-128x_{29}x_{83}\\*
&{}-128x_{29}x_{84}-736x_{30}-128x_{30}^{2}-128x_{30}x_{83}-128x_{30}x_{84}\\*
&{}-368x_{83}-32x_{83}^{2}-64x_{83}x_{84}-336x_{84}-32x_{84}^{2}-2x_{87}\\*
&{}-2x_{87}^{2}+x_{88}\\[3pt]
R_{65}={}&-276-30x_{29}-2x_{29}^{2}-16x_{29}x_{31}-8x_{29}x_{83}-30x_{30}-2x_{30}^{2}\\*
&{}-16x_{30}x_{31}-8x_{30}x_{84}-272x_{31}-64x_{31}^{2}-32x_{31}x_{83}\\*
&{}-32x_{31}x_{84}-68x_{83}-8x_{83}^{2}-68x_{84}-8x_{84}^{2}+x_{89}\\[3pt]
R_{66}={}&-552-120x_{30}-8x_{30}^{2}-64x_{30}x_{31}-32x_{30}x_{84}-544x_{31}\\*
&{}-128x_{31}^{2}-128x_{31}x_{84}-272x_{84}-32x_{84}^{2}-2x_{89}-2x_{89}^{2}+x_{90}\\[3pt]
R_{67}={}&-247-4x_{29}-352x_{30}-128x_{30}^{2}-64x_{30}x_{79}-64x_{30}x_{83}\\*
&{}-92x_{79}-8x_{79}^{2}-16x_{79}x_{83}-84x_{83}-8x_{83}^{2}-2x_{84}+x_{91}\\[3pt]
R_{68}={}&-960-1408x_{30}-512x_{30}^{2}-256x_{30}x_{79}-256x_{30}x_{83}-368x_{79}\\*
&{}-32x_{79}^{2}-64x_{79}x_{83}-336x_{83}-32x_{83}^{2}-2x_{91}-2x_{91}^{2}+x_{92}\\[3pt]
R_{69}={}&-56-60x_{30}-16x_{30}^{2}-16x_{30}x_{79}-30x_{79}-4x_{79}^{2}+x_{93}\\[3pt]
R_{70}={}&-30-11x_{29}-x_{29}^{2}-4x_{29}x_{93}-22x_{93}-4x_{93}^{2}+x_{94}\\[3pt]
R_{71}={}&-380-312x_{93}-64x_{93}^{2}-32x_{93}x_{94}-78x_{94}-4x_{94}^{2}+x_{95}\\[3pt]
R_{72}={}&-30-11x_{30}-x_{30}^{2}-4x_{30}x_{95}-22x_{95}-4x_{95}^{2}+x_{96}\\[3pt]
R_{73}={}&-349-420x_{30}-128x_{30}^{2}-128x_{30}x_{79}-208x_{79}-32x_{79}^{2}-2x_{83}\\*
&{}+x_{97}\\[3pt]
R_{74}={}&-1368-1664x_{30}-512x_{30}^{2}-512x_{30}x_{79}-832x_{79}-128x_{79}^{2}\\*
&{}-2x_{97}-2x_{97}^{2}+x_{98}\\[3pt]
R_{75}={}&-56-60x_{29}-16x_{29}^{2}-16x_{29}x_{84}-8x_{29}x_{98}-30x_{84}-4x_{84}^{2}\\*
&{}-4x_{84}x_{98}-15x_{98}-x_{98}^{2}+x_{99}\\[3pt]
R_{76}={}&-6-5x_{31}-x_{31}^{2}-2x_{31}x_{32}-5x_{32}-x_{32}^{2}+x_{100}\\[3pt]
R_{77}={}&-224-26x_{29}-2x_{29}^{2}-16x_{29}x_{32}-8x_{29}x_{100}-288x_{31}\\*
&{}-128x_{31}^{2}-64x_{31}x_{83}-64x_{31}x_{84}-88x_{32}-32x_{32}^{2}\\*
&{}-32x_{32}x_{100}-76x_{83}-8x_{83}^{2}-16x_{83}x_{84}-68x_{84}-8x_{84}^{2}\\*
&{}-44x_{100}-8x_{100}^{2}+x_{101}\\[3pt]
R_{78}={}&-640-1152x_{31}-512x_{31}^{2}-256x_{31}x_{83}-256x_{31}x_{84}-304x_{83}\\*
&{}-32x_{83}^{2}-64x_{83}x_{84}-272x_{84}-32x_{84}^{2}-2x_{101}-2x_{101}^{2}+x_{102}\\[3pt]
R_{79}={}&-241-168x_{29}-32x_{29}^{2}-64x_{29}x_{30}-32x_{29}x_{83}-32x_{29}x_{84}\\*
&{}-184x_{30}-32x_{30}^{2}-32x_{30}x_{83}-32x_{30}x_{84}-x_{31}-92x_{83}\\*
&{}-8x_{83}^{2}-16x_{83}x_{84}-84x_{84}-8x_{84}^{2}+x_{103}\\[3pt]
R_{80}={}&-42-52x_{31}-16x_{31}^{2}-16x_{31}x_{83}-26x_{83}-4x_{83}^{2}+x_{104}\\[3pt]
R_{81}={}&-30-11x_{31}-x_{31}^{2}-4x_{31}x_{104}-22x_{104}-4x_{104}^{2}+x_{105}\\[3pt]
R_{82}={}&-306-35x_{30}-x_{30}^{2}-16x_{30}x_{104}-4x_{30}x_{105}-280x_{104}-64x_{104}^{2}\\*
&{}-32x_{104}x_{105}-70x_{105}-4x_{105}^{2}+x_{106}\\[3pt]
R_{83}={}&-4556-4320x_{104}-1024x_{104}^{2}-512x_{104}x_{105}-128x_{104}x_{106}\\*
&{}-1080x_{105}-64x_{105}^{2}-32x_{105}x_{106}-270x_{106}-4x_{106}^{2}+x_{107}\\[3pt]
R_{84}={}&-30-11x_{31}-x_{31}^{2}-4x_{31}x_{107}-22x_{107}-4x_{107}^{2}+x_{108}\\[3pt]
R_{85}={}&-306-35x_{30}-x_{30}^{2}-16x_{30}x_{107}-4x_{30}x_{108}-280x_{107}-64x_{107}^{2}\\*
&{}-32x_{107}x_{108}-70x_{108}-4x_{108}^{2}+x_{109}\\[3pt]
R_{86}={}&-1028-672x_{29}-128x_{29}^{2}-256x_{29}x_{30}-128x_{29}x_{83}\\*
&{}-128x_{29}x_{84}-736x_{30}-128x_{30}^{2}-128x_{30}x_{83}-128x_{30}x_{84}\\*
&{}-368x_{83}-32x_{83}^{2}-64x_{83}x_{84}-336x_{84}-32x_{84}^{2}-2x_{103}\\*
&{}-2x_{103}^{2}-32x_{107}-8x_{108}-2x_{109}+x_{110}\\[3pt]
R_{87}={}&-264-128x_{107}-32x_{108}-8x_{109}-2x_{110}-2x_{110}^{2}+x_{111}\\[3pt]
R_{88}={}&-2-3x_{29}-x_{29}^{2}-2x_{29}x_{111}-3x_{111}-x_{111}^{2}+x_{112}\\[3pt]
R_{89}={}&-12-28x_{111}-16x_{111}^{2}-16x_{111}x_{112}-14x_{112}-4x_{112}^{2}+x_{113}\\[3pt]
R_{90}={}&-30-11x_{31}-x_{31}^{2}-4x_{31}x_{113}-22x_{113}-4x_{113}^{2}+x_{114}\\[3pt]
R_{91}={}&-306-35x_{29}-x_{29}^{2}-16x_{29}x_{113}-4x_{29}x_{114}-280x_{113}-64x_{113}^{2}\\*
&{}-32x_{113}x_{114}-70x_{114}-4x_{114}^{2}+x_{115}\\[3pt]
R_{92}={}&-4556-4320x_{113}-1024x_{113}^{2}-512x_{113}x_{114}-128x_{113}x_{115}\\*
&{}-1080x_{114}-64x_{114}^{2}-32x_{114}x_{115}-270x_{115}-4x_{115}^{2}+x_{116}\\[3pt]
R_{93}={}&-90-76x_{31}-16x_{31}^{2}-16x_{31}x_{83}-16x_{31}x_{116}-38x_{83}-4x_{83}^{2}\\*
&{}-8x_{83}x_{116}-38x_{116}-4x_{116}^{2}+x_{117}\\[3pt]
R_{94}={}&-4556-4320x_{107}-1024x_{107}^{2}-512x_{107}x_{108}-128x_{107}x_{109}\\*
&{}-1080x_{108}-64x_{108}^{2}-32x_{108}x_{109}-270x_{109}-4x_{109}^{2}+x_{118}\\[3pt]
R_{95}={}&-682-70x_{30}-2x_{30}^{2}-32x_{30}x_{116}-8x_{30}x_{117}-592x_{116}\\*
&{}-128x_{116}^{2}-64x_{116}x_{117}-148x_{117}-8x_{117}^{2}-2x_{118}+x_{119}\\[3pt]
R_{96}={}&-x_{119}-x_{119}^{2}+x_{120}\\[3pt]
R_{97}={}&-42-52x_{30}-16x_{30}^{2}-8x_{30}x_{31}-16x_{30}x_{79}-13x_{31}-x_{31}^{2}\\*
&{}-4x_{31}x_{79}-26x_{79}-4x_{79}^{2}+x_{121}\\[3pt]
R_{98}={}&-2278-2264x_{30}-800x_{30}^{2}-800x_{30}x_{79}-160x_{30}x_{121}-70x_{31}\\*
&{}-2x_{31}^{2}-32x_{31}x_{118}-8x_{31}x_{120}-1132x_{79}-200x_{79}^{2}\\*
&{}-80x_{79}x_{121}-592x_{118}-128x_{118}^{2}-64x_{118}x_{120}-148x_{120}-8x_{120}^{2}\\*
&{}-228x_{121}-8x_{121}^{2}+x_{122}\\[3pt]
R_{99}={}&-2712-280x_{31}-8x_{31}^{2}-128x_{31}x_{118}-32x_{31}x_{120}-2368x_{118}\\*
&{}-512x_{118}^{2}-256x_{118}x_{120}-592x_{120}-32x_{120}^{2}-2x_{122}-2x_{122}^{2}+x_{123}\\[3pt]
R_{100}={}&-2-3x_{123}-x_{123}^{2}+x_{124}\\[3pt]
R_{101}={}&-1-x_{2}-x_{2}x_{34}-x_{34}+x_{125}\\[3pt]
R_{102}={}&-x_{33}^{2}+2x_{33}x_{34}-x_{34}^{2}+x_{126}\\[3pt]
R_{103}={}&-2x_{82}x_{125}-4x_{125}+x_{127}\\[3pt]
R_{104}={}&-x_{33}^{2}+2x_{33}x_{34}+2x_{33}x_{127}-x_{34}^{2}-2x_{34}x_{127}-x_{127}^{2}+x_{128}\\[3pt]
R_{105}={}&-x_{86}x_{125}+x_{129}\\[3pt]
R_{106}={}&-x_{33}^{2}+2x_{33}x_{34}+2x_{33}x_{129}-x_{34}^{2}-2x_{34}x_{129}-x_{129}^{2}+x_{130}\\[3pt]
R_{107}={}&-x_{88}x_{125}+x_{131}\\[3pt]
R_{108}={}&-x_{33}^{2}+2x_{33}x_{34}+2x_{33}x_{131}-x_{34}^{2}-2x_{34}x_{131}-x_{131}^{2}+x_{132}\\[3pt]
R_{109}={}&-x_{90}x_{125}+x_{133}\\[3pt]
R_{110}={}&-x_{33}^{2}+2x_{33}x_{34}+2x_{33}x_{133}-x_{34}^{2}-2x_{34}x_{133}-x_{133}^{2}+x_{134}\\[3pt]
R_{111}={}&-x_{92}x_{125}+x_{135}\\[3pt]
R_{112}={}&-x_{33}^{2}+2x_{33}x_{34}+2x_{33}x_{135}-x_{34}^{2}-2x_{34}x_{135}-x_{135}^{2}+x_{136}\\[3pt]
R_{113}={}&-4x_{30}x_{125}-2x_{96}x_{125}-4x_{125}+x_{137}\\[3pt]
R_{114}={}&-x_{33}^{2}+2x_{33}x_{34}+2x_{33}x_{137}-x_{34}^{2}-2x_{34}x_{137}-x_{137}^{2}+x_{138}\\[3pt]
R_{115}={}&-4x_{98}x_{125}-2x_{99}x_{125}+x_{139}\\[3pt]
R_{116}={}&-x_{33}^{2}+2x_{33}x_{34}+2x_{33}x_{139}-x_{34}^{2}-2x_{34}x_{139}-x_{139}^{2}+x_{140}\\[3pt]
R_{117}={}&-x_{102}x_{125}+x_{141}\\[3pt]
R_{118}={}&-x_{33}^{2}+2x_{33}x_{34}+2x_{33}x_{141}-x_{34}^{2}-2x_{34}x_{141}-x_{141}^{2}+x_{142}\\[3pt]
R_{119}={}&-x_{29}^{2}+2x_{29}x_{31}-x_{31}^{2}+x_{143}\\[3pt]
R_{120}={}&-x_{30}^{2}+2x_{30}x_{31}-x_{31}^{2}+x_{144}\\[3pt]
R_{121}={}&-x_{143}x_{144}+x_{145}\\[3pt]
R_{122}={}&-4-4x_{32}-x_{32}^{2}+x_{146}\\[3pt]
R_{123}={}&-x_{145}^{2}+2x_{145}x_{146}-x_{146}^{2}+x_{147}\\[3pt]
R_{124}={}&-2x_{124}x_{125}-4x_{125}+x_{148}\\[3pt]
R_{125}={}&-x_{33}^{2}+2x_{33}x_{34}+2x_{33}x_{148}-x_{34}^{2}-2x_{34}x_{148}-x_{148}^{2}+x_{149}\\[3pt]
R_{126}={}&-1-x_{2}-x_{2}x_{30}-x_{30}+x_{150}\\[3pt]
R_{127}={}&-1-x_{34}-x_{34}x_{150}-x_{150}+x_{151}\\[3pt]
R_{128}={}&-1-x_{2}-x_{2}x_{29}-x_{29}+x_{152}\\[3pt]
R_{129}={}&-4-4x_{2}-4x_{2}x_{30}-2x_{2}x_{79}-4x_{30}-2x_{79}+x_{153}\\[3pt]
R_{130}={}&-1-x_{33}-x_{33}x_{152}-x_{152}+x_{154}\\[3pt]
R_{131}={}&-x_{153}x_{154}-x_{154}+x_{155}\\[3pt]
R_{132}={}&-x_{151}^{2}+2x_{151}x_{155}-x_{155}^{2}+x_{156}\\[3pt]
R_{133}={}&-6-6x_{2}-4x_{2}x_{30}-2x_{2}x_{79}-4x_{30}-2x_{79}+x_{157}\\[3pt]
R_{134}={}&-1-x_{34}-x_{34}x_{157}-x_{157}+x_{158}\\[3pt]
R_{135}={}&-1-x_{33}-x_{33}x_{150}-x_{150}+x_{159}\\[3pt]
R_{136}={}&-x_{158}^{2}+2x_{158}x_{159}-x_{159}^{2}+x_{160}\\[3pt]
R_{137}={}&-x_{126}x_{128}+x_{161}\\[3pt]
R_{138}={}&-x_{130}x_{161}+x_{162}\\[3pt]
R_{139}={}&-x_{132}x_{162}+x_{163}\\[3pt]
R_{140}={}&-x_{134}x_{163}+x_{164}\\[3pt]
R_{141}={}&-x_{136}x_{164}+x_{165}\\[3pt]
R_{142}={}&-x_{138}x_{165}+x_{166}\\[3pt]
R_{143}={}&-x_{140}x_{166}+x_{167}\\[3pt]
R_{144}={}&-x_{142}x_{167}+x_{168}\\[3pt]
R_{145}={}&-x_{147}x_{168}-x_{149}x_{168}+x_{169}\\[3pt]
R_{146}={}&-x_{156}x_{169}+x_{170}\\[3pt]
R_{147}={}&-2-x_{2}-x_{2}x_{37}-x_{26}-x_{26}x_{37}-2x_{37}+x_{171}\\[3pt]
R_{148}={}&-1-x_{35}-x_{35}x_{171}-x_{171}+x_{172}\\[3pt]
R_{149}={}&-3-x_{2}-x_{2}x_{36}-x_{26}-x_{26}x_{36}-3x_{36}-x_{36}x_{171}-x_{171}\\*
&{}+x_{173}\\[3pt]
R_{150}={}&2-4x_{1}+2x_{1}^{2}-2x_{1}x_{33}-2x_{1}x_{34}-2x_{1}x_{172}\\*
&{}-2x_{1}x_{173}+2x_{33}+x_{33}^{2}+2x_{33}x_{172}+2x_{34}+x_{34}^{2}+2x_{34}x_{173}\\*
&{}-x_{76}x_{174}+x_{160}x_{170}+2x_{172}+x_{172}^{2}+2x_{173}+x_{173}^{2}\\[3pt]
R_{151}={}&-8-x_{29}-x_{30}-x_{31}-x_{32}-x_{33}-x_{34}-x_{35}-x_{36}+8x_{76}\\*
&{}-x_{175}\\[3pt]
R_{152}={}&2-x_{63}x_{76}+x_{64}+x_{65}-x_{76}+x_{179}\\[3pt]
R_{153}={}&-x_{64}x_{176}-x_{65}x_{176}-2x_{176}-x_{176}x_{179}+x_{180}\\[3pt]
R_{154}={}&-x_{64}x_{181}-x_{65}x_{181}-x_{179}x_{181}-2x_{181}+x_{185}\\[3pt]
R_{155}={}&x_{64}+x_{65}-x_{72}-x_{73}+x_{179}-x_{184}\\[3pt]
R_{156}={}&-x_{64}x_{186}-x_{65}x_{186}-x_{179}x_{186}-2x_{186}+x_{190}\\[3pt]
R_{157}={}&x_{182}+x_{183}-x_{187}-x_{188}\\[3pt]
R_{158}={}&1+x_{64}+x_{65}-x_{74}+x_{179}-x_{189}\\[3pt]
R_{159}={}&-x_{76}-x_{76}x_{192}-x_{76}x_{193}+x_{194}\\[3pt]
R_{160}={}&x_{192}+x_{193}-x_{196}-x_{197}\\[3pt]
R_{161}={}&x_{192}+x_{193}-x_{199}-x_{200}\\[3pt]
R_{162}={}&x_{192}+x_{193}-x_{202}-x_{203}\\[3pt]
R_{163}={}&-x_{76}-x_{76}x_{205}-x_{76}x_{206}+x_{207}\\[3pt]
R_{164}={}&x_{205}+x_{206}-x_{209}-x_{210}\\[3pt]
R_{165}={}&-x_{76}-x_{76}x_{212}-x_{76}x_{213}+x_{214}\\[3pt]
R_{166}={}&x_{212}+x_{213}-x_{216}-x_{217}\\[3pt]
R_{167}={}&-x_{218}-x_{218}x_{220}+x_{221}\\[3pt]
R_{168}={}&-x_{221}^{2}+x_{222}\\[3pt]
R_{169}={}&1-x_{219}^{2}+x_{250}\\[3pt]
R_{170}={}&-1-2x_{218}-x_{218}^{2}+x_{223}\\[3pt]
R_{171}={}&x_{222}-x_{222}x_{223}+x_{250}\\[3pt]
R_{172}={}&-x_{225}^{2}+x_{228}\\[3pt]
R_{173}={}&-1+x_{224}^{2}-x_{228}x_{250}\\[3pt]
R_{174}={}&-4x_{225}-4x_{225}x_{226}+x_{229}\\[3pt]
R_{175}={}&1-x_{219}-x_{224}x_{227}+x_{229}\\[3pt]
R_{176}={}&-1-2x_{0}-x_{0}^{2}-2x_{0}x_{233}-2x_{233}-x_{233}^{2}+x_{234}\\[3pt]
R_{177}={}&-1+x_{230}^{2}-x_{234}x_{250}\\[3pt]
R_{178}={}&-x_{232}^{2}+x_{235}\\[3pt]
R_{179}={}&-1-2x_{229}-x_{229}^{2}+x_{236}\\[3pt]
R_{180}={}&-1+x_{231}^{2}+x_{235}-x_{235}x_{236}\\[3pt]
R_{181}={}&-x_{224}x_{237}-x_{230}+x_{231}\\[3pt]
R_{182}={}&-1-x_{0}-4x_{0}x_{238}+x_{232}-4x_{233}x_{238}-4x_{238}\\[3pt]
R_{183}={}&-5-4x_{0}-x_{0}^{2}+2x_{0}x_{219}+4x_{219}-x_{239}\\[3pt]
R_{184}={}&-x_{0}x_{219}-x_{219}-x_{219}x_{233}+x_{241}\\[3pt]
R_{185}={}&2+3x_{0}+x_{0}^{2}+x_{0}x_{233}-x_{3}+x_{230}+2x_{233}-x_{239}x_{240}\\*
&{}-x_{241}\\[3pt]
R_{186}={}&-x_{60}^{2}-2x_{60}x_{245}-x_{245}^{2}+x_{246}\\[3pt]
R_{187}={}&-1+x_{242}^{2}-x_{246}x_{250}\\[3pt]
R_{188}={}&-x_{244}^{2}+x_{247}\\[3pt]
R_{189}={}&-1-x_{236}x_{247}+x_{243}^{2}+x_{247}\\[3pt]
R_{190}={}&-x_{224}x_{248}-x_{242}+x_{243}\\[3pt]
R_{191}={}&-x_{60}-4x_{60}x_{249}+x_{244}-4x_{245}x_{249}\\[3pt]
R_{192}={}&-2x_{60}x_{219}-2x_{219}x_{245}+x_{251}\\[3pt]
R_{193}={}&-2x_{219}x_{242}+x_{252}\\[3pt]
R_{194}={}&-x_{60}-2x_{60}x_{250}-x_{245}-2x_{245}x_{250}+x_{253}\\[3pt]
R_{195}={}&-5-4x_{60}-x_{60}^{2}+2x_{60}x_{219}+4x_{219}-x_{254}\\[3pt]
R_{196}={}&x_{60}x_{252}+x_{60}x_{253}-x_{62}-x_{219}x_{252}-x_{219}x_{253}+x_{242}\\*
&{}+2x_{242}x_{250}+x_{250}x_{251}+2x_{252}+2x_{253}-x_{254}x_{255}\\[3pt]
R_{197}={}&-2-2x_{62}-x_{62}^{2}+2x_{62}x_{219}+2x_{219}-x_{256}\\[3pt]
R_{198}={}&x_{60}+x_{60}x_{62}-x_{60}x_{219}+x_{62}x_{245}-x_{63}-x_{219}x_{245}+x_{242}\\*
&{}+x_{245}-x_{256}x_{257}\\[3pt]
R_{199}={}&-2-2x_{66}-x_{66}^{2}+2x_{66}x_{219}+2x_{219}-x_{258}\\[3pt]
R_{200}={}&1+x_{0}+x_{0}x_{66}+x_{66}+x_{66}x_{233}-x_{67}+x_{230}+x_{233}-x_{241}\\*
&{}-x_{258}x_{259}\\[3pt]
R_{201}={}&-1-x_{61}^{2}+2x_{61}x_{219}-x_{260}\\[3pt]
R_{202}={}&x_{0}x_{61}+x_{61}+x_{61}x_{233}-x_{68}+x_{230}-x_{241}-x_{260}x_{261}\\[3pt]
R_{203}={}&-2x_{0}x_{219}-2x_{219}-2x_{219}x_{233}+x_{262}\\[3pt]
R_{204}={}&-2x_{219}x_{230}+x_{263}\\[3pt]
R_{205}={}&-1-x_{0}-2x_{0}x_{250}-x_{233}-2x_{233}x_{250}-2x_{250}+x_{264}\\[3pt]
R_{206}={}&-10-6x_{0}-x_{0}^{2}+2x_{0}x_{219}+6x_{219}-x_{265}\\[3pt]
R_{207}={}&x_{0}x_{263}+x_{0}x_{264}-x_{70}-x_{219}x_{263}-x_{219}x_{264}+x_{230}\\*
&{}+2x_{230}x_{250}+x_{250}x_{262}+3x_{263}+3x_{264}-x_{265}x_{266}\\[3pt]
R_{208}={}&-2-2x_{70}-x_{70}^{2}+2x_{70}x_{219}+2x_{219}-x_{267}\\[3pt]
R_{209}={}&1+x_{0}+x_{0}x_{70}+x_{70}+x_{70}x_{233}-x_{71}+x_{230}+x_{233}-x_{241}\\*
&{}-x_{267}x_{268}\\[3pt]
R_{210}={}&-x_{60}^{2}-2x_{60}x_{62}-2x_{60}x_{272}-x_{62}^{2}-2x_{62}x_{272}-x_{272}^{2}+x_{273}\\[3pt]
R_{211}={}&-1-x_{250}x_{273}+x_{269}^{2}\\[3pt]
R_{212}={}&-x_{271}^{2}+x_{274}\\[3pt]
R_{213}={}&-1-x_{236}x_{274}+x_{270}^{2}+x_{274}\\[3pt]
R_{214}={}&-x_{224}x_{275}-x_{269}+x_{270}\\[3pt]
R_{215}={}&-x_{60}-4x_{60}x_{276}-x_{62}-4x_{62}x_{276}+x_{271}-4x_{272}x_{276}\\[3pt]
R_{216}={}&-10-6x_{64}-x_{64}^{2}-2x_{64}x_{65}-2x_{64}x_{179}+2x_{64}x_{219}-6x_{65}\\*
&{}-x_{65}^{2}-2x_{65}x_{179}+2x_{65}x_{219}-6x_{179}-x_{179}^{2}+2x_{179}x_{219}+6x_{219}\\*
&{}-x_{277}\\[3pt]
R_{217}={}&-1+3x_{60}+x_{60}x_{64}+x_{60}x_{65}+x_{60}x_{179}-x_{60}x_{219}+3x_{62}\\*
&{}+x_{62}x_{64}+x_{62}x_{65}+x_{62}x_{179}-x_{62}x_{219}-x_{64}-x_{64}x_{177}\\*
&{}-x_{64}x_{178}+x_{64}x_{272}-x_{65}-x_{65}x_{177}-x_{65}x_{178}+x_{65}x_{272}\\*
&{}-2x_{177}-x_{177}x_{180}-x_{178}-x_{178}x_{180}+x_{179}x_{272}-x_{180}-x_{219}x_{272}\\*
&{}+x_{269}+3x_{272}-x_{277}x_{278}\\[3pt]
R_{218}={}&-5-4x_{64}-x_{64}^{2}-2x_{64}x_{65}-2x_{64}x_{179}+2x_{64}x_{219}-4x_{65}\\*
&{}-x_{65}^{2}-2x_{65}x_{179}+2x_{65}x_{219}-4x_{179}-x_{179}^{2}+2x_{179}x_{219}+4x_{219}\\*
&{}-x_{279}\\[3pt]
R_{219}={}&-1+2x_{60}+x_{60}x_{64}+x_{60}x_{65}+x_{60}x_{179}-x_{60}x_{219}+x_{64}x_{245}\\*
&{}+x_{65}x_{245}-x_{177}-x_{178}+x_{179}x_{245}-x_{219}x_{245}+x_{242}+2x_{245}\\*
&{}-x_{279}x_{280}\\[3pt]
R_{220}={}&-1-2x_{0}-x_{0}^{2}-2x_{0}x_{70}-2x_{0}x_{284}-2x_{70}-x_{70}^{2}\\*
&{}-2x_{70}x_{284}-2x_{284}-x_{284}^{2}+x_{285}\\[3pt]
R_{221}={}&-1-x_{250}x_{285}+x_{281}^{2}\\[3pt]
R_{222}={}&-x_{283}^{2}+x_{286}\\[3pt]
R_{223}={}&-1-x_{236}x_{286}+x_{282}^{2}+x_{286}\\[3pt]
R_{224}={}&-x_{224}x_{287}-x_{281}+x_{282}\\[3pt]
R_{225}={}&-1-x_{0}-4x_{0}x_{288}-x_{70}-4x_{70}x_{288}+x_{283}-4x_{284}x_{288}-4x_{288}\\[3pt]
R_{226}={}&2+3x_{0}+x_{0}x_{64}+x_{0}x_{65}+x_{0}x_{179}-x_{0}x_{219}+x_{64}\\*
&{}+x_{64}x_{70}+x_{64}x_{284}+x_{65}+x_{65}x_{70}+x_{65}x_{284}+3x_{70}\\*
&{}+x_{70}x_{179}-x_{70}x_{219}-x_{72}-x_{72}x_{182}-x_{72}x_{183}-x_{73}\\*
&{}-x_{73}x_{182}-x_{73}x_{183}+x_{179}+x_{179}x_{284}-2x_{182}-x_{182}x_{185}-x_{183}\\*
&{}-x_{183}x_{185}-x_{185}-x_{219}-x_{219}x_{284}-x_{277}x_{289}+x_{281}+3x_{284}\\[3pt]
R_{227}={}&1+2x_{0}+x_{0}x_{64}+x_{0}x_{65}+x_{0}x_{179}+x_{64}+x_{64}x_{233}\\*
&{}+x_{65}+x_{65}x_{233}+x_{179}+x_{179}x_{233}-x_{182}-x_{183}+x_{230}+2x_{233}\\*
&{}-x_{241}-x_{279}x_{290}\\[3pt]
R_{228}={}&-9-6x_{0}-x_{0}^{2}-2x_{0}x_{67}-2x_{0}x_{294}-6x_{67}-x_{67}^{2}\\*
&{}-2x_{67}x_{294}-6x_{294}-x_{294}^{2}+x_{295}\\[3pt]
R_{229}={}&-1-x_{250}x_{295}+x_{291}^{2}\\[3pt]
R_{230}={}&-x_{293}^{2}+x_{296}\\[3pt]
R_{231}={}&-1-x_{236}x_{296}+x_{292}^{2}+x_{296}\\[3pt]
R_{232}={}&-x_{224}x_{297}-x_{291}+x_{292}\\[3pt]
R_{233}={}&-3-x_{0}-4x_{0}x_{298}-x_{67}-4x_{67}x_{298}+x_{293}-4x_{294}x_{298}\\*
&{}-12x_{298}\\[3pt]
R_{234}={}&9+3x_{0}+x_{0}x_{64}+x_{0}x_{65}+x_{0}x_{179}-x_{0}x_{219}+3x_{64}\\*
&{}+x_{64}x_{67}+x_{64}x_{294}+3x_{65}+x_{65}x_{67}+x_{65}x_{294}+3x_{67}\\*
&{}+x_{67}x_{179}-x_{67}x_{219}-x_{74}-x_{74}x_{182}-x_{74}x_{183}+3x_{179}\\*
&{}+x_{179}x_{294}-x_{182}x_{190}-x_{183}x_{190}-x_{187}-x_{190}-3x_{219}-x_{219}x_{294}\\*
&{}-x_{277}x_{299}+x_{291}+3x_{294}\\[3pt]
R_{235}={}&-1-2x_{29}-x_{29}^{2}-2x_{29}x_{303}-2x_{303}-x_{303}^{2}+x_{304}\\[3pt]
R_{236}={}&-1-x_{250}x_{304}+x_{300}^{2}\\[3pt]
R_{237}={}&-x_{302}^{2}+x_{305}\\[3pt]
R_{238}={}&-1-x_{236}x_{305}+x_{301}^{2}+x_{305}\\[3pt]
R_{239}={}&-x_{224}x_{306}-x_{300}+x_{301}\\[3pt]
R_{240}={}&-1-x_{29}-4x_{29}x_{307}+x_{302}-4x_{303}x_{307}-4x_{307}\\[3pt]
R_{241}={}&3+3x_{29}+x_{29}x_{64}+x_{29}x_{65}+x_{29}x_{179}-x_{29}x_{219}+x_{64}\\*
&{}+x_{64}x_{303}+x_{65}+x_{65}x_{303}+x_{179}+x_{179}x_{303}-x_{191}x_{194}-x_{192}\\*
&{}-x_{219}-x_{219}x_{303}-x_{277}x_{308}+x_{300}+3x_{303}\\[3pt]
R_{242}={}&-1+x_{0}x_{250}+x_{64}x_{230}+x_{64}x_{241}+x_{65}x_{230}+x_{65}x_{241}\\*
&{}+x_{179}x_{230}+x_{179}x_{241}-x_{192}-x_{193}-x_{219}x_{241}+2x_{230}+x_{233}x_{250}\\*
&{}+2x_{241}+x_{250}-x_{279}x_{309}\\[3pt]
R_{243}={}&-1-2x_{30}-x_{30}^{2}-2x_{30}x_{313}-2x_{313}-x_{313}^{2}+x_{314}\\[3pt]
R_{244}={}&-1-x_{250}x_{314}+x_{310}^{2}\\[3pt]
R_{245}={}&-x_{312}^{2}+x_{315}\\[3pt]
R_{246}={}&-1-x_{236}x_{315}+x_{311}^{2}+x_{315}\\[3pt]
R_{247}={}&-x_{224}x_{316}-x_{310}+x_{311}\\[3pt]
R_{248}={}&-1-x_{30}-4x_{30}x_{317}+x_{312}-4x_{313}x_{317}-4x_{317}\\[3pt]
R_{249}={}&3+3x_{30}+x_{30}x_{64}+x_{30}x_{65}+x_{30}x_{179}-x_{30}x_{219}+x_{64}\\*
&{}+x_{64}x_{313}+x_{65}+x_{65}x_{313}+x_{179}+x_{179}x_{313}-x_{194}x_{195}-x_{196}\\*
&{}-x_{219}-x_{219}x_{313}-x_{277}x_{318}+x_{310}+3x_{313}\\[3pt]
R_{250}={}&-1-2x_{31}-x_{31}^{2}-2x_{31}x_{322}-2x_{322}-x_{322}^{2}+x_{323}\\[3pt]
R_{251}={}&-1-x_{250}x_{323}+x_{319}^{2}\\[3pt]
R_{252}={}&-x_{321}^{2}+x_{324}\\[3pt]
R_{253}={}&-1-x_{236}x_{324}+x_{320}^{2}+x_{324}\\[3pt]
R_{254}={}&-x_{224}x_{325}-x_{319}+x_{320}\\[3pt]
R_{255}={}&-1-x_{31}-4x_{31}x_{326}+x_{321}-4x_{322}x_{326}-4x_{326}\\[3pt]
R_{256}={}&3+3x_{31}+x_{31}x_{64}+x_{31}x_{65}+x_{31}x_{179}-x_{31}x_{219}+x_{64}\\*
&{}+x_{64}x_{322}+x_{65}+x_{65}x_{322}+x_{179}+x_{179}x_{322}-x_{194}x_{198}-x_{199}\\*
&{}-x_{219}-x_{219}x_{322}-x_{277}x_{327}+x_{319}+3x_{322}\\[3pt]
R_{257}={}&-1-2x_{32}-x_{32}^{2}-2x_{32}x_{331}-2x_{331}-x_{331}^{2}+x_{332}\\[3pt]
R_{258}={}&-1-x_{250}x_{332}+x_{328}^{2}\\[3pt]
R_{259}={}&-x_{330}^{2}+x_{333}\\[3pt]
R_{260}={}&-1-x_{236}x_{333}+x_{329}^{2}+x_{333}\\[3pt]
R_{261}={}&-x_{224}x_{334}-x_{328}+x_{329}\\[3pt]
R_{262}={}&-1-x_{32}-4x_{32}x_{335}+x_{330}-4x_{331}x_{335}-4x_{335}\\[3pt]
R_{263}={}&3+3x_{32}+x_{32}x_{64}+x_{32}x_{65}+x_{32}x_{179}-x_{32}x_{219}+x_{64}\\*
&{}+x_{64}x_{331}+x_{65}+x_{65}x_{331}+x_{179}+x_{179}x_{331}-x_{194}x_{201}-x_{202}\\*
&{}-x_{219}-x_{219}x_{331}-x_{277}x_{336}+x_{328}+3x_{331}\\[3pt]
R_{264}={}&-1-2x_{33}-x_{33}^{2}-2x_{33}x_{340}-2x_{340}-x_{340}^{2}+x_{341}\\[3pt]
R_{265}={}&-1-x_{250}x_{341}+x_{337}^{2}\\[3pt]
R_{266}={}&-x_{339}^{2}+x_{342}\\[3pt]
R_{267}={}&-1-x_{236}x_{342}+x_{338}^{2}+x_{342}\\[3pt]
R_{268}={}&-x_{224}x_{343}-x_{337}+x_{338}\\[3pt]
R_{269}={}&-1-x_{33}-4x_{33}x_{344}+x_{339}-4x_{340}x_{344}-4x_{344}\\[3pt]
R_{270}={}&3+3x_{33}+x_{33}x_{64}+x_{33}x_{65}+x_{33}x_{179}-x_{33}x_{219}+x_{64}\\*
&{}+x_{64}x_{340}+x_{65}+x_{65}x_{340}+x_{179}+x_{179}x_{340}-x_{204}x_{207}-x_{205}\\*
&{}-x_{219}-x_{219}x_{340}-x_{277}x_{345}+x_{337}+3x_{340}\\[3pt]
R_{271}={}&-9-6x_{2}-x_{2}^{2}-2x_{2}x_{26}-2x_{2}x_{349}-6x_{26}-x_{26}^{2}\\*
&{}-2x_{26}x_{349}-6x_{349}-x_{349}^{2}+x_{350}\\[3pt]
R_{272}={}&-1-x_{250}x_{350}+x_{346}^{2}\\[3pt]
R_{273}={}&-x_{348}^{2}+x_{351}\\[3pt]
R_{274}={}&-1-x_{236}x_{351}+x_{347}^{2}+x_{351}\\[3pt]
R_{275}={}&-x_{224}x_{352}-x_{346}+x_{347}\\[3pt]
R_{276}={}&-3-x_{2}-4x_{2}x_{353}-x_{26}-4x_{26}x_{353}+x_{348}-4x_{349}x_{353}\\*
&{}-12x_{353}\\[3pt]
R_{277}={}&5+2x_{2}+x_{2}x_{64}+x_{2}x_{65}+x_{2}x_{179}-x_{2}x_{219}+2x_{26}\\*
&{}+x_{26}x_{64}+x_{26}x_{65}+x_{26}x_{179}-x_{26}x_{219}+3x_{64}+x_{64}x_{349}\\*
&{}+3x_{65}+x_{65}x_{349}+3x_{179}+x_{179}x_{349}-x_{205}-x_{206}-3x_{219}\\*
&{}-x_{219}x_{349}-x_{279}x_{354}+x_{346}+2x_{349}\\[3pt]
R_{278}={}&-1-2x_{34}-x_{34}^{2}-2x_{34}x_{358}-2x_{358}-x_{358}^{2}+x_{359}\\[3pt]
R_{279}={}&-1-x_{250}x_{359}+x_{355}^{2}\\[3pt]
R_{280}={}&-x_{357}^{2}+x_{360}\\[3pt]
R_{281}={}&-1-x_{236}x_{360}+x_{356}^{2}+x_{360}\\[3pt]
R_{282}={}&-x_{224}x_{361}-x_{355}+x_{356}\\[3pt]
R_{283}={}&-1-x_{34}-4x_{34}x_{362}+x_{357}-4x_{358}x_{362}-4x_{362}\\[3pt]
R_{284}={}&3+3x_{34}+x_{34}x_{64}+x_{34}x_{65}+x_{34}x_{179}-x_{34}x_{219}+x_{64}\\*
&{}+x_{64}x_{358}+x_{65}+x_{65}x_{358}+x_{179}+x_{179}x_{358}-x_{207}x_{208}-x_{209}\\*
&{}-x_{219}-x_{219}x_{358}-x_{277}x_{363}+x_{355}+3x_{358}\\[3pt]
R_{285}={}&-1-2x_{35}-x_{35}^{2}-2x_{35}x_{367}-2x_{367}-x_{367}^{2}+x_{368}\\[3pt]
R_{286}={}&-1-x_{250}x_{368}+x_{364}^{2}\\[3pt]
R_{287}={}&-x_{366}^{2}+x_{369}\\[3pt]
R_{288}={}&-1-x_{236}x_{369}+x_{365}^{2}+x_{369}\\[3pt]
R_{289}={}&-x_{224}x_{370}-x_{364}+x_{365}\\[3pt]
R_{290}={}&-1-x_{35}-4x_{35}x_{371}+x_{366}-4x_{367}x_{371}-4x_{371}\\[3pt]
R_{291}={}&3+3x_{35}+x_{35}x_{64}+x_{35}x_{65}+x_{35}x_{179}-x_{35}x_{219}+x_{64}\\*
&{}+x_{64}x_{367}+x_{65}+x_{65}x_{367}+x_{179}+x_{179}x_{367}-x_{211}x_{214}-x_{212}\\*
&{}-x_{219}-x_{219}x_{367}-x_{277}x_{372}+x_{364}+3x_{367}\\[3pt]
R_{292}={}&-1-2x_{1}-x_{1}^{2}-2x_{1}x_{376}-2x_{376}-x_{376}^{2}+x_{377}\\[3pt]
R_{293}={}&-1-x_{250}x_{377}+x_{373}^{2}\\[3pt]
R_{294}={}&-x_{375}^{2}+x_{378}\\[3pt]
R_{295}={}&-1-x_{236}x_{378}+x_{374}^{2}+x_{378}\\[3pt]
R_{296}={}&-x_{224}x_{379}-x_{373}+x_{374}\\[3pt]
R_{297}={}&-1-x_{1}-4x_{1}x_{380}+x_{375}-4x_{376}x_{380}-4x_{380}\\[3pt]
R_{298}={}&1+2x_{1}+x_{1}x_{64}+x_{1}x_{65}+x_{1}x_{179}-x_{1}x_{219}+x_{64}\\*
&{}+x_{64}x_{376}+x_{65}+x_{65}x_{376}+x_{179}+x_{179}x_{376}-x_{212}-x_{213}-x_{219}\\*
&{}-x_{219}x_{376}-x_{279}x_{381}+x_{373}+2x_{376}\\[3pt]
R_{299}={}&-1-2x_{36}-x_{36}^{2}-2x_{36}x_{385}-2x_{385}-x_{385}^{2}+x_{386}\\[3pt]
R_{300}={}&-1-x_{250}x_{386}+x_{382}^{2}\\[3pt]
R_{301}={}&-x_{384}^{2}+x_{387}\\[3pt]
R_{302}={}&-1-x_{236}x_{387}+x_{383}^{2}+x_{387}\\[3pt]
R_{303}={}&-x_{224}x_{388}-x_{382}+x_{383}\\[3pt]
R_{304}={}&-1-x_{36}-4x_{36}x_{389}+x_{384}-4x_{385}x_{389}-4x_{389}\\[3pt]
R_{305}={}&3+3x_{36}+x_{36}x_{64}+x_{36}x_{65}+x_{36}x_{179}-x_{36}x_{219}+x_{64}\\*
&{}+x_{64}x_{385}+x_{65}+x_{65}x_{385}+x_{179}+x_{179}x_{385}-x_{214}x_{215}-x_{216}\\*
&{}-x_{219}-x_{219}x_{385}-x_{277}x_{390}+x_{382}+3x_{385}\\[3pt]
R_{306}={}&-x_{234}x_{246}+x_{391}\\[3pt]
R_{307}={}&-x_{273}x_{391}+x_{392}\\[3pt]
R_{308}={}&-x_{285}x_{392}+x_{393}\\[3pt]
R_{309}={}&-x_{295}x_{393}+x_{394}\\[3pt]
R_{310}={}&-x_{304}x_{394}+x_{395}\\[3pt]
R_{311}={}&-x_{314}x_{395}+x_{396}\\[3pt]
R_{312}={}&-x_{323}x_{396}+x_{397}\\[3pt]
R_{313}={}&-x_{332}x_{397}+x_{398}\\[3pt]
R_{314}={}&-x_{341}x_{398}+x_{399}\\[3pt]
R_{315}={}&-x_{350}x_{399}+x_{400}\\[3pt]
R_{316}={}&-x_{359}x_{400}+x_{401}\\[3pt]
R_{317}={}&-x_{368}x_{401}+x_{402}\\[3pt]
R_{318}={}&-x_{377}x_{402}+x_{403}\\[3pt]
R_{319}={}&-x_{386}x_{403}+x_{404}\\[3pt]
R_{320}={}&x_{225}-x_{404}-x_{404}x_{405}\\[3pt]
R_{321}={}&-40-5x_{0}-x_{1}-x_{2}-x_{3}-x_{26}-x_{29}-x_{30}-x_{31}-x_{32}\\*
&{}-x_{33}-x_{34}-x_{35}-x_{36}-3x_{60}-x_{61}-3x_{62}-x_{63}-3x_{64}\\*
&{}-x_{64}x_{177}-x_{64}x_{178}-3x_{65}-x_{65}x_{177}-x_{65}x_{178}-x_{66}-2x_{67}\\*
&{}-x_{68}-3x_{70}-x_{71}-x_{72}-x_{72}x_{182}-x_{72}x_{183}-x_{73}-x_{73}x_{182}\\*
&{}-x_{73}x_{183}-x_{74}-x_{74}x_{182}-x_{74}x_{183}-3x_{177}-x_{177}x_{180}-2x_{178}\\*
&{}-x_{178}x_{180}-2x_{179}-x_{180}-3x_{182}-x_{182}x_{185}-x_{182}x_{190}-2x_{183}\\*
&{}-x_{183}x_{185}-x_{183}x_{190}-x_{185}-x_{187}-x_{190}-x_{191}x_{194}-2x_{192}-x_{193}\\*
&{}-x_{194}x_{195}-x_{194}x_{198}-x_{194}x_{201}-x_{196}-x_{199}-x_{202}-x_{204}x_{207}\\*
&{}-2x_{205}-x_{206}-x_{207}x_{208}-x_{209}-x_{211}x_{214}-x_{212}+x_{406}\\[3pt]
R_{322}={}&-1-x_{212}-x_{213}-x_{214}x_{215}-x_{216}+x_{218}-x_{406}\\[3pt]
\end{align*}
\endgroup


\clearpage
\section{The complete 888-integer equation}\label{app:integer888}

Here every $x_i$ ranges independently over $\Z$, and the indices run
from 0 to 887. The following complete list defines the single equation
\[
 \boxed{\displaystyle Q_{888}(x_0,\ldots,x_{887})
                  =\sum_{j=0}^{469}S_j(x_0,\ldots,x_{887})^2=0.}
\]
Its natural-coordinate meaning and all inverse substitutions are certified
by the files in \file{equations/integer888/}. These coordinates are
independent of the coordinates in Appendix~\ref{app:rows}; the repeated
letter $x$ denotes the argument of the equation currently being defined.

% Generated from the checked literal residual data.
\begingroup\small
\allowdisplaybreaks[4]
\setlength{\jot}{2pt}
\begin{align*}
S_{0}={}&-x_{4}^{2}+x_{5}\\[3pt]
S_{1}={}&-x_{5}^{2}+x_{6}\\[3pt]
S_{2}={}&-x_{6}^{2}+x_{7}\\[3pt]
S_{3}={}&-x_{7}^{2}+x_{8}\\[3pt]
S_{4}={}&-x_{8}^{2}+x_{9}\\[3pt]
S_{5}={}&-x_{9}^{2}+x_{10}\\[3pt]
S_{6}={}&-x_{10}^{2}+x_{11}\\[3pt]
S_{7}={}&-x_{11}^{2}+x_{12}\\[3pt]
S_{8}={}&-x_{12}^{2}+x_{13}\\[3pt]
S_{9}={}&-x_{13}^{2}+x_{14}\\[3pt]
S_{10}={}&-x_{14}^{2}+x_{15}\\[3pt]
S_{11}={}&-x_{15}^{2}+x_{16}\\[3pt]
S_{12}={}&-x_{16}^{2}+x_{17}\\[3pt]
S_{13}={}&-x_{17}^{2}+x_{18}\\[3pt]
S_{14}={}&-x_{18}^{2}+x_{19}\\[3pt]
S_{15}={}&-x_{19}^{2}+x_{20}\\[3pt]
S_{16}={}&-x_{20}^{2}+x_{21}\\[3pt]
S_{17}={}&-x_{21}^{2}+x_{22}\\[3pt]
S_{18}={}&-x_{22}^{2}+x_{23}\\[3pt]
S_{19}={}&-x_{23}^{2}+x_{24}\\[3pt]
S_{20}={}&-x_{24}^{2}+x_{25}\\[3pt]
S_{21}={}&1+x_{3}-x_{25}^{2}\\[3pt]
S_{22}={}&-x_{0}x_{2}-2x_{2}+x_{28}\\[3pt]
S_{23}={}&-2+x_{1}-x_{2}x_{26}-x_{3}-x_{26}\\[3pt]
S_{24}={}&-1+x_{1}-x_{27}-x_{27}x_{28}\\[3pt]
S_{25}={}&-9-6x_{0}-x_{0}^{2}-2x_{0}x_{1}-2x_{0}x_{2}-2x_{0}x_{3}-6x_{1}\\*
&{}-x_{1}^{2}-2x_{1}x_{2}-2x_{1}x_{3}-6x_{2}-x_{2}^{2}-2x_{2}x_{3}-6x_{3}\\*
&{}-x_{3}^{2}+x_{38}\\[3pt]
S_{26}={}&-x_{38}^{2}+x_{39}\\[3pt]
S_{27}={}&-x_{39}^{2}+x_{40}\\[3pt]
S_{28}={}&-x_{40}^{2}+x_{41}\\[3pt]
S_{29}={}&-x_{41}^{2}+x_{42}\\[3pt]
S_{30}={}&-x_{42}^{2}+x_{43}\\[3pt]
S_{31}={}&-x_{43}^{2}+x_{44}\\[3pt]
S_{32}={}&-x_{44}^{2}+x_{45}\\[3pt]
S_{33}={}&-x_{45}^{2}+x_{46}\\[3pt]
S_{34}={}&-x_{46}^{2}+x_{47}\\[3pt]
S_{35}={}&-x_{47}^{2}+x_{48}\\[3pt]
S_{36}={}&-x_{48}^{2}+x_{49}\\[3pt]
S_{37}={}&-x_{49}^{2}+x_{50}\\[3pt]
S_{38}={}&-x_{50}^{2}+x_{51}\\[3pt]
S_{39}={}&-x_{51}^{2}+x_{52}\\[3pt]
S_{40}={}&-x_{52}^{2}+x_{53}\\[3pt]
S_{41}={}&-x_{53}^{2}+x_{54}\\[3pt]
S_{42}={}&-x_{54}^{2}+x_{55}\\[3pt]
S_{43}={}&-x_{55}^{2}+x_{56}\\[3pt]
S_{44}={}&-x_{56}^{2}+x_{57}\\[3pt]
S_{45}={}&-x_{57}^{2}+x_{58}\\[3pt]
S_{46}={}&-x_{58}^{2}+x_{59}\\[3pt]
S_{47}={}&-x_{59}^{2}+x_{222}-x_{225}\\[3pt]
S_{48}={}&-x_{0}x_{60}-2x_{60}+x_{64}\\[3pt]
S_{49}={}&-x_{66}x_{67}+x_{71}\\[3pt]
S_{50}={}&-x_{70}x_{71}+x_{72}x_{74}+x_{73}\\[3pt]
S_{51}={}&1-2x_{60}-x_{60}x_{65}+x_{72}\\[3pt]
S_{52}={}&-1-x_{37}x_{60}-x_{60}+x_{73}\\[3pt]
S_{53}={}&-6-5x_{29}-x_{29}^{2}-2x_{29}x_{30}-5x_{30}-x_{30}^{2}+x_{75}\\[3pt]
S_{54}={}&-72-68x_{30}-16x_{30}^{2}-16x_{30}x_{75}-34x_{75}-4x_{75}^{2}+x_{76}\\[3pt]
S_{55}={}&-30-11x_{29}-x_{29}^{2}-4x_{29}x_{76}-22x_{76}-4x_{76}^{2}+x_{77}\\[3pt]
S_{56}={}&-380-312x_{76}-64x_{76}^{2}-32x_{76}x_{77}-78x_{77}-4x_{77}^{2}+x_{78}\\[3pt]
S_{57}={}&-6-5x_{30}-x_{30}^{2}-2x_{30}x_{31}-5x_{31}-x_{31}^{2}+x_{79}\\[3pt]
S_{58}={}&-30-11x_{29}-x_{29}^{2}-8x_{29}x_{31}-4x_{29}x_{79}-44x_{31}-16x_{31}^{2}\\*
&{}-16x_{31}x_{79}-22x_{79}-4x_{79}^{2}+x_{80}\\[3pt]
S_{59}={}&-6-5x_{29}-x_{29}^{2}-2x_{29}x_{31}-5x_{31}-x_{31}^{2}+x_{81}\\[3pt]
S_{60}={}&-156-100x_{30}-16x_{30}^{2}-32x_{30}x_{31}-16x_{30}x_{75}-16x_{30}x_{81}\\*
&{}-100x_{31}-16x_{31}^{2}-16x_{31}x_{75}-16x_{31}x_{81}-50x_{75}-4x_{75}^{2}\\*
&{}-8x_{75}x_{81}-50x_{81}-4x_{81}^{2}+x_{82}\\[3pt]
S_{61}={}&-1806-2720x_{31}-1024x_{31}^{2}-512x_{31}x_{79}-128x_{31}x_{80}\\*
&{}-512x_{31}x_{81}-128x_{31}x_{82}-680x_{79}-64x_{79}^{2}-32x_{79}x_{80}\\*
&{}-128x_{79}x_{81}-32x_{79}x_{82}-170x_{80}-4x_{80}^{2}-32x_{80}x_{81}\\*
&{}-8x_{80}x_{82}-680x_{81}-64x_{81}^{2}-32x_{81}x_{82}-170x_{82}-4x_{82}^{2}\\*
&{}+x_{83}\\[3pt]
S_{62}={}&-110-84x_{29}-16x_{29}^{2}-32x_{29}x_{30}-16x_{29}x_{79}-16x_{29}x_{81}\\*
&{}-84x_{30}-16x_{30}^{2}-16x_{30}x_{79}-16x_{30}x_{81}-42x_{79}-4x_{79}^{2}\\*
&{}-8x_{79}x_{81}-42x_{81}-4x_{81}^{2}+x_{84}\\[3pt]
S_{63}={}&-756-1100x_{30}-400x_{30}^{2}-80x_{30}x_{75}-320x_{30}x_{79}-80x_{30}x_{84}\\*
&{}-110x_{75}-4x_{75}^{2}-32x_{75}x_{79}-8x_{75}x_{84}-440x_{79}-64x_{79}^{2}\\*
&{}-32x_{79}x_{84}-110x_{84}-4x_{84}^{2}+x_{85}\\[3pt]
S_{64}={}&-56-15x_{29}-x_{29}^{2}-8x_{29}x_{31}-4x_{29}x_{79}-60x_{31}-16x_{31}^{2}\\*
&{}-16x_{31}x_{79}-30x_{79}-4x_{79}^{2}+x_{86}\\[3pt]
S_{65}={}&-56-15x_{30}-x_{30}^{2}-8x_{30}x_{31}-4x_{30}x_{81}-60x_{31}-16x_{31}^{2}\\*
&{}-16x_{31}x_{81}-30x_{81}-4x_{81}^{2}+x_{87}\\[3pt]
S_{66}={}&-2756-3360x_{31}-1024x_{31}^{2}-512x_{31}x_{79}-512x_{31}x_{81}\\*
&{}-128x_{31}x_{86}-128x_{31}x_{87}-840x_{79}-64x_{79}^{2}-128x_{79}x_{81}\\*
&{}-32x_{79}x_{86}-32x_{79}x_{87}-840x_{81}-64x_{81}^{2}-32x_{81}x_{86}\\*
&{}-32x_{81}x_{87}-210x_{86}-4x_{86}^{2}-8x_{86}x_{87}-210x_{87}-4x_{87}^{2}\\*
&{}+x_{88}\\[3pt]
S_{67}={}&-110-168x_{30}-64x_{30}^{2}-32x_{30}x_{75}-32x_{30}x_{79}-42x_{75}\\*
&{}-4x_{75}^{2}-8x_{75}x_{79}-42x_{79}-4x_{79}^{2}+x_{89}\\[3pt]
S_{68}={}&-756-220x_{29}-16x_{29}^{2}-128x_{29}x_{30}-64x_{29}x_{75}-16x_{29}x_{81}\\*
&{}-16x_{29}x_{89}-880x_{30}-256x_{30}^{2}-256x_{30}x_{75}-64x_{30}x_{81}\\*
&{}-64x_{30}x_{89}-440x_{75}-64x_{75}^{2}-32x_{75}x_{81}-32x_{75}x_{89}\\*
&{}-110x_{81}-4x_{81}^{2}-8x_{81}x_{89}-110x_{89}-4x_{89}^{2}+x_{90}\\[3pt]
S_{69}={}&-56-60x_{30}-16x_{30}^{2}-16x_{30}x_{75}-30x_{75}-4x_{75}^{2}+x_{91}\\[3pt]
S_{70}={}&-30-11x_{29}-x_{29}^{2}-4x_{29}x_{91}-22x_{91}-4x_{91}^{2}+x_{92}\\[3pt]
S_{71}={}&-380-312x_{91}-64x_{91}^{2}-32x_{91}x_{92}-78x_{92}-4x_{92}^{2}+x_{93}\\[3pt]
S_{72}={}&-30-11x_{30}-x_{30}^{2}-4x_{30}x_{93}-22x_{93}-4x_{93}^{2}+x_{94}\\[3pt]
S_{73}={}&-156-200x_{30}-64x_{30}^{2}-64x_{30}x_{75}-100x_{75}-16x_{75}^{2}+x_{95}\\[3pt]
S_{74}={}&-1406-1500x_{30}-400x_{30}^{2}-320x_{30}x_{75}-80x_{30}x_{79}\\*
&{}-80x_{30}x_{95}-600x_{75}-64x_{75}^{2}-32x_{75}x_{79}-32x_{75}x_{95}\\*
&{}-150x_{79}-4x_{79}^{2}-8x_{79}x_{95}-150x_{95}-4x_{95}^{2}+x_{96}\\[3pt]
S_{75}={}&-16256-1020x_{29}-16x_{29}^{2}-512x_{29}x_{30}-256x_{29}x_{75}\\*
&{}-16x_{29}x_{81}-64x_{29}x_{95}-16x_{29}x_{96}-16320x_{30}-4096x_{30}^{2}\\*
&{}-4096x_{30}x_{75}-256x_{30}x_{81}-1024x_{30}x_{95}-256x_{30}x_{96}\\*
&{}-8160x_{75}-1024x_{75}^{2}-128x_{75}x_{81}-512x_{75}x_{95}-128x_{75}x_{96}\\*
&{}-510x_{81}-4x_{81}^{2}-32x_{81}x_{95}-8x_{81}x_{96}-2040x_{95}-64x_{95}^{2}\\*
&{}-32x_{95}x_{96}-510x_{96}-4x_{96}^{2}+x_{97}\\[3pt]
S_{76}={}&-6-5x_{31}-x_{31}^{2}-2x_{31}x_{32}-5x_{32}-x_{32}^{2}+x_{98}\\[3pt]
S_{77}={}&-30-11x_{29}-x_{29}^{2}-8x_{29}x_{32}-4x_{29}x_{98}-44x_{32}-16x_{32}^{2}\\*
&{}-16x_{32}x_{98}-22x_{98}-4x_{98}^{2}+x_{99}\\[3pt]
S_{78}={}&-72-136x_{31}-64x_{31}^{2}-32x_{31}x_{79}-32x_{31}x_{81}-34x_{79}-4x_{79}^{2}\\*
&{}-8x_{79}x_{81}-34x_{81}-4x_{81}^{2}+x_{100}\\[3pt]
S_{79}={}&-420-164x_{29}-16x_{29}^{2}-128x_{29}x_{31}-64x_{29}x_{79}-16x_{29}x_{99}\\*
&{}-16x_{29}x_{100}-656x_{31}-256x_{31}^{2}-256x_{31}x_{79}-64x_{31}x_{99}\\*
&{}-64x_{31}x_{100}-328x_{79}-64x_{79}^{2}-32x_{79}x_{99}-32x_{79}x_{100}\\*
&{}-82x_{99}-4x_{99}^{2}-8x_{99}x_{100}-82x_{100}-4x_{100}^{2}+x_{101}\\[3pt]
S_{80}={}&-462-688x_{30}-256x_{30}^{2}-32x_{30}x_{31}-256x_{30}x_{79}-64x_{30}x_{84}\\*
&{}-43x_{31}-x_{31}^{2}-16x_{31}x_{79}-4x_{31}x_{84}-344x_{79}-64x_{79}^{2}\\*
&{}-32x_{79}x_{84}-86x_{84}-4x_{84}^{2}+x_{102}\\[3pt]
S_{81}={}&-42-52x_{31}-16x_{31}^{2}-16x_{31}x_{79}-26x_{79}-4x_{79}^{2}+x_{103}\\[3pt]
S_{82}={}&-30-11x_{31}-x_{31}^{2}-4x_{31}x_{103}-22x_{103}-4x_{103}^{2}+x_{104}\\[3pt]
S_{83}={}&-306-35x_{30}-x_{30}^{2}-16x_{30}x_{103}-4x_{30}x_{104}-280x_{103}-64x_{103}^{2}\\*
&{}-32x_{103}x_{104}-70x_{104}-4x_{104}^{2}+x_{105}\\[3pt]
S_{84}={}&-4556-4320x_{103}-1024x_{103}^{2}-512x_{103}x_{104}-128x_{103}x_{105}\\*
&{}-1080x_{104}-64x_{104}^{2}-32x_{104}x_{105}-270x_{105}-4x_{105}^{2}+x_{106}\\[3pt]
S_{85}={}&-30-11x_{31}-x_{31}^{2}-4x_{31}x_{106}-22x_{106}-4x_{106}^{2}+x_{107}\\[3pt]
S_{86}={}&-306-35x_{30}-x_{30}^{2}-16x_{30}x_{106}-4x_{30}x_{107}-280x_{106}-64x_{106}^{2}\\*
&{}-32x_{106}x_{107}-70x_{107}-4x_{107}^{2}+x_{108}\\[3pt]
S_{87}={}&-22052-19008x_{30}-4096x_{30}^{2}-4096x_{30}x_{79}-1024x_{30}x_{84}\\*
&{}-256x_{30}x_{102}-4096x_{30}x_{106}-1024x_{30}x_{107}-256x_{30}x_{108}\\*
&{}-9504x_{79}-1024x_{79}^{2}-512x_{79}x_{84}-128x_{79}x_{102}-2048x_{79}x_{106}\\*
&{}-512x_{79}x_{107}-128x_{79}x_{108}-2376x_{84}-64x_{84}^{2}-32x_{84}x_{102}\\*
&{}-512x_{84}x_{106}-128x_{84}x_{107}-32x_{84}x_{108}-594x_{102}-4x_{102}^{2}\\*
&{}-128x_{102}x_{106}-32x_{102}x_{107}-8x_{102}x_{108}-9504x_{106}-1024x_{106}^{2}\\*
&{}-512x_{106}x_{107}-128x_{106}x_{108}-2376x_{107}-64x_{107}^{2}-32x_{107}x_{108}\\*
&{}-594x_{108}-4x_{108}^{2}+x_{109}\\[3pt]
S_{88}={}&-70490-531x_{29}-x_{29}^{2}-256x_{29}x_{106}-64x_{29}x_{107}-16x_{29}x_{108}\\*
&{}-4x_{29}x_{109}-67968x_{106}-16384x_{106}^{2}-8192x_{106}x_{107}-2048x_{106}x_{108}\\*
&{}-512x_{106}x_{109}-16992x_{107}-1024x_{107}^{2}-512x_{107}x_{108}-128x_{107}x_{109}\\*
&{}-4248x_{108}-64x_{108}^{2}-32x_{108}x_{109}-1062x_{109}-4x_{109}^{2}+x_{110}\\[3pt]
S_{89}={}&-1122540-1084928x_{106}-262144x_{106}^{2}-131072x_{106}x_{107}\\*
&{}-32768x_{106}x_{108}-8192x_{106}x_{109}-2048x_{106}x_{110}-271232x_{107}\\*
&{}-16384x_{107}^{2}-8192x_{107}x_{108}-2048x_{107}x_{109}-512x_{107}x_{110}\\*
&{}-67808x_{108}-1024x_{108}^{2}-512x_{108}x_{109}-128x_{108}x_{110}-16952x_{109}\\*
&{}-64x_{109}^{2}-32x_{109}x_{110}-4238x_{110}-4x_{110}^{2}+x_{111}\\[3pt]
S_{90}={}&-30-11x_{31}-x_{31}^{2}-4x_{31}x_{111}-22x_{111}-4x_{111}^{2}+x_{112}\\[3pt]
S_{91}={}&-306-35x_{29}-x_{29}^{2}-16x_{29}x_{111}-4x_{29}x_{112}-280x_{111}-64x_{111}^{2}\\*
&{}-32x_{111}x_{112}-70x_{112}-4x_{112}^{2}+x_{113}\\[3pt]
S_{92}={}&-4556-4320x_{111}-1024x_{111}^{2}-512x_{111}x_{112}-128x_{111}x_{113}\\*
&{}-1080x_{112}-64x_{112}^{2}-32x_{112}x_{113}-270x_{113}-4x_{113}^{2}+x_{114}\\[3pt]
S_{93}={}&-90-76x_{31}-16x_{31}^{2}-16x_{31}x_{79}-16x_{31}x_{114}-38x_{79}-4x_{79}^{2}\\*
&{}-8x_{79}x_{114}-38x_{114}-4x_{114}^{2}+x_{115}\\[3pt]
S_{94}={}&-306-35x_{30}-x_{30}^{2}-16x_{30}x_{114}-4x_{30}x_{115}-280x_{114}-64x_{114}^{2}\\*
&{}-32x_{114}x_{115}-70x_{115}-4x_{115}^{2}+x_{116}\\[3pt]
S_{95}={}&-4556-4320x_{106}-1024x_{106}^{2}-512x_{106}x_{107}-128x_{106}x_{108}\\*
&{}-1080x_{107}-64x_{107}^{2}-32x_{107}x_{108}-270x_{108}-4x_{108}^{2}+x_{117}\\[3pt]
S_{96}={}&-4970-4512x_{114}-1024x_{114}^{2}-512x_{114}x_{115}-128x_{114}x_{116}\\*
&{}-128x_{114}x_{117}-1128x_{115}-64x_{115}^{2}-32x_{115}x_{116}-32x_{115}x_{117}\\*
&{}-282x_{116}-4x_{116}^{2}-8x_{116}x_{117}-282x_{117}-4x_{117}^{2}+x_{118}\\[3pt]
S_{97}={}&-306-35x_{31}-x_{31}^{2}-16x_{31}x_{117}-4x_{31}x_{118}-280x_{117}-64x_{117}^{2}\\*
&{}-32x_{117}x_{118}-70x_{118}-4x_{118}^{2}+x_{119}\\[3pt]
S_{98}={}&-42-52x_{30}-16x_{30}^{2}-8x_{30}x_{31}-16x_{30}x_{75}-13x_{31}-x_{31}^{2}\\*
&{}-4x_{31}x_{75}-26x_{75}-4x_{75}^{2}+x_{120}\\[3pt]
S_{99}={}&-756-1100x_{30}-400x_{30}^{2}-400x_{30}x_{75}-80x_{30}x_{120}-550x_{75}\\*
&{}-100x_{75}^{2}-40x_{75}x_{120}-110x_{120}-4x_{120}^{2}+x_{121}\\[3pt]
S_{100}={}&-23870-19776x_{30}-4096x_{30}^{2}-4096x_{30}x_{75}-4096x_{30}x_{117}\\*
&{}-1024x_{30}x_{118}-256x_{30}x_{119}-1024x_{30}x_{120}-256x_{30}x_{121}\\*
&{}-9888x_{75}-1024x_{75}^{2}-2048x_{75}x_{117}-512x_{75}x_{118}-128x_{75}x_{119}\\*
&{}-512x_{75}x_{120}-128x_{75}x_{121}-9888x_{117}-1024x_{117}^{2}-512x_{117}x_{118}\\*
&{}-128x_{117}x_{119}-512x_{117}x_{120}-128x_{117}x_{121}-2472x_{118}-64x_{118}^{2}\\*
&{}-32x_{118}x_{119}-128x_{118}x_{120}-32x_{118}x_{121}-618x_{119}-4x_{119}^{2}\\*
&{}-32x_{119}x_{120}-8x_{119}x_{121}-2472x_{120}-64x_{120}^{2}-32x_{120}x_{121}\\*
&{}-618x_{121}-4x_{121}^{2}+x_{122}\\[3pt]
S_{101}={}&-70490-67968x_{117}-16384x_{117}^{2}-8192x_{117}x_{118}-2048x_{117}x_{119}\\*
&{}-512x_{117}x_{122}-16992x_{118}-1024x_{118}^{2}-512x_{118}x_{119}-128x_{118}x_{122}\\*
&{}-4248x_{119}-64x_{119}^{2}-32x_{119}x_{122}-1062x_{122}-4x_{122}^{2}+x_{123}\\[3pt]
S_{102}={}&-1-x_{3}-x_{3}x_{34}-x_{34}+x_{124}\\[3pt]
S_{103}={}&-x_{33}^{2}+2x_{33}x_{34}-x_{34}^{2}+x_{125}\\[3pt]
S_{104}={}&-2x_{78}x_{124}-4x_{124}+x_{126}\\[3pt]
S_{105}={}&-x_{33}^{2}+2x_{33}x_{34}+2x_{33}x_{126}-x_{34}^{2}-2x_{34}x_{126}-x_{126}^{2}+x_{127}\\[3pt]
S_{106}={}&-64x_{31}x_{124}-32x_{81}x_{124}-8x_{82}x_{124}-2x_{83}x_{124}-96x_{124}+x_{128}\\[3pt]
S_{107}={}&-x_{33}^{2}+2x_{33}x_{34}+2x_{33}x_{128}-x_{34}^{2}-2x_{34}x_{128}-x_{128}^{2}+x_{129}\\[3pt]
S_{108}={}&-64x_{30}x_{124}-32x_{79}x_{124}-8x_{84}x_{124}-2x_{85}x_{124}-80x_{124}+x_{130}\\[3pt]
S_{109}={}&-x_{33}^{2}+2x_{33}x_{34}+2x_{33}x_{130}-x_{34}^{2}-2x_{34}x_{130}-x_{130}^{2}+x_{131}\\[3pt]
S_{110}={}&-64x_{31}x_{124}-32x_{81}x_{124}-8x_{87}x_{124}-2x_{88}x_{124}-104x_{124}+x_{132}\\[3pt]
S_{111}={}&-x_{33}^{2}+2x_{33}x_{34}+2x_{33}x_{132}-x_{34}^{2}-2x_{34}x_{132}-x_{132}^{2}+x_{133}\\[3pt]
S_{112}={}&-64x_{30}x_{124}-32x_{75}x_{124}-8x_{89}x_{124}-2x_{90}x_{124}-80x_{124}+x_{134}\\[3pt]
S_{113}={}&-x_{33}^{2}+2x_{33}x_{34}+2x_{33}x_{134}-x_{34}^{2}-2x_{34}x_{134}-x_{134}^{2}+x_{135}\\[3pt]
S_{114}={}&-4x_{30}x_{124}-2x_{94}x_{124}-4x_{124}+x_{136}\\[3pt]
S_{115}={}&-x_{33}^{2}+2x_{33}x_{34}+2x_{33}x_{136}-x_{34}^{2}-2x_{34}x_{136}-x_{136}^{2}+x_{137}\\[3pt]
S_{116}={}&-256x_{30}x_{124}-128x_{75}x_{124}-32x_{95}x_{124}-8x_{96}x_{124}-2x_{97}x_{124}\\*
&{}-480x_{124}+x_{138}\\[3pt]
S_{117}={}&-x_{33}^{2}+2x_{33}x_{34}+2x_{33}x_{138}-x_{34}^{2}-2x_{34}x_{138}-x_{138}^{2}+x_{139}\\[3pt]
S_{118}={}&-64x_{31}x_{124}-32x_{79}x_{124}-8x_{100}x_{124}-2x_{101}x_{124}-64x_{124}+x_{140}\\[3pt]
S_{119}={}&-x_{33}^{2}+2x_{33}x_{34}+2x_{33}x_{140}-x_{34}^{2}-2x_{34}x_{140}-x_{140}^{2}+x_{141}\\[3pt]
S_{120}={}&-x_{29}^{2}+2x_{29}x_{31}-x_{31}^{2}+x_{142}\\[3pt]
S_{121}={}&-x_{30}^{2}+2x_{30}x_{31}-x_{31}^{2}+x_{143}\\[3pt]
S_{122}={}&-x_{142}x_{143}+x_{144}\\[3pt]
S_{123}={}&-4-4x_{32}-x_{32}^{2}+x_{145}\\[3pt]
S_{124}={}&-x_{144}^{2}+2x_{144}x_{145}-x_{145}^{2}+x_{146}\\[3pt]
S_{125}={}&-2x_{123}x_{124}-4x_{124}+x_{147}\\[3pt]
S_{126}={}&-x_{33}^{2}+2x_{33}x_{34}+2x_{33}x_{147}-x_{34}^{2}-2x_{34}x_{147}-x_{147}^{2}+x_{148}\\[3pt]
S_{127}={}&-1-x_{3}-x_{3}x_{30}-x_{30}+x_{149}\\[3pt]
S_{128}={}&-1-x_{34}-x_{34}x_{149}-x_{149}+x_{150}\\[3pt]
S_{129}={}&-1-x_{3}-x_{3}x_{29}-x_{29}+x_{151}\\[3pt]
S_{130}={}&-4-4x_{3}-4x_{3}x_{30}-2x_{3}x_{75}-4x_{30}-2x_{75}+x_{152}\\[3pt]
S_{131}={}&-1-x_{33}-x_{33}x_{151}-x_{151}+x_{153}\\[3pt]
S_{132}={}&-x_{152}x_{153}-x_{153}+x_{154}\\[3pt]
S_{133}={}&-x_{150}^{2}+2x_{150}x_{154}-x_{154}^{2}+x_{155}\\[3pt]
S_{134}={}&-6-6x_{3}-4x_{3}x_{30}-2x_{3}x_{75}-4x_{30}-2x_{75}+x_{156}\\[3pt]
S_{135}={}&-1-x_{34}-x_{34}x_{156}-x_{156}+x_{157}\\[3pt]
S_{136}={}&-1-x_{33}-x_{33}x_{149}-x_{149}+x_{158}\\[3pt]
S_{137}={}&-x_{157}^{2}+2x_{157}x_{158}-x_{158}^{2}+x_{159}\\[3pt]
S_{138}={}&-x_{125}x_{127}+x_{160}\\[3pt]
S_{139}={}&-x_{129}x_{160}+x_{161}\\[3pt]
S_{140}={}&-x_{131}x_{161}+x_{162}\\[3pt]
S_{141}={}&-x_{133}x_{162}+x_{163}\\[3pt]
S_{142}={}&-x_{135}x_{163}+x_{164}\\[3pt]
S_{143}={}&-x_{137}x_{164}+x_{165}\\[3pt]
S_{144}={}&-x_{139}x_{165}+x_{166}\\[3pt]
S_{145}={}&-x_{141}x_{166}+x_{167}\\[3pt]
S_{146}={}&-x_{146}x_{167}-x_{148}x_{167}+x_{168}\\[3pt]
S_{147}={}&-x_{155}x_{168}+x_{169}\\[3pt]
S_{148}={}&-x_{2}-x_{2}x_{37}+x_{170}\\[3pt]
S_{149}={}&-1-x_{35}-x_{35}x_{170}-x_{170}+x_{171}\\[3pt]
S_{150}={}&-2+4x_{1}-2x_{1}^{2}+2x_{1}x_{33}+2x_{1}x_{34}+2x_{1}x_{171}\\*
&{}+2x_{1}x_{172}-2x_{33}-x_{33}^{2}-2x_{33}x_{171}-2x_{34}-x_{34}^{2}-2x_{34}x_{172}\\*
&{}+x_{73}x_{173}-x_{159}x_{169}-2x_{171}-x_{171}^{2}-2x_{172}-x_{172}^{2}\\[3pt]
S_{151}={}&-1-x_{2}-x_{2}x_{36}-x_{36}-x_{36}x_{170}-x_{170}+x_{172}\\[3pt]
S_{152}={}&1+x_{62}-x_{63}x_{73}-x_{73}+x_{176}\\[3pt]
S_{153}={}&-x_{62}x_{175}-x_{175}-x_{175}x_{176}+x_{177}\\[3pt]
S_{154}={}&-x_{62}x_{179}-x_{176}x_{179}-x_{179}+x_{180}\\[3pt]
S_{155}={}&-x_{62}x_{181}-x_{176}x_{181}-x_{181}+x_{182}\\[3pt]
S_{156}={}&-x_{73}x_{183}+x_{185}\\[3pt]
S_{157}={}&-x_{73}x_{189}+x_{191}\\[3pt]
S_{158}={}&-x_{73}x_{193}+x_{195}\\[3pt]
S_{159}={}&-x_{197}-x_{197}x_{199}+x_{200}\\[3pt]
S_{160}={}&-x_{200}^{2}+x_{201}\\[3pt]
S_{161}={}&1-x_{198}^{2}+x_{230}\\[3pt]
S_{162}={}&-1-2x_{197}-x_{197}^{2}+x_{202}\\[3pt]
S_{163}={}&x_{201}-x_{201}x_{202}+x_{230}\\[3pt]
S_{164}={}&-x_{204}^{2}+x_{207}\\[3pt]
S_{165}={}&1-x_{203}^{2}+x_{207}x_{230}\\[3pt]
S_{166}={}&-4x_{204}-4x_{204}x_{205}+x_{208}\\[3pt]
S_{167}={}&1-x_{198}-x_{203}x_{206}+x_{208}\\[3pt]
S_{168}={}&-x_{210}^{2}+x_{213}\\[3pt]
S_{169}={}&1-x_{209}^{2}+x_{213}x_{230}\\[3pt]
S_{170}={}&-x_{212}^{2}+x_{214}\\[3pt]
S_{171}={}&1-x_{211}^{2}-x_{214}+x_{214}x_{215}\\[3pt]
S_{172}={}&-1-2x_{208}-x_{208}^{2}+x_{215}\\[3pt]
S_{173}={}&-x_{203}x_{216}-x_{209}+x_{211}\\[3pt]
S_{174}={}&-1-x_{0}-4x_{210}x_{217}+x_{212}\\[3pt]
S_{175}={}&5+4x_{0}+x_{0}^{2}-2x_{0}x_{198}-4x_{198}+x_{218}\\[3pt]
S_{176}={}&-x_{0}x_{210}+x_{4}-x_{209}-2x_{210}+x_{218}x_{219}+x_{220}\\[3pt]
S_{177}={}&-x_{198}x_{210}+x_{220}\\[3pt]
S_{178}={}&-x_{222}^{2}+x_{226}\\[3pt]
S_{179}={}&1-x_{221}^{2}+x_{226}x_{230}\\[3pt]
S_{180}={}&-x_{224}^{2}+x_{227}\\[3pt]
S_{181}={}&1+x_{215}x_{227}-x_{223}^{2}-x_{227}\\[3pt]
S_{182}={}&-x_{203}x_{228}-x_{221}+x_{223}\\[3pt]
S_{183}={}&-x_{222}-4x_{222}x_{229}+x_{224}+x_{225}\\[3pt]
S_{184}={}&x_{61}+x_{198}x_{232}+x_{198}x_{233}-x_{221}-2x_{221}x_{230}-x_{222}x_{232}\\*
&{}-x_{222}x_{233}+x_{225}x_{232}+x_{225}x_{233}-x_{230}x_{231}-2x_{232}-2x_{233}\\*
&{}+x_{234}x_{235}\\[3pt]
S_{185}={}&-2x_{198}x_{222}+x_{231}\\[3pt]
S_{186}={}&-2x_{198}x_{221}+x_{232}\\[3pt]
S_{187}={}&-x_{222}-2x_{222}x_{230}+x_{233}\\[3pt]
S_{188}={}&-5+4x_{198}+2x_{198}x_{222}-2x_{198}x_{225}-4x_{222}-x_{222}^{2}+2x_{222}x_{225}\\*
&{}+4x_{225}-x_{225}^{2}-x_{234}\\[3pt]
S_{189}={}&2+2x_{61}+x_{61}^{2}-2x_{61}x_{198}-2x_{198}+x_{236}\\[3pt]
S_{190}={}&-x_{61}x_{222}+x_{63}+x_{198}x_{222}-x_{221}-x_{222}+x_{236}x_{237}\\[3pt]
S_{191}={}&2+2x_{64}+x_{64}^{2}-2x_{64}x_{198}-2x_{198}+x_{238}\\[3pt]
S_{192}={}&-x_{64}x_{210}+x_{65}-x_{209}-x_{210}+x_{220}+x_{238}x_{239}\\[3pt]
S_{193}={}&1+x_{60}^{2}-2x_{60}x_{198}+x_{240}\\[3pt]
S_{194}={}&-x_{60}x_{210}+x_{66}-x_{209}+x_{220}+x_{240}x_{241}\\[3pt]
S_{195}={}&-x_{0}x_{243}-x_{0}x_{244}+x_{68}+x_{198}x_{243}+x_{198}x_{244}-x_{209}\\*
&{}-2x_{209}x_{230}-x_{230}x_{242}-3x_{243}-3x_{244}+x_{245}x_{246}\\[3pt]
S_{196}={}&-2x_{198}x_{210}+x_{242}\\[3pt]
S_{197}={}&-2x_{198}x_{209}+x_{243}\\[3pt]
S_{198}={}&-x_{210}-2x_{210}x_{230}+x_{244}\\[3pt]
S_{199}={}&10+6x_{0}+x_{0}^{2}-2x_{0}x_{198}-6x_{198}+x_{245}\\[3pt]
S_{200}={}&2+2x_{68}+x_{68}^{2}-2x_{68}x_{198}-2x_{198}+x_{247}\\[3pt]
S_{201}={}&-x_{68}x_{210}-x_{209}-x_{210}+x_{220}+x_{247}x_{248}+x_{459}^{2}+x_{460}^{2}+x_{461}\\*
&{}+x_{461}^{2}\\[3pt]
S_{202}={}&-x_{250}^{2}+x_{253}\\[3pt]
S_{203}={}&1+x_{230}x_{253}-x_{249}^{2}\\[3pt]
S_{204}={}&-x_{252}^{2}+x_{254}\\[3pt]
S_{205}={}&1+x_{215}x_{254}-x_{251}^{2}-x_{254}\\[3pt]
S_{206}={}&-x_{203}x_{255}-x_{249}+x_{251}\\[3pt]
S_{207}={}&-x_{61}-x_{222}+x_{225}-4x_{250}x_{256}+x_{252}\\[3pt]
S_{208}={}&-5-4x_{62}-x_{62}^{2}-2x_{62}x_{176}+2x_{62}x_{198}-4x_{176}-x_{176}^{2}\\*
&{}+2x_{176}x_{198}+4x_{198}-x_{257}\\[3pt]
S_{209}={}&x_{62}x_{174}-x_{62}x_{250}+x_{174}x_{177}-x_{176}x_{250}+x_{198}x_{250}-x_{249}\\*
&{}-2x_{250}+x_{257}x_{258}+x_{486}^{2}+x_{487}^{2}+x_{488}+x_{488}^{2}\\[3pt]
S_{210}={}&-2-2x_{62}-x_{62}^{2}-2x_{62}x_{176}+2x_{62}x_{198}-2x_{176}-x_{176}^{2}\\*
&{}+2x_{176}x_{198}+2x_{198}-x_{259}\\[3pt]
S_{211}={}&-x_{62}x_{222}+x_{174}-x_{176}x_{222}+x_{198}x_{222}-x_{221}-x_{222}+x_{259}x_{260}\\[3pt]
S_{212}={}&-x_{262}^{2}+x_{265}\\[3pt]
S_{213}={}&1+x_{230}x_{265}-x_{261}^{2}\\[3pt]
S_{214}={}&-x_{264}^{2}+x_{266}\\[3pt]
S_{215}={}&1+x_{215}x_{266}-x_{263}^{2}-x_{266}\\[3pt]
S_{216}={}&-x_{203}x_{267}-x_{261}+x_{263}\\[3pt]
S_{217}={}&-1-x_{0}-x_{68}-4x_{262}x_{268}+x_{264}\\[3pt]
S_{218}={}&-x_{62}x_{262}+x_{69}x_{178}-x_{176}x_{262}+x_{178}x_{180}+x_{198}x_{262}+x_{257}x_{269}\\*
&{}-x_{261}-2x_{262}+x_{498}^{2}+x_{499}^{2}+x_{500}+x_{500}^{2}\\[3pt]
S_{219}={}&-x_{62}x_{210}-x_{176}x_{210}+x_{178}-x_{209}-x_{210}+x_{220}+x_{259}x_{270}\\[3pt]
S_{220}={}&-x_{272}^{2}+x_{275}\\[3pt]
S_{221}={}&1+x_{230}x_{275}-x_{271}^{2}\\[3pt]
S_{222}={}&-x_{274}^{2}+x_{276}\\[3pt]
S_{223}={}&1+x_{215}x_{276}-x_{273}^{2}-x_{276}\\[3pt]
S_{224}={}&-x_{203}x_{277}-x_{271}+x_{273}\\[3pt]
S_{225}={}&-3-x_{0}-x_{65}-4x_{272}x_{278}+x_{274}\\[3pt]
S_{226}={}&-x_{62}x_{272}+x_{70}x_{178}-x_{176}x_{272}+x_{178}x_{182}+x_{198}x_{272}+x_{257}x_{279}\\*
&{}-x_{271}-2x_{272}+x_{510}^{2}+x_{511}^{2}+x_{512}+x_{512}^{2}\\[3pt]
S_{227}={}&-x_{281}^{2}+x_{284}\\[3pt]
S_{228}={}&1+x_{230}x_{284}-x_{280}^{2}\\[3pt]
S_{229}={}&-x_{283}^{2}+x_{285}\\[3pt]
S_{230}={}&1+x_{215}x_{285}-x_{282}^{2}-x_{285}\\[3pt]
S_{231}={}&-x_{203}x_{286}-x_{280}+x_{282}\\[3pt]
S_{232}={}&-1-x_{29}-4x_{281}x_{287}+x_{283}\\[3pt]
S_{233}={}&-x_{62}x_{281}-x_{176}x_{281}+x_{184}x_{185}+x_{198}x_{281}+x_{257}x_{288}-x_{280}\\*
&{}-2x_{281}+x_{522}^{2}+x_{523}^{2}+x_{524}+x_{524}^{2}\\[3pt]
S_{234}={}&-x_{62}x_{209}-x_{62}x_{220}-x_{176}x_{209}-x_{176}x_{220}+x_{183}+x_{198}x_{220}\\*
&{}-x_{209}-x_{210}x_{230}-x_{220}+x_{259}x_{289}\\[3pt]
S_{235}={}&-x_{291}^{2}+x_{294}\\[3pt]
S_{236}={}&1+x_{230}x_{294}-x_{290}^{2}\\[3pt]
S_{237}={}&-x_{293}^{2}+x_{295}\\[3pt]
S_{238}={}&1+x_{215}x_{295}-x_{292}^{2}-x_{295}\\[3pt]
S_{239}={}&-x_{203}x_{296}-x_{290}+x_{292}\\[3pt]
S_{240}={}&-1-x_{30}-4x_{291}x_{297}+x_{293}\\[3pt]
S_{241}={}&-x_{62}x_{291}-x_{176}x_{291}+x_{185}x_{186}+x_{198}x_{291}+x_{257}x_{298}-x_{290}\\*
&{}-2x_{291}+x_{531}^{2}+x_{532}^{2}+x_{533}+x_{533}^{2}\\[3pt]
S_{242}={}&-x_{300}^{2}+x_{303}\\[3pt]
S_{243}={}&1+x_{230}x_{303}-x_{299}^{2}\\[3pt]
S_{244}={}&-x_{302}^{2}+x_{304}\\[3pt]
S_{245}={}&1+x_{215}x_{304}-x_{301}^{2}-x_{304}\\[3pt]
S_{246}={}&-x_{203}x_{305}-x_{299}+x_{301}\\[3pt]
S_{247}={}&-1-x_{31}-4x_{300}x_{306}+x_{302}\\[3pt]
S_{248}={}&-x_{62}x_{300}-x_{176}x_{300}+x_{185}x_{187}+x_{198}x_{300}+x_{257}x_{307}-x_{299}\\*
&{}-2x_{300}+x_{540}^{2}+x_{541}^{2}+x_{542}+x_{542}^{2}\\[3pt]
S_{249}={}&-x_{309}^{2}+x_{312}\\[3pt]
S_{250}={}&1+x_{230}x_{312}-x_{308}^{2}\\[3pt]
S_{251}={}&-x_{311}^{2}+x_{313}\\[3pt]
S_{252}={}&1+x_{215}x_{313}-x_{310}^{2}-x_{313}\\[3pt]
S_{253}={}&-x_{203}x_{314}-x_{308}+x_{310}\\[3pt]
S_{254}={}&-1-x_{32}-4x_{309}x_{315}+x_{311}\\[3pt]
S_{255}={}&-x_{62}x_{309}-x_{176}x_{309}+x_{185}x_{188}+x_{198}x_{309}+x_{257}x_{316}-x_{308}\\*
&{}-2x_{309}+x_{549}^{2}+x_{550}^{2}+x_{551}+x_{551}^{2}\\[3pt]
S_{256}={}&-x_{318}^{2}+x_{321}\\[3pt]
S_{257}={}&1+x_{230}x_{321}-x_{317}^{2}\\[3pt]
S_{258}={}&-x_{320}^{2}+x_{322}\\[3pt]
S_{259}={}&1+x_{215}x_{322}-x_{319}^{2}-x_{322}\\[3pt]
S_{260}={}&-x_{203}x_{323}-x_{317}+x_{319}\\[3pt]
S_{261}={}&-1-x_{33}-4x_{318}x_{324}+x_{320}\\[3pt]
S_{262}={}&-x_{62}x_{318}-x_{176}x_{318}+x_{190}x_{191}+x_{198}x_{318}+x_{257}x_{325}-x_{317}\\*
&{}-2x_{318}+x_{558}^{2}+x_{559}^{2}+x_{560}+x_{560}^{2}\\[3pt]
S_{263}={}&-x_{327}^{2}+x_{330}\\[3pt]
S_{264}={}&1+x_{230}x_{330}-x_{326}^{2}\\[3pt]
S_{265}={}&-x_{329}^{2}+x_{331}\\[3pt]
S_{266}={}&1+x_{215}x_{331}-x_{328}^{2}-x_{331}\\[3pt]
S_{267}={}&-x_{203}x_{332}-x_{326}+x_{328}\\[3pt]
S_{268}={}&-1-x_{2}-4x_{327}x_{333}+x_{329}\\[3pt]
S_{269}={}&-x_{62}x_{327}-x_{176}x_{327}+x_{189}+x_{198}x_{327}+x_{259}x_{334}-x_{326}-x_{327}\\[3pt]
S_{270}={}&-x_{336}^{2}+x_{339}\\[3pt]
S_{271}={}&1+x_{230}x_{339}-x_{335}^{2}\\[3pt]
S_{272}={}&-x_{338}^{2}+x_{340}\\[3pt]
S_{273}={}&1+x_{215}x_{340}-x_{337}^{2}-x_{340}\\[3pt]
S_{274}={}&-x_{203}x_{341}-x_{335}+x_{337}\\[3pt]
S_{275}={}&-1-x_{34}-4x_{336}x_{342}+x_{338}\\[3pt]
S_{276}={}&-x_{62}x_{336}-x_{176}x_{336}+x_{191}x_{192}+x_{198}x_{336}+x_{257}x_{343}-x_{335}\\*
&{}-2x_{336}+x_{567}^{2}+x_{568}^{2}+x_{569}+x_{569}^{2}\\[3pt]
S_{277}={}&-x_{345}^{2}+x_{348}\\[3pt]
S_{278}={}&1+x_{230}x_{348}-x_{344}^{2}\\[3pt]
S_{279}={}&-x_{347}^{2}+x_{349}\\[3pt]
S_{280}={}&1+x_{215}x_{349}-x_{346}^{2}-x_{349}\\[3pt]
S_{281}={}&-x_{203}x_{350}-x_{344}+x_{346}\\[3pt]
S_{282}={}&-1-x_{35}-4x_{345}x_{351}+x_{347}\\[3pt]
S_{283}={}&-x_{62}x_{345}-x_{176}x_{345}+x_{194}x_{195}+x_{198}x_{345}+x_{257}x_{352}-x_{344}\\*
&{}-2x_{345}+x_{576}^{2}+x_{577}^{2}+x_{578}+x_{578}^{2}\\[3pt]
S_{284}={}&-x_{354}^{2}+x_{357}\\[3pt]
S_{285}={}&1+x_{230}x_{357}-x_{353}^{2}\\[3pt]
S_{286}={}&-x_{356}^{2}+x_{358}\\[3pt]
S_{287}={}&1+x_{215}x_{358}-x_{355}^{2}-x_{358}\\[3pt]
S_{288}={}&-x_{203}x_{359}-x_{353}+x_{355}\\[3pt]
S_{289}={}&-1-x_{1}-4x_{354}x_{360}+x_{356}\\[3pt]
S_{290}={}&-x_{62}x_{354}-x_{176}x_{354}+x_{193}+x_{198}x_{354}+x_{259}x_{361}-x_{353}-x_{354}\\[3pt]
S_{291}={}&-x_{363}^{2}+x_{366}\\[3pt]
S_{292}={}&1+x_{230}x_{366}-x_{362}^{2}\\[3pt]
S_{293}={}&-x_{365}^{2}+x_{367}\\[3pt]
S_{294}={}&1+x_{215}x_{367}-x_{364}^{2}-x_{367}\\[3pt]
S_{295}={}&-x_{203}x_{368}-x_{362}+x_{364}\\[3pt]
S_{296}={}&-1-x_{36}-4x_{363}x_{369}+x_{365}\\[3pt]
S_{297}={}&-x_{62}x_{363}-x_{176}x_{363}+x_{195}x_{196}+x_{198}x_{363}+x_{257}x_{370}-x_{362}\\*
&{}-2x_{363}+x_{585}^{2}+x_{586}^{2}+x_{587}+x_{587}^{2}\\[3pt]
S_{298}={}&-x_{213}x_{226}+x_{371}\\[3pt]
S_{299}={}&-x_{253}x_{371}+x_{372}\\[3pt]
S_{300}={}&-x_{265}x_{372}+x_{373}\\[3pt]
S_{301}={}&-x_{275}x_{373}+x_{374}\\[3pt]
S_{302}={}&-x_{284}x_{374}+x_{375}\\[3pt]
S_{303}={}&-x_{294}x_{375}+x_{376}\\[3pt]
S_{304}={}&-x_{303}x_{376}+x_{377}\\[3pt]
S_{305}={}&-x_{312}x_{377}+x_{378}\\[3pt]
S_{306}={}&-x_{321}x_{378}+x_{379}\\[3pt]
S_{307}={}&-x_{330}x_{379}+x_{380}\\[3pt]
S_{308}={}&-x_{339}x_{380}+x_{381}\\[3pt]
S_{309}={}&-x_{348}x_{381}+x_{382}\\[3pt]
S_{310}={}&-x_{357}x_{382}+x_{383}\\[3pt]
S_{311}={}&-x_{366}x_{383}+x_{384}\\[3pt]
S_{312}={}&x_{204}-x_{384}-x_{384}x_{385}\\[3pt]
S_{313}={}&-27-5x_{0}-x_{1}-x_{2}-x_{4}-x_{29}-x_{30}-x_{31}-x_{32}-x_{33}\\*
&{}-x_{34}-x_{35}-x_{36}-x_{60}-3x_{61}-x_{63}-x_{64}-2x_{65}-x_{66}\\*
&{}-3x_{68}-3x_{222}+3x_{225}+x_{386}-x_{459}^{2}-x_{460}^{2}-x_{461}-x_{461}^{2}\\[3pt]
S_{314}={}&-3-2x_{62}-x_{62}x_{174}-x_{69}x_{178}-x_{70}x_{178}-x_{174}-x_{174}x_{177}\\*
&{}-2x_{176}-x_{178}-x_{178}x_{180}-x_{178}x_{182}-x_{183}-x_{184}x_{185}-x_{185}x_{186}\\*
&{}-x_{185}x_{187}-x_{185}x_{188}-x_{189}-x_{190}x_{191}-x_{191}x_{192}-x_{193}\\*
&{}-x_{194}x_{195}-x_{195}x_{196}+x_{197}-x_{386}-x_{486}^{2}-x_{487}^{2}-x_{488}-x_{488}^{2}\\*
&{}-x_{498}^{2}-x_{499}^{2}-x_{500}-x_{500}^{2}-x_{510}^{2}-x_{511}^{2}-x_{512}-x_{512}^{2}-x_{522}^{2}\\*
&{}-x_{523}^{2}-x_{524}-x_{524}^{2}-x_{531}^{2}-x_{532}^{2}-x_{533}-x_{533}^{2}-x_{540}^{2}-x_{541}^{2}\\*
&{}-x_{542}-x_{542}^{2}-x_{549}^{2}-x_{550}^{2}-x_{551}-x_{551}^{2}-x_{558}^{2}-x_{559}^{2}-x_{560}\\*
&{}-x_{560}^{2}-x_{567}^{2}-x_{568}^{2}-x_{569}-x_{569}^{2}-x_{576}^{2}-x_{577}^{2}-x_{578}-x_{578}^{2}\\*
&{}-x_{585}^{2}-x_{586}^{2}-x_{587}-x_{587}^{2}\\[3pt]
S_{315}={}&x_{0}-x_{387}^{2}-x_{388}^{2}-x_{389}-x_{389}^{2}\\[3pt]
S_{316}={}&x_{3}-x_{390}^{2}-x_{391}^{2}-x_{392}-x_{392}^{2}\\[3pt]
S_{317}={}&x_{4}-x_{393}^{2}-x_{394}^{2}-x_{395}-x_{395}^{2}\\[3pt]
S_{318}={}&x_{26}-x_{396}^{2}-x_{397}^{2}-x_{398}-x_{398}^{2}\\[3pt]
S_{319}={}&-2+x_{2}-x_{3}-x_{399}^{2}-x_{400}^{2}-x_{401}-x_{401}^{2}\\[3pt]
S_{320}={}&x_{27}-x_{402}^{2}-x_{403}^{2}-x_{404}-x_{404}^{2}\\[3pt]
S_{321}={}&x_{29}-x_{405}^{2}-x_{406}^{2}-x_{407}-x_{407}^{2}\\[3pt]
S_{322}={}&x_{30}-x_{408}^{2}-x_{409}^{2}-x_{410}-x_{410}^{2}\\[3pt]
S_{323}={}&x_{31}-x_{411}^{2}-x_{412}^{2}-x_{413}-x_{413}^{2}\\[3pt]
S_{324}={}&x_{32}-x_{414}^{2}-x_{415}^{2}-x_{416}-x_{416}^{2}\\[3pt]
S_{325}={}&x_{33}-x_{417}^{2}-x_{418}^{2}-x_{419}-x_{419}^{2}\\[3pt]
S_{326}={}&x_{34}-x_{420}^{2}-x_{421}^{2}-x_{422}-x_{422}^{2}\\[3pt]
S_{327}={}&x_{35}-x_{423}^{2}-x_{424}^{2}-x_{425}-x_{425}^{2}\\[3pt]
S_{328}={}&x_{36}-x_{426}^{2}-x_{427}^{2}-x_{428}-x_{428}^{2}\\[3pt]
S_{329}={}&x_{37}-x_{429}^{2}-x_{430}^{2}-x_{431}-x_{431}^{2}\\[3pt]
S_{330}={}&x_{60}-x_{432}^{2}-x_{433}^{2}-x_{434}-x_{434}^{2}\\[3pt]
S_{331}={}&x_{61}-x_{435}^{2}-x_{436}^{2}-x_{437}-x_{437}^{2}\\[3pt]
S_{332}={}&x_{63}-x_{438}^{2}-x_{439}^{2}-x_{440}-x_{440}^{2}\\[3pt]
S_{333}={}&-1-x_{60}x_{62}+x_{62}+x_{63}-x_{441}^{2}-x_{442}^{2}-x_{443}-x_{443}^{2}\\[3pt]
S_{334}={}&x_{60}x_{62}-x_{63}-x_{444}^{2}-x_{445}^{2}-x_{446}-x_{446}^{2}\\[3pt]
S_{335}={}&x_{65}-x_{447}^{2}-x_{448}^{2}-x_{449}-x_{449}^{2}\\[3pt]
S_{336}={}&x_{66}-x_{450}^{2}-x_{451}^{2}-x_{452}-x_{452}^{2}\\[3pt]
S_{337}={}&x_{67}-x_{453}^{2}-x_{454}^{2}-x_{455}-x_{455}^{2}\\[3pt]
S_{338}={}&x_{68}-x_{456}^{2}-x_{457}^{2}-x_{458}-x_{458}^{2}\\[3pt]
S_{339}={}&-1-x_{67}x_{69}+x_{69}+x_{459}^{2}+x_{460}^{2}+x_{461}+x_{461}^{2}-x_{462}^{2}-x_{463}^{2}\\*
&{}-x_{464}-x_{464}^{2}\\[3pt]
S_{340}={}&x_{67}x_{69}-x_{459}^{2}-x_{460}^{2}-x_{461}-x_{461}^{2}-x_{465}^{2}-x_{466}^{2}-x_{467}\\*
&{}-x_{467}^{2}\\[3pt]
S_{341}={}&x_{70}-x_{468}^{2}-x_{469}^{2}-x_{470}-x_{470}^{2}\\[3pt]
S_{342}={}&x_{74}-x_{471}^{2}-x_{472}^{2}-x_{473}-x_{473}^{2}\\[3pt]
S_{343}={}&-1+x_{72}-x_{73}-x_{474}^{2}-x_{475}^{2}-x_{476}-x_{476}^{2}\\[3pt]
S_{344}={}&x_{173}-x_{477}^{2}-x_{478}^{2}-x_{479}-x_{479}^{2}\\[3pt]
S_{345}={}&-8-x_{29}-x_{30}-x_{31}-x_{32}-x_{33}-x_{34}-x_{35}-x_{36}+8x_{73}\\*
&{}-x_{480}^{2}-x_{481}^{2}-x_{482}-x_{482}^{2}\\[3pt]
S_{346}={}&x_{175}-x_{483}^{2}-x_{484}^{2}-x_{485}-x_{485}^{2}\\[3pt]
S_{347}={}&-1+x_{174}-x_{486}^{2}-x_{487}^{2}-x_{488}-x_{488}^{2}-x_{489}^{2}-x_{490}^{2}-x_{491}-x_{491}^{2}\\[3pt]
S_{348}={}&x_{176}-x_{492}^{2}-x_{493}^{2}-x_{494}-x_{494}^{2}\\[3pt]
S_{349}={}&x_{179}-x_{495}^{2}-x_{496}^{2}-x_{497}-x_{497}^{2}\\[3pt]
S_{350}={}&-1+x_{178}-x_{498}^{2}-x_{499}^{2}-x_{500}-x_{500}^{2}-x_{501}^{2}-x_{502}^{2}-x_{503}-x_{503}^{2}\\[3pt]
S_{351}={}&x_{62}-x_{69}+x_{176}-x_{504}^{2}-x_{505}^{2}-x_{506}-x_{506}^{2}\\[3pt]
S_{352}={}&x_{181}-x_{507}^{2}-x_{508}^{2}-x_{509}-x_{509}^{2}\\[3pt]
S_{353}={}&-1+x_{178}-x_{510}^{2}-x_{511}^{2}-x_{512}-x_{512}^{2}-x_{513}^{2}-x_{514}^{2}-x_{515}-x_{515}^{2}\\[3pt]
S_{354}={}&x_{62}-x_{70}+x_{176}-x_{516}^{2}-x_{517}^{2}-x_{518}-x_{518}^{2}\\[3pt]
S_{355}={}&x_{184}-x_{519}^{2}-x_{520}^{2}-x_{521}-x_{521}^{2}\\[3pt]
S_{356}={}&-1+x_{183}-x_{522}^{2}-x_{523}^{2}-x_{524}-x_{524}^{2}-x_{525}^{2}-x_{526}^{2}-x_{527}-x_{527}^{2}\\[3pt]
S_{357}={}&x_{186}-x_{528}^{2}-x_{529}^{2}-x_{530}-x_{530}^{2}\\[3pt]
S_{358}={}&-1+x_{183}-x_{531}^{2}-x_{532}^{2}-x_{533}-x_{533}^{2}-x_{534}^{2}-x_{535}^{2}-x_{536}-x_{536}^{2}\\[3pt]
S_{359}={}&x_{187}-x_{537}^{2}-x_{538}^{2}-x_{539}-x_{539}^{2}\\[3pt]
S_{360}={}&-1+x_{183}-x_{540}^{2}-x_{541}^{2}-x_{542}-x_{542}^{2}-x_{543}^{2}-x_{544}^{2}-x_{545}-x_{545}^{2}\\[3pt]
S_{361}={}&x_{188}-x_{546}^{2}-x_{547}^{2}-x_{548}-x_{548}^{2}\\[3pt]
S_{362}={}&-1+x_{183}-x_{549}^{2}-x_{550}^{2}-x_{551}-x_{551}^{2}-x_{552}^{2}-x_{553}^{2}-x_{554}-x_{554}^{2}\\[3pt]
S_{363}={}&x_{190}-x_{555}^{2}-x_{556}^{2}-x_{557}-x_{557}^{2}\\[3pt]
S_{364}={}&-1+x_{189}-x_{558}^{2}-x_{559}^{2}-x_{560}-x_{560}^{2}-x_{561}^{2}-x_{562}^{2}-x_{563}-x_{563}^{2}\\[3pt]
S_{365}={}&x_{192}-x_{564}^{2}-x_{565}^{2}-x_{566}-x_{566}^{2}\\[3pt]
S_{366}={}&-1+x_{189}-x_{567}^{2}-x_{568}^{2}-x_{569}-x_{569}^{2}-x_{570}^{2}-x_{571}^{2}-x_{572}-x_{572}^{2}\\[3pt]
S_{367}={}&x_{194}-x_{573}^{2}-x_{574}^{2}-x_{575}-x_{575}^{2}\\[3pt]
S_{368}={}&-1+x_{193}-x_{576}^{2}-x_{577}^{2}-x_{578}-x_{578}^{2}-x_{579}^{2}-x_{580}^{2}-x_{581}-x_{581}^{2}\\[3pt]
S_{369}={}&x_{196}-x_{582}^{2}-x_{583}^{2}-x_{584}-x_{584}^{2}\\[3pt]
S_{370}={}&-1+x_{193}-x_{585}^{2}-x_{586}^{2}-x_{587}-x_{587}^{2}-x_{588}^{2}-x_{589}^{2}-x_{590}-x_{590}^{2}\\[3pt]
S_{371}={}&x_{198}-x_{591}^{2}-x_{592}^{2}-x_{593}-x_{593}^{2}\\[3pt]
S_{372}={}&x_{199}-x_{594}^{2}-x_{595}^{2}-x_{596}-x_{596}^{2}\\[3pt]
S_{373}={}&x_{203}-x_{597}^{2}-x_{598}^{2}-x_{599}-x_{599}^{2}\\[3pt]
S_{374}={}&x_{205}-x_{600}^{2}-x_{601}^{2}-x_{602}-x_{602}^{2}\\[3pt]
S_{375}={}&x_{206}-x_{603}^{2}-x_{604}^{2}-x_{605}-x_{605}^{2}\\[3pt]
S_{376}={}&x_{209}-x_{606}^{2}-x_{607}^{2}-x_{608}-x_{608}^{2}\\[3pt]
S_{377}={}&-1-x_{0}+x_{210}-x_{609}^{2}-x_{610}^{2}-x_{611}-x_{611}^{2}\\[3pt]
S_{378}={}&x_{216}-x_{612}^{2}-x_{613}^{2}-x_{614}-x_{614}^{2}\\[3pt]
S_{379}={}&x_{217}-x_{615}^{2}-x_{616}^{2}-x_{617}-x_{617}^{2}\\[3pt]
S_{380}={}&x_{218}-x_{618}^{2}-x_{619}^{2}-x_{620}-x_{620}^{2}\\[3pt]
S_{381}={}&x_{219}-x_{621}^{2}-x_{622}^{2}-x_{623}-x_{623}^{2}\\[3pt]
S_{382}={}&x_{221}-x_{624}^{2}-x_{625}^{2}-x_{626}-x_{626}^{2}\\[3pt]
S_{383}={}&x_{225}-x_{627}^{2}-x_{628}^{2}-x_{629}-x_{629}^{2}\\[3pt]
S_{384}={}&x_{228}-x_{630}^{2}-x_{631}^{2}-x_{632}-x_{632}^{2}\\[3pt]
S_{385}={}&x_{229}-x_{633}^{2}-x_{634}^{2}-x_{635}-x_{635}^{2}\\[3pt]
S_{386}={}&x_{230}-x_{636}^{2}-x_{637}^{2}-x_{638}-x_{638}^{2}\\[3pt]
S_{387}={}&x_{234}-x_{639}^{2}-x_{640}^{2}-x_{641}-x_{641}^{2}\\[3pt]
S_{388}={}&x_{235}-x_{642}^{2}-x_{643}^{2}-x_{644}-x_{644}^{2}\\[3pt]
S_{389}={}&x_{236}-x_{645}^{2}-x_{646}^{2}-x_{647}-x_{647}^{2}\\[3pt]
S_{390}={}&x_{237}-x_{648}^{2}-x_{649}^{2}-x_{650}-x_{650}^{2}\\[3pt]
S_{391}={}&x_{238}-x_{651}^{2}-x_{652}^{2}-x_{653}-x_{653}^{2}\\[3pt]
S_{392}={}&x_{239}-x_{654}^{2}-x_{655}^{2}-x_{656}-x_{656}^{2}\\[3pt]
S_{393}={}&x_{240}-x_{657}^{2}-x_{658}^{2}-x_{659}-x_{659}^{2}\\[3pt]
S_{394}={}&x_{241}-x_{660}^{2}-x_{661}^{2}-x_{662}-x_{662}^{2}\\[3pt]
S_{395}={}&x_{245}-x_{663}^{2}-x_{664}^{2}-x_{665}-x_{665}^{2}\\[3pt]
S_{396}={}&x_{246}-x_{666}^{2}-x_{667}^{2}-x_{668}-x_{668}^{2}\\[3pt]
S_{397}={}&x_{247}-x_{669}^{2}-x_{670}^{2}-x_{671}-x_{671}^{2}\\[3pt]
S_{398}={}&x_{248}-x_{672}^{2}-x_{673}^{2}-x_{674}-x_{674}^{2}\\[3pt]
S_{399}={}&x_{249}-x_{675}^{2}-x_{676}^{2}-x_{677}-x_{677}^{2}\\[3pt]
S_{400}={}&-x_{61}-x_{222}+x_{225}+x_{250}-x_{678}^{2}-x_{679}^{2}-x_{680}-x_{680}^{2}\\[3pt]
S_{401}={}&x_{255}-x_{681}^{2}-x_{682}^{2}-x_{683}-x_{683}^{2}\\[3pt]
S_{402}={}&x_{256}-x_{684}^{2}-x_{685}^{2}-x_{686}-x_{686}^{2}\\[3pt]
S_{403}={}&x_{257}-x_{687}^{2}-x_{688}^{2}-x_{689}-x_{689}^{2}\\[3pt]
S_{404}={}&x_{258}-x_{690}^{2}-x_{691}^{2}-x_{692}-x_{692}^{2}\\[3pt]
S_{405}={}&x_{259}-x_{693}^{2}-x_{694}^{2}-x_{695}-x_{695}^{2}\\[3pt]
S_{406}={}&x_{260}-x_{696}^{2}-x_{697}^{2}-x_{698}-x_{698}^{2}\\[3pt]
S_{407}={}&x_{261}-x_{699}^{2}-x_{700}^{2}-x_{701}-x_{701}^{2}\\[3pt]
S_{408}={}&-1-x_{0}-x_{68}+x_{262}-x_{702}^{2}-x_{703}^{2}-x_{704}-x_{704}^{2}\\[3pt]
S_{409}={}&x_{267}-x_{705}^{2}-x_{706}^{2}-x_{707}-x_{707}^{2}\\[3pt]
S_{410}={}&x_{268}-x_{708}^{2}-x_{709}^{2}-x_{710}-x_{710}^{2}\\[3pt]
S_{411}={}&x_{269}-x_{711}^{2}-x_{712}^{2}-x_{713}-x_{713}^{2}\\[3pt]
S_{412}={}&x_{270}-x_{714}^{2}-x_{715}^{2}-x_{716}-x_{716}^{2}\\[3pt]
S_{413}={}&x_{271}-x_{717}^{2}-x_{718}^{2}-x_{719}-x_{719}^{2}\\[3pt]
S_{414}={}&-3-x_{0}-x_{65}+x_{272}-x_{720}^{2}-x_{721}^{2}-x_{722}-x_{722}^{2}\\[3pt]
S_{415}={}&x_{277}-x_{723}^{2}-x_{724}^{2}-x_{725}-x_{725}^{2}\\[3pt]
S_{416}={}&x_{278}-x_{726}^{2}-x_{727}^{2}-x_{728}-x_{728}^{2}\\[3pt]
S_{417}={}&x_{279}-x_{729}^{2}-x_{730}^{2}-x_{731}-x_{731}^{2}\\[3pt]
S_{418}={}&x_{280}-x_{732}^{2}-x_{733}^{2}-x_{734}-x_{734}^{2}\\[3pt]
S_{419}={}&-1-x_{29}+x_{281}-x_{735}^{2}-x_{736}^{2}-x_{737}-x_{737}^{2}\\[3pt]
S_{420}={}&x_{286}-x_{738}^{2}-x_{739}^{2}-x_{740}-x_{740}^{2}\\[3pt]
S_{421}={}&x_{287}-x_{741}^{2}-x_{742}^{2}-x_{743}-x_{743}^{2}\\[3pt]
S_{422}={}&x_{288}-x_{744}^{2}-x_{745}^{2}-x_{746}-x_{746}^{2}\\[3pt]
S_{423}={}&x_{289}-x_{747}^{2}-x_{748}^{2}-x_{749}-x_{749}^{2}\\[3pt]
S_{424}={}&x_{290}-x_{750}^{2}-x_{751}^{2}-x_{752}-x_{752}^{2}\\[3pt]
S_{425}={}&-1-x_{30}+x_{291}-x_{753}^{2}-x_{754}^{2}-x_{755}-x_{755}^{2}\\[3pt]
S_{426}={}&x_{296}-x_{756}^{2}-x_{757}^{2}-x_{758}-x_{758}^{2}\\[3pt]
S_{427}={}&x_{297}-x_{759}^{2}-x_{760}^{2}-x_{761}-x_{761}^{2}\\[3pt]
S_{428}={}&x_{298}-x_{762}^{2}-x_{763}^{2}-x_{764}-x_{764}^{2}\\[3pt]
S_{429}={}&x_{299}-x_{765}^{2}-x_{766}^{2}-x_{767}-x_{767}^{2}\\[3pt]
S_{430}={}&-1-x_{31}+x_{300}-x_{768}^{2}-x_{769}^{2}-x_{770}-x_{770}^{2}\\[3pt]
S_{431}={}&x_{305}-x_{771}^{2}-x_{772}^{2}-x_{773}-x_{773}^{2}\\[3pt]
S_{432}={}&x_{306}-x_{774}^{2}-x_{775}^{2}-x_{776}-x_{776}^{2}\\[3pt]
S_{433}={}&x_{307}-x_{777}^{2}-x_{778}^{2}-x_{779}-x_{779}^{2}\\[3pt]
S_{434}={}&x_{308}-x_{780}^{2}-x_{781}^{2}-x_{782}-x_{782}^{2}\\[3pt]
S_{435}={}&-1-x_{32}+x_{309}-x_{783}^{2}-x_{784}^{2}-x_{785}-x_{785}^{2}\\[3pt]
S_{436}={}&x_{314}-x_{786}^{2}-x_{787}^{2}-x_{788}-x_{788}^{2}\\[3pt]
S_{437}={}&x_{315}-x_{789}^{2}-x_{790}^{2}-x_{791}-x_{791}^{2}\\[3pt]
S_{438}={}&x_{316}-x_{792}^{2}-x_{793}^{2}-x_{794}-x_{794}^{2}\\[3pt]
S_{439}={}&x_{317}-x_{795}^{2}-x_{796}^{2}-x_{797}-x_{797}^{2}\\[3pt]
S_{440}={}&-1-x_{33}+x_{318}-x_{798}^{2}-x_{799}^{2}-x_{800}-x_{800}^{2}\\[3pt]
S_{441}={}&x_{323}-x_{801}^{2}-x_{802}^{2}-x_{803}-x_{803}^{2}\\[3pt]
S_{442}={}&x_{324}-x_{804}^{2}-x_{805}^{2}-x_{806}-x_{806}^{2}\\[3pt]
S_{443}={}&x_{325}-x_{807}^{2}-x_{808}^{2}-x_{809}-x_{809}^{2}\\[3pt]
S_{444}={}&x_{326}-x_{810}^{2}-x_{811}^{2}-x_{812}-x_{812}^{2}\\[3pt]
S_{445}={}&-1-x_{2}+x_{327}-x_{813}^{2}-x_{814}^{2}-x_{815}-x_{815}^{2}\\[3pt]
S_{446}={}&x_{332}-x_{816}^{2}-x_{817}^{2}-x_{818}-x_{818}^{2}\\[3pt]
S_{447}={}&x_{333}-x_{819}^{2}-x_{820}^{2}-x_{821}-x_{821}^{2}\\[3pt]
S_{448}={}&x_{334}-x_{822}^{2}-x_{823}^{2}-x_{824}-x_{824}^{2}\\[3pt]
S_{449}={}&x_{335}-x_{825}^{2}-x_{826}^{2}-x_{827}-x_{827}^{2}\\[3pt]
S_{450}={}&-1-x_{34}+x_{336}-x_{828}^{2}-x_{829}^{2}-x_{830}-x_{830}^{2}\\[3pt]
S_{451}={}&x_{341}-x_{831}^{2}-x_{832}^{2}-x_{833}-x_{833}^{2}\\[3pt]
S_{452}={}&x_{342}-x_{834}^{2}-x_{835}^{2}-x_{836}-x_{836}^{2}\\[3pt]
S_{453}={}&x_{343}-x_{837}^{2}-x_{838}^{2}-x_{839}-x_{839}^{2}\\[3pt]
S_{454}={}&x_{344}-x_{840}^{2}-x_{841}^{2}-x_{842}-x_{842}^{2}\\[3pt]
S_{455}={}&-1-x_{35}+x_{345}-x_{843}^{2}-x_{844}^{2}-x_{845}-x_{845}^{2}\\[3pt]
S_{456}={}&x_{350}-x_{846}^{2}-x_{847}^{2}-x_{848}-x_{848}^{2}\\[3pt]
S_{457}={}&x_{351}-x_{849}^{2}-x_{850}^{2}-x_{851}-x_{851}^{2}\\[3pt]
S_{458}={}&x_{352}-x_{852}^{2}-x_{853}^{2}-x_{854}-x_{854}^{2}\\[3pt]
S_{459}={}&x_{353}-x_{855}^{2}-x_{856}^{2}-x_{857}-x_{857}^{2}\\[3pt]
S_{460}={}&-1-x_{1}+x_{354}-x_{858}^{2}-x_{859}^{2}-x_{860}-x_{860}^{2}\\[3pt]
S_{461}={}&x_{359}-x_{861}^{2}-x_{862}^{2}-x_{863}-x_{863}^{2}\\[3pt]
S_{462}={}&x_{360}-x_{864}^{2}-x_{865}^{2}-x_{866}-x_{866}^{2}\\[3pt]
S_{463}={}&x_{361}-x_{867}^{2}-x_{868}^{2}-x_{869}-x_{869}^{2}\\[3pt]
S_{464}={}&x_{362}-x_{870}^{2}-x_{871}^{2}-x_{872}-x_{872}^{2}\\[3pt]
S_{465}={}&-1-x_{36}+x_{363}-x_{873}^{2}-x_{874}^{2}-x_{875}-x_{875}^{2}\\[3pt]
S_{466}={}&x_{368}-x_{876}^{2}-x_{877}^{2}-x_{878}-x_{878}^{2}\\[3pt]
S_{467}={}&x_{369}-x_{879}^{2}-x_{880}^{2}-x_{881}-x_{881}^{2}\\[3pt]
S_{468}={}&x_{370}-x_{882}^{2}-x_{883}^{2}-x_{884}-x_{884}^{2}\\[3pt]
S_{469}={}&x_{385}-x_{885}^{2}-x_{886}^{2}-x_{887}-x_{887}^{2}\\[3pt]
\end{align*}
\endgroup


\clearpage
\section{The complete small-coefficient 917-integer equation}\label{app:integer917}

Here every $x_i$ ranges independently over $\Z$, and the indices run
from 0 to 916. The equation is
\[
 \boxed{\displaystyle Q_{917}(x_0,\ldots,x_{916})
                  =\sum_{j=0}^{498}T_j(x_0,\ldots,x_{916})^2=0.}
\]
The following complete list corresponds to
\file{equations/integer917/equation.json}. Its expanded polynomial has
6,986 monomials and coefficient height 63,906,036.

% Generated from the checked literal residual data.
\begingroup\small
\allowdisplaybreaks[4]
\setlength{\jot}{2pt}
\begin{align*}
T_{0}={}&-x_{4}^{2}+x_{5}\\[3pt]
T_{1}={}&-x_{5}^{2}+x_{6}\\[3pt]
T_{2}={}&-x_{6}^{2}+x_{7}\\[3pt]
T_{3}={}&-x_{7}^{2}+x_{8}\\[3pt]
T_{4}={}&-x_{8}^{2}+x_{9}\\[3pt]
T_{5}={}&-x_{9}^{2}+x_{10}\\[3pt]
T_{6}={}&-x_{10}^{2}+x_{11}\\[3pt]
T_{7}={}&-x_{11}^{2}+x_{12}\\[3pt]
T_{8}={}&-x_{12}^{2}+x_{13}\\[3pt]
T_{9}={}&-x_{13}^{2}+x_{14}\\[3pt]
T_{10}={}&-x_{14}^{2}+x_{15}\\[3pt]
T_{11}={}&-x_{15}^{2}+x_{16}\\[3pt]
T_{12}={}&-x_{16}^{2}+x_{17}\\[3pt]
T_{13}={}&-x_{17}^{2}+x_{18}\\[3pt]
T_{14}={}&-x_{18}^{2}+x_{19}\\[3pt]
T_{15}={}&-x_{19}^{2}+x_{20}\\[3pt]
T_{16}={}&-x_{20}^{2}+x_{21}\\[3pt]
T_{17}={}&-x_{21}^{2}+x_{22}\\[3pt]
T_{18}={}&-x_{22}^{2}+x_{23}\\[3pt]
T_{19}={}&-x_{23}^{2}+x_{24}\\[3pt]
T_{20}={}&-x_{24}^{2}+x_{25}\\[3pt]
T_{21}={}&1+x_{3}-x_{25}^{2}\\[3pt]
T_{22}={}&-x_{0}x_{2}-2x_{2}+x_{28}\\[3pt]
T_{23}={}&-2+x_{1}-x_{2}x_{26}-x_{3}-x_{26}\\[3pt]
T_{24}={}&-1+x_{1}-x_{27}-x_{27}x_{28}\\[3pt]
T_{25}={}&-3-x_{0}-x_{1}-x_{2}-x_{3}+x_{38}\\[3pt]
T_{26}={}&-x_{38}^{2}+x_{39}\\[3pt]
T_{27}={}&-x_{39}^{2}+x_{40}\\[3pt]
T_{28}={}&-x_{40}^{2}+x_{41}\\[3pt]
T_{29}={}&-x_{41}^{2}+x_{42}\\[3pt]
T_{30}={}&-x_{42}^{2}+x_{43}\\[3pt]
T_{31}={}&-x_{43}^{2}+x_{44}\\[3pt]
T_{32}={}&-x_{44}^{2}+x_{45}\\[3pt]
T_{33}={}&-x_{45}^{2}+x_{46}\\[3pt]
T_{34}={}&-x_{46}^{2}+x_{47}\\[3pt]
T_{35}={}&-x_{47}^{2}+x_{48}\\[3pt]
T_{36}={}&-x_{48}^{2}+x_{49}\\[3pt]
T_{37}={}&-x_{49}^{2}+x_{50}\\[3pt]
T_{38}={}&-x_{50}^{2}+x_{51}\\[3pt]
T_{39}={}&-x_{51}^{2}+x_{52}\\[3pt]
T_{40}={}&-x_{52}^{2}+x_{53}\\[3pt]
T_{41}={}&-x_{53}^{2}+x_{54}\\[3pt]
T_{42}={}&-x_{54}^{2}+x_{55}\\[3pt]
T_{43}={}&-x_{55}^{2}+x_{56}\\[3pt]
T_{44}={}&-x_{56}^{2}+x_{57}\\[3pt]
T_{45}={}&-x_{57}^{2}+x_{58}\\[3pt]
T_{46}={}&-x_{58}^{2}+x_{59}\\[3pt]
T_{47}={}&-x_{59}^{2}+x_{60}\\[3pt]
T_{48}={}&-x_{60}^{2}+x_{236}-x_{239}\\[3pt]
T_{49}={}&-x_{0}x_{61}-2x_{61}+x_{65}\\[3pt]
T_{50}={}&-x_{67}x_{68}+x_{72}\\[3pt]
T_{51}={}&-x_{71}x_{72}+x_{73}x_{75}+x_{74}\\[3pt]
T_{52}={}&1-2x_{61}-x_{61}x_{66}+x_{73}\\[3pt]
T_{53}={}&-1-x_{37}x_{61}-x_{61}+x_{74}\\[3pt]
T_{54}={}&-6-5x_{29}-x_{29}^{2}-2x_{29}x_{30}-5x_{30}-x_{30}^{2}+x_{76}\\[3pt]
T_{55}={}&-72-68x_{30}-16x_{30}^{2}-16x_{30}x_{76}-34x_{76}-4x_{76}^{2}+x_{77}\\[3pt]
T_{56}={}&-30-11x_{29}-x_{29}^{2}-4x_{29}x_{77}-22x_{77}-4x_{77}^{2}+x_{78}\\[3pt]
T_{57}={}&-380-312x_{77}-64x_{77}^{2}-32x_{77}x_{78}-78x_{78}-4x_{78}^{2}+x_{79}\\[3pt]
T_{58}={}&-6-5x_{30}-x_{30}^{2}-2x_{30}x_{31}-5x_{31}-x_{31}^{2}+x_{80}\\[3pt]
T_{59}={}&-30-11x_{29}-x_{29}^{2}-8x_{29}x_{31}-4x_{29}x_{80}-44x_{31}-16x_{31}^{2}\\*
&{}-16x_{31}x_{80}-22x_{80}-4x_{80}^{2}+x_{81}\\[3pt]
T_{60}={}&-6-5x_{29}-x_{29}^{2}-2x_{29}x_{31}-5x_{31}-x_{31}^{2}+x_{82}\\[3pt]
T_{61}={}&-12-4x_{30}-4x_{31}-2x_{76}-2x_{82}+x_{83}\\[3pt]
T_{62}={}&-42-32x_{31}-8x_{80}-2x_{81}-8x_{82}-2x_{83}-2x_{83}^{2}+x_{84}\\[3pt]
T_{63}={}&-96-64x_{31}-32x_{82}-8x_{83}-8x_{83}^{2}-2x_{84}-2x_{84}^{2}+x_{85}\\[3pt]
T_{64}={}&-10-4x_{29}-4x_{30}-2x_{80}-2x_{82}+x_{86}\\[3pt]
T_{65}={}&-27-20x_{30}-2x_{76}-8x_{80}-2x_{86}-2x_{86}^{2}+x_{87}\\[3pt]
T_{66}={}&-80-64x_{30}-32x_{80}-8x_{86}-8x_{86}^{2}-2x_{87}-2x_{87}^{2}+x_{88}\\[3pt]
T_{67}={}&-56-15x_{29}-x_{29}^{2}-8x_{29}x_{31}-4x_{29}x_{80}-60x_{31}-16x_{31}^{2}\\*
&{}-16x_{31}x_{80}-30x_{80}-4x_{80}^{2}+x_{89}\\[3pt]
T_{68}={}&-56-15x_{30}-x_{30}^{2}-8x_{30}x_{31}-4x_{30}x_{82}-60x_{31}-16x_{31}^{2}\\*
&{}-16x_{31}x_{82}-30x_{82}-4x_{82}^{2}+x_{90}\\[3pt]
T_{69}={}&-52-32x_{31}-8x_{80}-8x_{82}-2x_{89}-2x_{90}+x_{91}\\[3pt]
T_{70}={}&-104-64x_{31}-32x_{82}-8x_{90}-2x_{91}-2x_{91}^{2}+x_{92}\\[3pt]
T_{71}={}&-110-168x_{30}-64x_{30}^{2}-32x_{30}x_{76}-32x_{30}x_{80}-42x_{76}\\*
&{}-4x_{76}^{2}-8x_{76}x_{80}-42x_{80}-4x_{80}^{2}+x_{93}\\[3pt]
T_{72}={}&-27-4x_{29}-16x_{30}-8x_{76}-2x_{82}-2x_{93}+x_{94}\\[3pt]
T_{73}={}&-80-64x_{30}-32x_{76}-8x_{93}-2x_{94}-2x_{94}^{2}+x_{95}\\[3pt]
T_{74}={}&-56-60x_{30}-16x_{30}^{2}-16x_{30}x_{76}-30x_{76}-4x_{76}^{2}+x_{96}\\[3pt]
T_{75}={}&-30-11x_{29}-x_{29}^{2}-4x_{29}x_{96}-22x_{96}-4x_{96}^{2}+x_{97}\\[3pt]
T_{76}={}&-380-312x_{96}-64x_{96}^{2}-32x_{96}x_{97}-78x_{97}-4x_{97}^{2}+x_{98}\\[3pt]
T_{77}={}&-30-11x_{30}-x_{30}^{2}-4x_{30}x_{98}-22x_{98}-4x_{98}^{2}+x_{99}\\[3pt]
T_{78}={}&-156-200x_{30}-64x_{30}^{2}-64x_{30}x_{76}-100x_{76}-16x_{76}^{2}+x_{100}\\[3pt]
T_{79}={}&-37-20x_{30}-8x_{76}-2x_{80}-2x_{100}+x_{101}\\[3pt]
T_{80}={}&-120-64x_{30}-32x_{76}-8x_{100}-2x_{101}-2x_{101}^{2}+x_{102}\\[3pt]
T_{81}={}&-56-60x_{29}-16x_{29}^{2}-16x_{29}x_{82}-8x_{29}x_{102}-30x_{82}-4x_{82}^{2}\\*
&{}-4x_{82}x_{102}-15x_{102}-x_{102}^{2}+x_{103}\\[3pt]
T_{82}={}&-6-5x_{31}-x_{31}^{2}-2x_{31}x_{32}-5x_{32}-x_{32}^{2}+x_{104}\\[3pt]
T_{83}={}&-30-11x_{29}-x_{29}^{2}-8x_{29}x_{32}-4x_{29}x_{104}-44x_{32}-16x_{32}^{2}\\*
&{}-16x_{32}x_{104}-22x_{104}-4x_{104}^{2}+x_{105}\\[3pt]
T_{84}={}&-72-136x_{31}-64x_{31}^{2}-32x_{31}x_{80}-32x_{31}x_{82}-34x_{80}-4x_{80}^{2}\\*
&{}-8x_{80}x_{82}-34x_{82}-4x_{82}^{2}+x_{106}\\[3pt]
T_{85}={}&-20-4x_{29}-16x_{31}-8x_{80}-2x_{105}-2x_{106}+x_{107}\\[3pt]
T_{86}={}&-64-64x_{31}-32x_{80}-8x_{106}-2x_{107}-2x_{107}^{2}+x_{108}\\[3pt]
T_{87}={}&-21-16x_{30}-x_{31}-8x_{80}-2x_{86}-2x_{86}^{2}+x_{109}\\[3pt]
T_{88}={}&-80-64x_{30}-32x_{80}-8x_{86}-8x_{86}^{2}-2x_{109}-2x_{109}^{2}+x_{110}\\[3pt]
T_{89}={}&-42-52x_{31}-16x_{31}^{2}-16x_{31}x_{80}-26x_{80}-4x_{80}^{2}+x_{111}\\[3pt]
T_{90}={}&-30-11x_{31}-x_{31}^{2}-4x_{31}x_{111}-22x_{111}-4x_{111}^{2}+x_{112}\\[3pt]
T_{91}={}&-306-35x_{30}-x_{30}^{2}-16x_{30}x_{111}-4x_{30}x_{112}-280x_{111}-64x_{111}^{2}\\*
&{}-32x_{111}x_{112}-70x_{112}-4x_{112}^{2}+x_{113}\\[3pt]
T_{92}={}&-4556-4320x_{111}-1024x_{111}^{2}-512x_{111}x_{112}-128x_{111}x_{113}\\*
&{}-1080x_{112}-64x_{112}^{2}-32x_{112}x_{113}-270x_{113}-4x_{113}^{2}+x_{114}\\[3pt]
T_{93}={}&-30-11x_{31}-x_{31}^{2}-4x_{31}x_{114}-22x_{114}-4x_{114}^{2}+x_{115}\\[3pt]
T_{94}={}&-306-35x_{30}-x_{30}^{2}-16x_{30}x_{114}-4x_{30}x_{115}-280x_{114}-64x_{114}^{2}\\*
&{}-32x_{114}x_{115}-70x_{115}-4x_{115}^{2}+x_{116}\\[3pt]
T_{95}={}&-68-x_{110}-32x_{114}-8x_{115}-2x_{116}+x_{117}\\[3pt]
T_{96}={}&-264-128x_{114}-32x_{115}-8x_{116}-2x_{117}-2x_{117}^{2}+x_{118}\\[3pt]
T_{97}={}&-2-3x_{29}-x_{29}^{2}-2x_{29}x_{118}-3x_{118}-x_{118}^{2}+x_{119}\\[3pt]
T_{98}={}&-12-28x_{118}-16x_{118}^{2}-16x_{118}x_{119}-14x_{119}-4x_{119}^{2}+x_{120}\\[3pt]
T_{99}={}&-30-11x_{31}-x_{31}^{2}-4x_{31}x_{120}-22x_{120}-4x_{120}^{2}+x_{121}\\[3pt]
T_{100}={}&-306-35x_{29}-x_{29}^{2}-16x_{29}x_{120}-4x_{29}x_{121}-280x_{120}-64x_{120}^{2}\\*
&{}-32x_{120}x_{121}-70x_{121}-4x_{121}^{2}+x_{122}\\[3pt]
T_{101}={}&-4556-4320x_{120}-1024x_{120}^{2}-512x_{120}x_{121}-128x_{120}x_{122}\\*
&{}-1080x_{121}-64x_{121}^{2}-32x_{121}x_{122}-270x_{122}-4x_{122}^{2}+x_{123}\\[3pt]
T_{102}={}&-90-76x_{31}-16x_{31}^{2}-16x_{31}x_{80}-16x_{31}x_{123}-38x_{80}-4x_{80}^{2}\\*
&{}-8x_{80}x_{123}-38x_{123}-4x_{123}^{2}+x_{124}\\[3pt]
T_{103}={}&-306-35x_{30}-x_{30}^{2}-16x_{30}x_{123}-4x_{30}x_{124}-280x_{123}-64x_{123}^{2}\\*
&{}-32x_{123}x_{124}-70x_{124}-4x_{124}^{2}+x_{125}\\[3pt]
T_{104}={}&-4556-4320x_{114}-1024x_{114}^{2}-512x_{114}x_{115}-128x_{114}x_{116}\\*
&{}-1080x_{115}-64x_{115}^{2}-32x_{115}x_{116}-270x_{116}-4x_{116}^{2}+x_{126}\\[3pt]
T_{105}={}&-70-32x_{123}-8x_{124}-2x_{125}-2x_{126}+x_{127}\\[3pt]
T_{106}={}&-x_{127}-x_{127}^{2}+x_{128}\\[3pt]
T_{107}={}&-306-35x_{31}-x_{31}^{2}-16x_{31}x_{126}-4x_{31}x_{128}-280x_{126}-64x_{126}^{2}\\*
&{}-32x_{126}x_{128}-70x_{128}-4x_{128}^{2}+x_{129}\\[3pt]
T_{108}={}&-42-52x_{30}-16x_{30}^{2}-8x_{30}x_{31}-16x_{30}x_{76}-13x_{31}-x_{31}^{2}\\*
&{}-4x_{31}x_{76}-26x_{76}-4x_{76}^{2}+x_{130}\\[3pt]
T_{109}={}&-756-1100x_{30}-400x_{30}^{2}-400x_{30}x_{76}-80x_{30}x_{130}-550x_{76}\\*
&{}-100x_{76}^{2}-40x_{76}x_{130}-110x_{130}-4x_{130}^{2}+x_{131}\\[3pt]
T_{110}={}&-88-64x_{30}-32x_{76}-8x_{130}-2x_{131}+x_{132}\\[3pt]
T_{111}={}&-66-32x_{126}-8x_{128}-2x_{129}-x_{132}+x_{133}\\[3pt]
T_{112}={}&-264-128x_{126}-32x_{128}-8x_{129}-2x_{133}-2x_{133}^{2}+x_{134}\\[3pt]
T_{113}={}&-2-3x_{134}-x_{134}^{2}+x_{135}\\[3pt]
T_{114}={}&-1-x_{3}-x_{3}x_{34}-x_{34}+x_{136}\\[3pt]
T_{115}={}&-x_{33}^{2}+2x_{33}x_{34}-x_{34}^{2}+x_{137}\\[3pt]
T_{116}={}&-2x_{79}x_{136}-4x_{136}+x_{138}\\[3pt]
T_{117}={}&-x_{33}^{2}+2x_{33}x_{34}+2x_{33}x_{138}-x_{34}^{2}-2x_{34}x_{138}-x_{138}^{2}+x_{139}\\[3pt]
T_{118}={}&-x_{85}x_{136}+x_{140}\\[3pt]
T_{119}={}&-x_{33}^{2}+2x_{33}x_{34}+2x_{33}x_{140}-x_{34}^{2}-2x_{34}x_{140}-x_{140}^{2}+x_{141}\\[3pt]
T_{120}={}&-x_{88}x_{136}+x_{142}\\[3pt]
T_{121}={}&-x_{33}^{2}+2x_{33}x_{34}+2x_{33}x_{142}-x_{34}^{2}-2x_{34}x_{142}-x_{142}^{2}+x_{143}\\[3pt]
T_{122}={}&-x_{92}x_{136}+x_{144}\\[3pt]
T_{123}={}&-x_{33}^{2}+2x_{33}x_{34}+2x_{33}x_{144}-x_{34}^{2}-2x_{34}x_{144}-x_{144}^{2}+x_{145}\\[3pt]
T_{124}={}&-x_{95}x_{136}+x_{146}\\[3pt]
T_{125}={}&-x_{33}^{2}+2x_{33}x_{34}+2x_{33}x_{146}-x_{34}^{2}-2x_{34}x_{146}-x_{146}^{2}+x_{147}\\[3pt]
T_{126}={}&-4x_{30}x_{136}-2x_{99}x_{136}-4x_{136}+x_{148}\\[3pt]
T_{127}={}&-x_{33}^{2}+2x_{33}x_{34}+2x_{33}x_{148}-x_{34}^{2}-2x_{34}x_{148}-x_{148}^{2}+x_{149}\\[3pt]
T_{128}={}&-4x_{102}x_{136}-2x_{103}x_{136}+x_{150}\\[3pt]
T_{129}={}&-x_{33}^{2}+2x_{33}x_{34}+2x_{33}x_{150}-x_{34}^{2}-2x_{34}x_{150}-x_{150}^{2}+x_{151}\\[3pt]
T_{130}={}&-x_{108}x_{136}+x_{152}\\[3pt]
T_{131}={}&-x_{33}^{2}+2x_{33}x_{34}+2x_{33}x_{152}-x_{34}^{2}-2x_{34}x_{152}-x_{152}^{2}+x_{153}\\[3pt]
T_{132}={}&-x_{29}^{2}+2x_{29}x_{31}-x_{31}^{2}+x_{154}\\[3pt]
T_{133}={}&-x_{30}^{2}+2x_{30}x_{31}-x_{31}^{2}+x_{155}\\[3pt]
T_{134}={}&-x_{154}x_{155}+x_{156}\\[3pt]
T_{135}={}&-4-4x_{32}-x_{32}^{2}+x_{157}\\[3pt]
T_{136}={}&-x_{156}^{2}+2x_{156}x_{157}-x_{157}^{2}+x_{158}\\[3pt]
T_{137}={}&-2x_{135}x_{136}-4x_{136}+x_{159}\\[3pt]
T_{138}={}&-x_{33}^{2}+2x_{33}x_{34}+2x_{33}x_{159}-x_{34}^{2}-2x_{34}x_{159}-x_{159}^{2}+x_{160}\\[3pt]
T_{139}={}&-1-x_{3}-x_{3}x_{30}-x_{30}+x_{161}\\[3pt]
T_{140}={}&-1-x_{34}-x_{34}x_{161}-x_{161}+x_{162}\\[3pt]
T_{141}={}&-1-x_{3}-x_{3}x_{29}-x_{29}+x_{163}\\[3pt]
T_{142}={}&-4-4x_{3}-4x_{3}x_{30}-2x_{3}x_{76}-4x_{30}-2x_{76}+x_{164}\\[3pt]
T_{143}={}&-1-x_{33}-x_{33}x_{163}-x_{163}+x_{165}\\[3pt]
T_{144}={}&-x_{164}x_{165}-x_{165}+x_{166}\\[3pt]
T_{145}={}&-x_{162}^{2}+2x_{162}x_{166}-x_{166}^{2}+x_{167}\\[3pt]
T_{146}={}&-6-6x_{3}-4x_{3}x_{30}-2x_{3}x_{76}-4x_{30}-2x_{76}+x_{168}\\[3pt]
T_{147}={}&-1-x_{34}-x_{34}x_{168}-x_{168}+x_{169}\\[3pt]
T_{148}={}&-1-x_{33}-x_{33}x_{161}-x_{161}+x_{170}\\[3pt]
T_{149}={}&-x_{169}^{2}+2x_{169}x_{170}-x_{170}^{2}+x_{171}\\[3pt]
T_{150}={}&-x_{137}x_{139}+x_{172}\\[3pt]
T_{151}={}&-x_{141}x_{172}+x_{173}\\[3pt]
T_{152}={}&-x_{143}x_{173}+x_{174}\\[3pt]
T_{153}={}&-x_{145}x_{174}+x_{175}\\[3pt]
T_{154}={}&-x_{147}x_{175}+x_{176}\\[3pt]
T_{155}={}&-x_{149}x_{176}+x_{177}\\[3pt]
T_{156}={}&-x_{151}x_{177}+x_{178}\\[3pt]
T_{157}={}&-x_{153}x_{178}+x_{179}\\[3pt]
T_{158}={}&-x_{158}x_{179}-x_{160}x_{179}+x_{180}\\[3pt]
T_{159}={}&-x_{167}x_{180}+x_{181}\\[3pt]
T_{160}={}&-x_{2}-x_{2}x_{37}+x_{182}\\[3pt]
T_{161}={}&-1-x_{35}-x_{35}x_{182}-x_{182}+x_{183}\\[3pt]
T_{162}={}&-1+2x_{1}-x_{1}^{2}+2x_{1}x_{33}+2x_{1}x_{183}-2x_{33}-x_{33}^{2}\\*
&{}-2x_{33}x_{183}+x_{74}x_{186}-x_{171}x_{181}-2x_{183}-x_{183}^{2}-x_{185}\\[3pt]
T_{163}={}&-1-x_{2}-x_{2}x_{36}-x_{36}-x_{36}x_{182}-x_{182}+x_{184}\\[3pt]
T_{164}={}&-1+2x_{1}-x_{1}^{2}+2x_{1}x_{34}+2x_{1}x_{184}-2x_{34}-x_{34}^{2}\\*
&{}-2x_{34}x_{184}-2x_{184}-x_{184}^{2}+x_{185}\\[3pt]
T_{165}={}&-4-x_{29}-x_{30}-x_{31}-x_{32}+x_{187}\\[3pt]
T_{166}={}&-4-x_{33}-x_{34}-x_{35}-x_{36}+8x_{74}-x_{187}-x_{509}^{2}-x_{510}^{2}-x_{511}\\*
&{}-x_{511}^{2}\\[3pt]
T_{167}={}&1+x_{63}-x_{64}x_{74}-x_{74}+x_{190}\\[3pt]
T_{168}={}&-x_{63}x_{189}-x_{189}-x_{189}x_{190}+x_{191}\\[3pt]
T_{169}={}&-x_{63}x_{193}-x_{190}x_{193}-x_{193}+x_{194}\\[3pt]
T_{170}={}&-x_{63}x_{195}-x_{190}x_{195}-x_{195}+x_{196}\\[3pt]
T_{171}={}&-x_{74}x_{197}+x_{199}\\[3pt]
T_{172}={}&-x_{74}x_{203}+x_{205}\\[3pt]
T_{173}={}&-x_{74}x_{207}+x_{209}\\[3pt]
T_{174}={}&-x_{211}-x_{211}x_{213}+x_{214}\\[3pt]
T_{175}={}&-x_{214}^{2}+x_{215}\\[3pt]
T_{176}={}&1-x_{212}^{2}+x_{244}\\[3pt]
T_{177}={}&-1-2x_{211}-x_{211}^{2}+x_{216}\\[3pt]
T_{178}={}&x_{215}-x_{215}x_{216}+x_{244}\\[3pt]
T_{179}={}&-x_{218}^{2}+x_{221}\\[3pt]
T_{180}={}&1-x_{217}^{2}+x_{221}x_{244}\\[3pt]
T_{181}={}&-4x_{218}-4x_{218}x_{219}+x_{222}\\[3pt]
T_{182}={}&1-x_{212}-x_{217}x_{220}+x_{222}\\[3pt]
T_{183}={}&-x_{224}^{2}+x_{227}\\[3pt]
T_{184}={}&1-x_{223}^{2}+x_{227}x_{244}\\[3pt]
T_{185}={}&-x_{226}^{2}+x_{228}\\[3pt]
T_{186}={}&1-x_{225}^{2}-x_{228}+x_{228}x_{229}\\[3pt]
T_{187}={}&-1-2x_{222}-x_{222}^{2}+x_{229}\\[3pt]
T_{188}={}&-x_{217}x_{230}-x_{223}+x_{225}\\[3pt]
T_{189}={}&-1-x_{0}-4x_{224}x_{231}+x_{226}\\[3pt]
T_{190}={}&5+4x_{0}+x_{0}^{2}-2x_{0}x_{212}-4x_{212}+x_{232}\\[3pt]
T_{191}={}&-x_{0}x_{224}+x_{4}-x_{223}-2x_{224}+x_{232}x_{233}+x_{234}\\[3pt]
T_{192}={}&-x_{212}x_{224}+x_{234}\\[3pt]
T_{193}={}&-x_{236}^{2}+x_{240}\\[3pt]
T_{194}={}&1-x_{235}^{2}+x_{240}x_{244}\\[3pt]
T_{195}={}&-x_{238}^{2}+x_{241}\\[3pt]
T_{196}={}&1+x_{229}x_{241}-x_{237}^{2}-x_{241}\\[3pt]
T_{197}={}&-x_{217}x_{242}-x_{235}+x_{237}\\[3pt]
T_{198}={}&-x_{236}-4x_{236}x_{243}+x_{238}+x_{239}\\[3pt]
T_{199}={}&-2x_{212}x_{236}+x_{245}\\[3pt]
T_{200}={}&-2x_{212}x_{235}+x_{246}\\[3pt]
T_{201}={}&-x_{236}-2x_{236}x_{244}+x_{247}\\[3pt]
T_{202}={}&1-4x_{212}-2x_{212}x_{236}+2x_{212}x_{239}+x_{248}+x_{249}\\[3pt]
T_{203}={}&-4-4x_{236}-x_{236}^{2}+2x_{236}x_{239}+4x_{239}-x_{239}^{2}+x_{249}\\[3pt]
T_{204}={}&-x_{236}x_{246}-x_{236}x_{247}+x_{239}x_{246}+x_{239}x_{247}-2x_{246}-2x_{247}+x_{251}\\[3pt]
T_{205}={}&-x_{62}-x_{212}x_{246}-x_{212}x_{247}+x_{235}+2x_{235}x_{244}+x_{244}x_{245}\\*
&{}-x_{248}x_{250}+x_{251}\\[3pt]
T_{206}={}&2+2x_{62}+x_{62}^{2}-2x_{62}x_{212}-2x_{212}+x_{252}\\[3pt]
T_{207}={}&-x_{62}x_{236}+x_{64}+x_{212}x_{236}-x_{235}-x_{236}+x_{252}x_{253}\\[3pt]
T_{208}={}&2+2x_{65}+x_{65}^{2}-2x_{65}x_{212}-2x_{212}+x_{254}\\[3pt]
T_{209}={}&-x_{65}x_{224}+x_{66}-x_{223}-x_{224}+x_{234}+x_{254}x_{255}\\[3pt]
T_{210}={}&1+x_{61}^{2}-2x_{61}x_{212}+x_{256}\\[3pt]
T_{211}={}&-x_{61}x_{224}+x_{67}-x_{223}+x_{234}+x_{256}x_{257}\\[3pt]
T_{212}={}&-2x_{212}x_{224}+x_{258}\\[3pt]
T_{213}={}&-2x_{212}x_{223}+x_{259}\\[3pt]
T_{214}={}&-x_{224}-2x_{224}x_{244}+x_{260}\\[3pt]
T_{215}={}&10+6x_{0}+x_{0}^{2}-2x_{0}x_{212}-6x_{212}+x_{261}\\[3pt]
T_{216}={}&x_{0}x_{259}+x_{0}x_{260}-x_{69}-x_{212}x_{259}-x_{212}x_{260}+x_{223}\\*
&{}+2x_{223}x_{244}+x_{244}x_{258}+3x_{259}+3x_{260}-x_{261}x_{262}\\[3pt]
T_{217}={}&2+2x_{69}+x_{69}^{2}-2x_{69}x_{212}-2x_{212}+x_{263}\\[3pt]
T_{218}={}&-x_{69}x_{224}-x_{223}-x_{224}+x_{234}+x_{263}x_{264}+x_{488}^{2}+x_{489}^{2}+x_{490}\\*
&{}+x_{490}^{2}\\[3pt]
T_{219}={}&-x_{266}^{2}+x_{269}\\[3pt]
T_{220}={}&1+x_{244}x_{269}-x_{265}^{2}\\[3pt]
T_{221}={}&-x_{268}^{2}+x_{270}\\[3pt]
T_{222}={}&1+x_{229}x_{270}-x_{267}^{2}-x_{270}\\[3pt]
T_{223}={}&-x_{217}x_{271}-x_{265}+x_{267}\\[3pt]
T_{224}={}&-x_{62}-x_{236}+x_{239}-4x_{266}x_{272}+x_{268}\\[3pt]
T_{225}={}&1-2x_{63}x_{212}-2x_{190}x_{212}-4x_{212}+x_{273}+x_{274}\\[3pt]
T_{226}={}&-4-4x_{63}-x_{63}^{2}-2x_{63}x_{190}-4x_{190}-x_{190}^{2}+x_{274}\\[3pt]
T_{227}={}&x_{63}x_{188}-x_{63}x_{266}+x_{188}x_{191}-x_{190}x_{266}+x_{212}x_{266}-x_{265}\\*
&{}-2x_{266}+x_{273}x_{275}+x_{515}^{2}+x_{516}^{2}+x_{517}+x_{517}^{2}\\[3pt]
T_{228}={}&1-2x_{63}x_{212}-2x_{190}x_{212}-2x_{212}+x_{276}+x_{277}\\[3pt]
T_{229}={}&-1-2x_{63}-x_{63}^{2}-2x_{63}x_{190}-2x_{190}-x_{190}^{2}+x_{277}\\[3pt]
T_{230}={}&-x_{63}x_{236}+x_{188}-x_{190}x_{236}+x_{212}x_{236}-x_{235}-x_{236}+x_{276}x_{278}\\[3pt]
T_{231}={}&-x_{280}^{2}+x_{283}\\[3pt]
T_{232}={}&1+x_{244}x_{283}-x_{279}^{2}\\[3pt]
T_{233}={}&-x_{282}^{2}+x_{284}\\[3pt]
T_{234}={}&1+x_{229}x_{284}-x_{281}^{2}-x_{284}\\[3pt]
T_{235}={}&-x_{217}x_{285}-x_{279}+x_{281}\\[3pt]
T_{236}={}&-1-x_{0}-x_{69}-4x_{280}x_{286}+x_{282}\\[3pt]
T_{237}={}&-x_{63}x_{280}+x_{70}x_{192}-x_{190}x_{280}+x_{192}x_{194}+x_{212}x_{280}+x_{273}x_{287}\\*
&{}-x_{279}-2x_{280}+x_{527}^{2}+x_{528}^{2}+x_{529}+x_{529}^{2}\\[3pt]
T_{238}={}&-x_{63}x_{224}-x_{190}x_{224}+x_{192}-x_{223}-x_{224}+x_{234}+x_{276}x_{288}\\[3pt]
T_{239}={}&-x_{290}^{2}+x_{293}\\[3pt]
T_{240}={}&1+x_{244}x_{293}-x_{289}^{2}\\[3pt]
T_{241}={}&-x_{292}^{2}+x_{294}\\[3pt]
T_{242}={}&1+x_{229}x_{294}-x_{291}^{2}-x_{294}\\[3pt]
T_{243}={}&-x_{217}x_{295}-x_{289}+x_{291}\\[3pt]
T_{244}={}&-3-x_{0}-x_{66}-4x_{290}x_{296}+x_{292}\\[3pt]
T_{245}={}&-x_{63}x_{290}+x_{71}x_{192}-x_{190}x_{290}+x_{192}x_{196}+x_{212}x_{290}+x_{273}x_{297}\\*
&{}-x_{289}-2x_{290}+x_{539}^{2}+x_{540}^{2}+x_{541}+x_{541}^{2}\\[3pt]
T_{246}={}&-x_{299}^{2}+x_{302}\\[3pt]
T_{247}={}&1+x_{244}x_{302}-x_{298}^{2}\\[3pt]
T_{248}={}&-x_{301}^{2}+x_{303}\\[3pt]
T_{249}={}&1+x_{229}x_{303}-x_{300}^{2}-x_{303}\\[3pt]
T_{250}={}&-x_{217}x_{304}-x_{298}+x_{300}\\[3pt]
T_{251}={}&-1-x_{29}-4x_{299}x_{305}+x_{301}\\[3pt]
T_{252}={}&-x_{63}x_{299}-x_{190}x_{299}+x_{198}x_{199}+x_{212}x_{299}+x_{273}x_{306}-x_{298}\\*
&{}-2x_{299}+x_{551}^{2}+x_{552}^{2}+x_{553}+x_{553}^{2}\\[3pt]
T_{253}={}&-x_{63}x_{223}-x_{63}x_{234}-x_{190}x_{223}-x_{190}x_{234}+x_{197}+x_{212}x_{234}\\*
&{}-x_{223}-x_{224}x_{244}-x_{234}+x_{276}x_{307}\\[3pt]
T_{254}={}&-x_{309}^{2}+x_{312}\\[3pt]
T_{255}={}&1+x_{244}x_{312}-x_{308}^{2}\\[3pt]
T_{256}={}&-x_{311}^{2}+x_{313}\\[3pt]
T_{257}={}&1+x_{229}x_{313}-x_{310}^{2}-x_{313}\\[3pt]
T_{258}={}&-x_{217}x_{314}-x_{308}+x_{310}\\[3pt]
T_{259}={}&-1-x_{30}-4x_{309}x_{315}+x_{311}\\[3pt]
T_{260}={}&-x_{63}x_{309}-x_{190}x_{309}+x_{199}x_{200}+x_{212}x_{309}+x_{273}x_{316}-x_{308}\\*
&{}-2x_{309}+x_{560}^{2}+x_{561}^{2}+x_{562}+x_{562}^{2}\\[3pt]
T_{261}={}&-x_{318}^{2}+x_{321}\\[3pt]
T_{262}={}&1+x_{244}x_{321}-x_{317}^{2}\\[3pt]
T_{263}={}&-x_{320}^{2}+x_{322}\\[3pt]
T_{264}={}&1+x_{229}x_{322}-x_{319}^{2}-x_{322}\\[3pt]
T_{265}={}&-x_{217}x_{323}-x_{317}+x_{319}\\[3pt]
T_{266}={}&-1-x_{31}-4x_{318}x_{324}+x_{320}\\[3pt]
T_{267}={}&-x_{63}x_{318}-x_{190}x_{318}+x_{199}x_{201}+x_{212}x_{318}+x_{273}x_{325}-x_{317}\\*
&{}-2x_{318}+x_{569}^{2}+x_{570}^{2}+x_{571}+x_{571}^{2}\\[3pt]
T_{268}={}&-x_{327}^{2}+x_{330}\\[3pt]
T_{269}={}&1+x_{244}x_{330}-x_{326}^{2}\\[3pt]
T_{270}={}&-x_{329}^{2}+x_{331}\\[3pt]
T_{271}={}&1+x_{229}x_{331}-x_{328}^{2}-x_{331}\\[3pt]
T_{272}={}&-x_{217}x_{332}-x_{326}+x_{328}\\[3pt]
T_{273}={}&-1-x_{32}-4x_{327}x_{333}+x_{329}\\[3pt]
T_{274}={}&-x_{63}x_{327}-x_{190}x_{327}+x_{199}x_{202}+x_{212}x_{327}+x_{273}x_{334}-x_{326}\\*
&{}-2x_{327}+x_{578}^{2}+x_{579}^{2}+x_{580}+x_{580}^{2}\\[3pt]
T_{275}={}&-x_{336}^{2}+x_{339}\\[3pt]
T_{276}={}&1+x_{244}x_{339}-x_{335}^{2}\\[3pt]
T_{277}={}&-x_{338}^{2}+x_{340}\\[3pt]
T_{278}={}&1+x_{229}x_{340}-x_{337}^{2}-x_{340}\\[3pt]
T_{279}={}&-x_{217}x_{341}-x_{335}+x_{337}\\[3pt]
T_{280}={}&-1-x_{33}-4x_{336}x_{342}+x_{338}\\[3pt]
T_{281}={}&-x_{63}x_{336}-x_{190}x_{336}+x_{204}x_{205}+x_{212}x_{336}+x_{273}x_{343}-x_{335}\\*
&{}-2x_{336}+x_{587}^{2}+x_{588}^{2}+x_{589}+x_{589}^{2}\\[3pt]
T_{282}={}&-x_{345}^{2}+x_{348}\\[3pt]
T_{283}={}&1+x_{244}x_{348}-x_{344}^{2}\\[3pt]
T_{284}={}&-x_{347}^{2}+x_{349}\\[3pt]
T_{285}={}&1+x_{229}x_{349}-x_{346}^{2}-x_{349}\\[3pt]
T_{286}={}&-x_{217}x_{350}-x_{344}+x_{346}\\[3pt]
T_{287}={}&-1-x_{2}-4x_{345}x_{351}+x_{347}\\[3pt]
T_{288}={}&-x_{63}x_{345}-x_{190}x_{345}+x_{203}+x_{212}x_{345}+x_{276}x_{352}-x_{344}-x_{345}\\[3pt]
T_{289}={}&-x_{354}^{2}+x_{357}\\[3pt]
T_{290}={}&1+x_{244}x_{357}-x_{353}^{2}\\[3pt]
T_{291}={}&-x_{356}^{2}+x_{358}\\[3pt]
T_{292}={}&1+x_{229}x_{358}-x_{355}^{2}-x_{358}\\[3pt]
T_{293}={}&-x_{217}x_{359}-x_{353}+x_{355}\\[3pt]
T_{294}={}&-1-x_{34}-4x_{354}x_{360}+x_{356}\\[3pt]
T_{295}={}&-x_{63}x_{354}-x_{190}x_{354}+x_{205}x_{206}+x_{212}x_{354}+x_{273}x_{361}-x_{353}\\*
&{}-2x_{354}+x_{596}^{2}+x_{597}^{2}+x_{598}+x_{598}^{2}\\[3pt]
T_{296}={}&-x_{363}^{2}+x_{366}\\[3pt]
T_{297}={}&1+x_{244}x_{366}-x_{362}^{2}\\[3pt]
T_{298}={}&-x_{365}^{2}+x_{367}\\[3pt]
T_{299}={}&1+x_{229}x_{367}-x_{364}^{2}-x_{367}\\[3pt]
T_{300}={}&-x_{217}x_{368}-x_{362}+x_{364}\\[3pt]
T_{301}={}&-1-x_{35}-4x_{363}x_{369}+x_{365}\\[3pt]
T_{302}={}&-x_{63}x_{363}-x_{190}x_{363}+x_{208}x_{209}+x_{212}x_{363}+x_{273}x_{370}-x_{362}\\*
&{}-2x_{363}+x_{605}^{2}+x_{606}^{2}+x_{607}+x_{607}^{2}\\[3pt]
T_{303}={}&-x_{372}^{2}+x_{375}\\[3pt]
T_{304}={}&1+x_{244}x_{375}-x_{371}^{2}\\[3pt]
T_{305}={}&-x_{374}^{2}+x_{376}\\[3pt]
T_{306}={}&1+x_{229}x_{376}-x_{373}^{2}-x_{376}\\[3pt]
T_{307}={}&-x_{217}x_{377}-x_{371}+x_{373}\\[3pt]
T_{308}={}&-1-x_{1}-4x_{372}x_{378}+x_{374}\\[3pt]
T_{309}={}&-x_{63}x_{372}-x_{190}x_{372}+x_{207}+x_{212}x_{372}+x_{276}x_{379}-x_{371}-x_{372}\\[3pt]
T_{310}={}&-x_{381}^{2}+x_{384}\\[3pt]
T_{311}={}&1+x_{244}x_{384}-x_{380}^{2}\\[3pt]
T_{312}={}&-x_{383}^{2}+x_{385}\\[3pt]
T_{313}={}&1+x_{229}x_{385}-x_{382}^{2}-x_{385}\\[3pt]
T_{314}={}&-x_{217}x_{386}-x_{380}+x_{382}\\[3pt]
T_{315}={}&-1-x_{36}-4x_{381}x_{387}+x_{383}\\[3pt]
T_{316}={}&-x_{63}x_{381}-x_{190}x_{381}+x_{209}x_{210}+x_{212}x_{381}+x_{273}x_{388}-x_{380}\\*
&{}-2x_{381}+x_{614}^{2}+x_{615}^{2}+x_{616}+x_{616}^{2}\\[3pt]
T_{317}={}&-x_{227}x_{240}+x_{389}\\[3pt]
T_{318}={}&-x_{269}x_{389}+x_{390}\\[3pt]
T_{319}={}&-x_{283}x_{390}+x_{391}\\[3pt]
T_{320}={}&-x_{293}x_{391}+x_{392}\\[3pt]
T_{321}={}&-x_{302}x_{392}+x_{393}\\[3pt]
T_{322}={}&-x_{312}x_{393}+x_{394}\\[3pt]
T_{323}={}&-x_{321}x_{394}+x_{395}\\[3pt]
T_{324}={}&-x_{330}x_{395}+x_{396}\\[3pt]
T_{325}={}&-x_{339}x_{396}+x_{397}\\[3pt]
T_{326}={}&-x_{348}x_{397}+x_{398}\\[3pt]
T_{327}={}&-x_{357}x_{398}+x_{399}\\[3pt]
T_{328}={}&-x_{366}x_{399}+x_{400}\\[3pt]
T_{329}={}&-x_{375}x_{400}+x_{401}\\[3pt]
T_{330}={}&-x_{384}x_{401}+x_{402}\\[3pt]
T_{331}={}&x_{218}-x_{402}-x_{402}x_{403}\\[3pt]
T_{332}={}&-7-2x_{0}-x_{2}-x_{4}-x_{29}+x_{404}\\[3pt]
T_{333}={}&-3-x_{30}-x_{31}-x_{32}-x_{404}+x_{405}\\[3pt]
T_{334}={}&-3-x_{33}-x_{34}-x_{35}-x_{405}+x_{406}\\[3pt]
T_{335}={}&-4-x_{0}-x_{36}-x_{236}+x_{239}-x_{406}+x_{407}\\[3pt]
T_{336}={}&-2-x_{61}-x_{62}-x_{236}+x_{239}-x_{407}+x_{408}\\[3pt]
T_{337}={}&-1-2x_{62}-x_{64}-x_{236}+x_{239}-x_{408}+x_{409}\\[3pt]
T_{338}={}&-1-x_{65}-x_{66}-x_{67}-x_{409}+x_{410}\\[3pt]
T_{339}={}&-2-x_{0}-3x_{69}-x_{410}+x_{411}-x_{488}^{2}-x_{489}^{2}-x_{490}-x_{490}^{2}\\[3pt]
T_{340}={}&-4-x_{0}-x_{1}-x_{66}-x_{411}+x_{412}\\[3pt]
T_{341}={}&-3-2x_{63}-x_{63}x_{188}-x_{188}-x_{188}x_{191}-2x_{190}-x_{412}+x_{413}\\*
&{}-x_{515}^{2}-x_{516}^{2}-x_{517}-x_{517}^{2}\\[3pt]
T_{342}={}&-x_{70}x_{192}-x_{71}x_{192}-x_{192}-x_{192}x_{194}-x_{192}x_{196}-x_{198}x_{199}\\*
&{}-x_{413}+x_{414}-x_{527}^{2}-x_{528}^{2}-x_{529}-x_{529}^{2}-x_{539}^{2}-x_{540}^{2}-x_{541}\\*
&{}-x_{541}^{2}-x_{551}^{2}-x_{552}^{2}-x_{553}-x_{553}^{2}\\[3pt]
T_{343}={}&-x_{197}-x_{199}x_{200}-x_{199}x_{201}-x_{199}x_{202}-x_{414}+x_{415}-x_{560}^{2}\\*
&{}-x_{561}^{2}-x_{562}-x_{562}^{2}-x_{569}^{2}-x_{570}^{2}-x_{571}-x_{571}^{2}-x_{578}^{2}-x_{579}^{2}\\*
&{}-x_{580}-x_{580}^{2}\\[3pt]
T_{344}={}&-x_{203}-x_{204}x_{205}-x_{205}x_{206}-x_{207}-x_{208}x_{209}-x_{209}x_{210}+x_{211}\\*
&{}-x_{415}-x_{587}^{2}-x_{588}^{2}-x_{589}-x_{589}^{2}-x_{596}^{2}-x_{597}^{2}-x_{598}-x_{598}^{2}\\*
&{}-x_{605}^{2}-x_{606}^{2}-x_{607}-x_{607}^{2}-x_{614}^{2}-x_{615}^{2}-x_{616}-x_{616}^{2}\\[3pt]
T_{345}={}&x_{0}-x_{416}^{2}-x_{417}^{2}-x_{418}-x_{418}^{2}\\[3pt]
T_{346}={}&x_{3}-x_{419}^{2}-x_{420}^{2}-x_{421}-x_{421}^{2}\\[3pt]
T_{347}={}&x_{4}-x_{422}^{2}-x_{423}^{2}-x_{424}-x_{424}^{2}\\[3pt]
T_{348}={}&x_{26}-x_{425}^{2}-x_{426}^{2}-x_{427}-x_{427}^{2}\\[3pt]
T_{349}={}&-2+x_{2}-x_{3}-x_{428}^{2}-x_{429}^{2}-x_{430}-x_{430}^{2}\\[3pt]
T_{350}={}&x_{27}-x_{431}^{2}-x_{432}^{2}-x_{433}-x_{433}^{2}\\[3pt]
T_{351}={}&x_{29}-x_{434}^{2}-x_{435}^{2}-x_{436}-x_{436}^{2}\\[3pt]
T_{352}={}&x_{30}-x_{437}^{2}-x_{438}^{2}-x_{439}-x_{439}^{2}\\[3pt]
T_{353}={}&x_{31}-x_{440}^{2}-x_{441}^{2}-x_{442}-x_{442}^{2}\\[3pt]
T_{354}={}&x_{32}-x_{443}^{2}-x_{444}^{2}-x_{445}-x_{445}^{2}\\[3pt]
T_{355}={}&x_{33}-x_{446}^{2}-x_{447}^{2}-x_{448}-x_{448}^{2}\\[3pt]
T_{356}={}&x_{34}-x_{449}^{2}-x_{450}^{2}-x_{451}-x_{451}^{2}\\[3pt]
T_{357}={}&x_{35}-x_{452}^{2}-x_{453}^{2}-x_{454}-x_{454}^{2}\\[3pt]
T_{358}={}&x_{36}-x_{455}^{2}-x_{456}^{2}-x_{457}-x_{457}^{2}\\[3pt]
T_{359}={}&x_{37}-x_{458}^{2}-x_{459}^{2}-x_{460}-x_{460}^{2}\\[3pt]
T_{360}={}&x_{61}-x_{461}^{2}-x_{462}^{2}-x_{463}-x_{463}^{2}\\[3pt]
T_{361}={}&x_{62}-x_{464}^{2}-x_{465}^{2}-x_{466}-x_{466}^{2}\\[3pt]
T_{362}={}&x_{64}-x_{467}^{2}-x_{468}^{2}-x_{469}-x_{469}^{2}\\[3pt]
T_{363}={}&-1-x_{61}x_{63}+x_{63}+x_{64}-x_{470}^{2}-x_{471}^{2}-x_{472}-x_{472}^{2}\\[3pt]
T_{364}={}&x_{61}x_{63}-x_{64}-x_{473}^{2}-x_{474}^{2}-x_{475}-x_{475}^{2}\\[3pt]
T_{365}={}&x_{66}-x_{476}^{2}-x_{477}^{2}-x_{478}-x_{478}^{2}\\[3pt]
T_{366}={}&x_{67}-x_{479}^{2}-x_{480}^{2}-x_{481}-x_{481}^{2}\\[3pt]
T_{367}={}&x_{68}-x_{482}^{2}-x_{483}^{2}-x_{484}-x_{484}^{2}\\[3pt]
T_{368}={}&x_{69}-x_{485}^{2}-x_{486}^{2}-x_{487}-x_{487}^{2}\\[3pt]
T_{369}={}&-1-x_{68}x_{70}+x_{70}+x_{488}^{2}+x_{489}^{2}+x_{490}+x_{490}^{2}-x_{491}^{2}-x_{492}^{2}\\*
&{}-x_{493}-x_{493}^{2}\\[3pt]
T_{370}={}&x_{68}x_{70}-x_{488}^{2}-x_{489}^{2}-x_{490}-x_{490}^{2}-x_{494}^{2}-x_{495}^{2}-x_{496}\\*
&{}-x_{496}^{2}\\[3pt]
T_{371}={}&x_{71}-x_{497}^{2}-x_{498}^{2}-x_{499}-x_{499}^{2}\\[3pt]
T_{372}={}&x_{75}-x_{500}^{2}-x_{501}^{2}-x_{502}-x_{502}^{2}\\[3pt]
T_{373}={}&-1+x_{73}-x_{74}-x_{503}^{2}-x_{504}^{2}-x_{505}-x_{505}^{2}\\[3pt]
T_{374}={}&x_{186}-x_{506}^{2}-x_{507}^{2}-x_{508}-x_{508}^{2}\\[3pt]
T_{375}={}&x_{189}-x_{512}^{2}-x_{513}^{2}-x_{514}-x_{514}^{2}\\[3pt]
T_{376}={}&-1+x_{188}-x_{515}^{2}-x_{516}^{2}-x_{517}-x_{517}^{2}-x_{518}^{2}-x_{519}^{2}-x_{520}-x_{520}^{2}\\[3pt]
T_{377}={}&x_{190}-x_{521}^{2}-x_{522}^{2}-x_{523}-x_{523}^{2}\\[3pt]
T_{378}={}&x_{193}-x_{524}^{2}-x_{525}^{2}-x_{526}-x_{526}^{2}\\[3pt]
T_{379}={}&-1+x_{192}-x_{527}^{2}-x_{528}^{2}-x_{529}-x_{529}^{2}-x_{530}^{2}-x_{531}^{2}-x_{532}-x_{532}^{2}\\[3pt]
T_{380}={}&x_{63}-x_{70}+x_{190}-x_{533}^{2}-x_{534}^{2}-x_{535}-x_{535}^{2}\\[3pt]
T_{381}={}&x_{195}-x_{536}^{2}-x_{537}^{2}-x_{538}-x_{538}^{2}\\[3pt]
T_{382}={}&-1+x_{192}-x_{539}^{2}-x_{540}^{2}-x_{541}-x_{541}^{2}-x_{542}^{2}-x_{543}^{2}-x_{544}-x_{544}^{2}\\[3pt]
T_{383}={}&x_{63}-x_{71}+x_{190}-x_{545}^{2}-x_{546}^{2}-x_{547}-x_{547}^{2}\\[3pt]
T_{384}={}&x_{198}-x_{548}^{2}-x_{549}^{2}-x_{550}-x_{550}^{2}\\[3pt]
T_{385}={}&-1+x_{197}-x_{551}^{2}-x_{552}^{2}-x_{553}-x_{553}^{2}-x_{554}^{2}-x_{555}^{2}-x_{556}-x_{556}^{2}\\[3pt]
T_{386}={}&x_{200}-x_{557}^{2}-x_{558}^{2}-x_{559}-x_{559}^{2}\\[3pt]
T_{387}={}&-1+x_{197}-x_{560}^{2}-x_{561}^{2}-x_{562}-x_{562}^{2}-x_{563}^{2}-x_{564}^{2}-x_{565}-x_{565}^{2}\\[3pt]
T_{388}={}&x_{201}-x_{566}^{2}-x_{567}^{2}-x_{568}-x_{568}^{2}\\[3pt]
T_{389}={}&-1+x_{197}-x_{569}^{2}-x_{570}^{2}-x_{571}-x_{571}^{2}-x_{572}^{2}-x_{573}^{2}-x_{574}-x_{574}^{2}\\[3pt]
T_{390}={}&x_{202}-x_{575}^{2}-x_{576}^{2}-x_{577}-x_{577}^{2}\\[3pt]
T_{391}={}&-1+x_{197}-x_{578}^{2}-x_{579}^{2}-x_{580}-x_{580}^{2}-x_{581}^{2}-x_{582}^{2}-x_{583}-x_{583}^{2}\\[3pt]
T_{392}={}&x_{204}-x_{584}^{2}-x_{585}^{2}-x_{586}-x_{586}^{2}\\[3pt]
T_{393}={}&-1+x_{203}-x_{587}^{2}-x_{588}^{2}-x_{589}-x_{589}^{2}-x_{590}^{2}-x_{591}^{2}-x_{592}-x_{592}^{2}\\[3pt]
T_{394}={}&x_{206}-x_{593}^{2}-x_{594}^{2}-x_{595}-x_{595}^{2}\\[3pt]
T_{395}={}&-1+x_{203}-x_{596}^{2}-x_{597}^{2}-x_{598}-x_{598}^{2}-x_{599}^{2}-x_{600}^{2}-x_{601}-x_{601}^{2}\\[3pt]
T_{396}={}&x_{208}-x_{602}^{2}-x_{603}^{2}-x_{604}-x_{604}^{2}\\[3pt]
T_{397}={}&-1+x_{207}-x_{605}^{2}-x_{606}^{2}-x_{607}-x_{607}^{2}-x_{608}^{2}-x_{609}^{2}-x_{610}-x_{610}^{2}\\[3pt]
T_{398}={}&x_{210}-x_{611}^{2}-x_{612}^{2}-x_{613}-x_{613}^{2}\\[3pt]
T_{399}={}&-1+x_{207}-x_{614}^{2}-x_{615}^{2}-x_{616}-x_{616}^{2}-x_{617}^{2}-x_{618}^{2}-x_{619}-x_{619}^{2}\\[3pt]
T_{400}={}&x_{212}-x_{620}^{2}-x_{621}^{2}-x_{622}-x_{622}^{2}\\[3pt]
T_{401}={}&x_{213}-x_{623}^{2}-x_{624}^{2}-x_{625}-x_{625}^{2}\\[3pt]
T_{402}={}&x_{217}-x_{626}^{2}-x_{627}^{2}-x_{628}-x_{628}^{2}\\[3pt]
T_{403}={}&x_{219}-x_{629}^{2}-x_{630}^{2}-x_{631}-x_{631}^{2}\\[3pt]
T_{404}={}&x_{220}-x_{632}^{2}-x_{633}^{2}-x_{634}-x_{634}^{2}\\[3pt]
T_{405}={}&x_{223}-x_{635}^{2}-x_{636}^{2}-x_{637}-x_{637}^{2}\\[3pt]
T_{406}={}&-1-x_{0}+x_{224}-x_{638}^{2}-x_{639}^{2}-x_{640}-x_{640}^{2}\\[3pt]
T_{407}={}&x_{230}-x_{641}^{2}-x_{642}^{2}-x_{643}-x_{643}^{2}\\[3pt]
T_{408}={}&x_{231}-x_{644}^{2}-x_{645}^{2}-x_{646}-x_{646}^{2}\\[3pt]
T_{409}={}&x_{232}-x_{647}^{2}-x_{648}^{2}-x_{649}-x_{649}^{2}\\[3pt]
T_{410}={}&x_{233}-x_{650}^{2}-x_{651}^{2}-x_{652}-x_{652}^{2}\\[3pt]
T_{411}={}&x_{235}-x_{653}^{2}-x_{654}^{2}-x_{655}-x_{655}^{2}\\[3pt]
T_{412}={}&x_{239}-x_{656}^{2}-x_{657}^{2}-x_{658}-x_{658}^{2}\\[3pt]
T_{413}={}&x_{242}-x_{659}^{2}-x_{660}^{2}-x_{661}-x_{661}^{2}\\[3pt]
T_{414}={}&x_{243}-x_{662}^{2}-x_{663}^{2}-x_{664}-x_{664}^{2}\\[3pt]
T_{415}={}&x_{244}-x_{665}^{2}-x_{666}^{2}-x_{667}-x_{667}^{2}\\[3pt]
T_{416}={}&x_{248}-x_{668}^{2}-x_{669}^{2}-x_{670}-x_{670}^{2}\\[3pt]
T_{417}={}&x_{250}-x_{671}^{2}-x_{672}^{2}-x_{673}-x_{673}^{2}\\[3pt]
T_{418}={}&x_{252}-x_{674}^{2}-x_{675}^{2}-x_{676}-x_{676}^{2}\\[3pt]
T_{419}={}&x_{253}-x_{677}^{2}-x_{678}^{2}-x_{679}-x_{679}^{2}\\[3pt]
T_{420}={}&x_{254}-x_{680}^{2}-x_{681}^{2}-x_{682}-x_{682}^{2}\\[3pt]
T_{421}={}&x_{255}-x_{683}^{2}-x_{684}^{2}-x_{685}-x_{685}^{2}\\[3pt]
T_{422}={}&x_{256}-x_{686}^{2}-x_{687}^{2}-x_{688}-x_{688}^{2}\\[3pt]
T_{423}={}&x_{257}-x_{689}^{2}-x_{690}^{2}-x_{691}-x_{691}^{2}\\[3pt]
T_{424}={}&x_{261}-x_{692}^{2}-x_{693}^{2}-x_{694}-x_{694}^{2}\\[3pt]
T_{425}={}&x_{262}-x_{695}^{2}-x_{696}^{2}-x_{697}-x_{697}^{2}\\[3pt]
T_{426}={}&x_{263}-x_{698}^{2}-x_{699}^{2}-x_{700}-x_{700}^{2}\\[3pt]
T_{427}={}&x_{264}-x_{701}^{2}-x_{702}^{2}-x_{703}-x_{703}^{2}\\[3pt]
T_{428}={}&x_{265}-x_{704}^{2}-x_{705}^{2}-x_{706}-x_{706}^{2}\\[3pt]
T_{429}={}&-x_{62}-x_{236}+x_{239}+x_{266}-x_{707}^{2}-x_{708}^{2}-x_{709}-x_{709}^{2}\\[3pt]
T_{430}={}&x_{271}-x_{710}^{2}-x_{711}^{2}-x_{712}-x_{712}^{2}\\[3pt]
T_{431}={}&x_{272}-x_{713}^{2}-x_{714}^{2}-x_{715}-x_{715}^{2}\\[3pt]
T_{432}={}&x_{273}-x_{716}^{2}-x_{717}^{2}-x_{718}-x_{718}^{2}\\[3pt]
T_{433}={}&x_{275}-x_{719}^{2}-x_{720}^{2}-x_{721}-x_{721}^{2}\\[3pt]
T_{434}={}&x_{276}-x_{722}^{2}-x_{723}^{2}-x_{724}-x_{724}^{2}\\[3pt]
T_{435}={}&x_{278}-x_{725}^{2}-x_{726}^{2}-x_{727}-x_{727}^{2}\\[3pt]
T_{436}={}&x_{279}-x_{728}^{2}-x_{729}^{2}-x_{730}-x_{730}^{2}\\[3pt]
T_{437}={}&-1-x_{0}-x_{69}+x_{280}-x_{731}^{2}-x_{732}^{2}-x_{733}-x_{733}^{2}\\[3pt]
T_{438}={}&x_{285}-x_{734}^{2}-x_{735}^{2}-x_{736}-x_{736}^{2}\\[3pt]
T_{439}={}&x_{286}-x_{737}^{2}-x_{738}^{2}-x_{739}-x_{739}^{2}\\[3pt]
T_{440}={}&x_{287}-x_{740}^{2}-x_{741}^{2}-x_{742}-x_{742}^{2}\\[3pt]
T_{441}={}&x_{288}-x_{743}^{2}-x_{744}^{2}-x_{745}-x_{745}^{2}\\[3pt]
T_{442}={}&x_{289}-x_{746}^{2}-x_{747}^{2}-x_{748}-x_{748}^{2}\\[3pt]
T_{443}={}&-3-x_{0}-x_{66}+x_{290}-x_{749}^{2}-x_{750}^{2}-x_{751}-x_{751}^{2}\\[3pt]
T_{444}={}&x_{295}-x_{752}^{2}-x_{753}^{2}-x_{754}-x_{754}^{2}\\[3pt]
T_{445}={}&x_{296}-x_{755}^{2}-x_{756}^{2}-x_{757}-x_{757}^{2}\\[3pt]
T_{446}={}&x_{297}-x_{758}^{2}-x_{759}^{2}-x_{760}-x_{760}^{2}\\[3pt]
T_{447}={}&x_{298}-x_{761}^{2}-x_{762}^{2}-x_{763}-x_{763}^{2}\\[3pt]
T_{448}={}&-1-x_{29}+x_{299}-x_{764}^{2}-x_{765}^{2}-x_{766}-x_{766}^{2}\\[3pt]
T_{449}={}&x_{304}-x_{767}^{2}-x_{768}^{2}-x_{769}-x_{769}^{2}\\[3pt]
T_{450}={}&x_{305}-x_{770}^{2}-x_{771}^{2}-x_{772}-x_{772}^{2}\\[3pt]
T_{451}={}&x_{306}-x_{773}^{2}-x_{774}^{2}-x_{775}-x_{775}^{2}\\[3pt]
T_{452}={}&x_{307}-x_{776}^{2}-x_{777}^{2}-x_{778}-x_{778}^{2}\\[3pt]
T_{453}={}&x_{308}-x_{779}^{2}-x_{780}^{2}-x_{781}-x_{781}^{2}\\[3pt]
T_{454}={}&-1-x_{30}+x_{309}-x_{782}^{2}-x_{783}^{2}-x_{784}-x_{784}^{2}\\[3pt]
T_{455}={}&x_{314}-x_{785}^{2}-x_{786}^{2}-x_{787}-x_{787}^{2}\\[3pt]
T_{456}={}&x_{315}-x_{788}^{2}-x_{789}^{2}-x_{790}-x_{790}^{2}\\[3pt]
T_{457}={}&x_{316}-x_{791}^{2}-x_{792}^{2}-x_{793}-x_{793}^{2}\\[3pt]
T_{458}={}&x_{317}-x_{794}^{2}-x_{795}^{2}-x_{796}-x_{796}^{2}\\[3pt]
T_{459}={}&-1-x_{31}+x_{318}-x_{797}^{2}-x_{798}^{2}-x_{799}-x_{799}^{2}\\[3pt]
T_{460}={}&x_{323}-x_{800}^{2}-x_{801}^{2}-x_{802}-x_{802}^{2}\\[3pt]
T_{461}={}&x_{324}-x_{803}^{2}-x_{804}^{2}-x_{805}-x_{805}^{2}\\[3pt]
T_{462}={}&x_{325}-x_{806}^{2}-x_{807}^{2}-x_{808}-x_{808}^{2}\\[3pt]
T_{463}={}&x_{326}-x_{809}^{2}-x_{810}^{2}-x_{811}-x_{811}^{2}\\[3pt]
T_{464}={}&-1-x_{32}+x_{327}-x_{812}^{2}-x_{813}^{2}-x_{814}-x_{814}^{2}\\[3pt]
T_{465}={}&x_{332}-x_{815}^{2}-x_{816}^{2}-x_{817}-x_{817}^{2}\\[3pt]
T_{466}={}&x_{333}-x_{818}^{2}-x_{819}^{2}-x_{820}-x_{820}^{2}\\[3pt]
T_{467}={}&x_{334}-x_{821}^{2}-x_{822}^{2}-x_{823}-x_{823}^{2}\\[3pt]
T_{468}={}&x_{335}-x_{824}^{2}-x_{825}^{2}-x_{826}-x_{826}^{2}\\[3pt]
T_{469}={}&-1-x_{33}+x_{336}-x_{827}^{2}-x_{828}^{2}-x_{829}-x_{829}^{2}\\[3pt]
T_{470}={}&x_{341}-x_{830}^{2}-x_{831}^{2}-x_{832}-x_{832}^{2}\\[3pt]
T_{471}={}&x_{342}-x_{833}^{2}-x_{834}^{2}-x_{835}-x_{835}^{2}\\[3pt]
T_{472}={}&x_{343}-x_{836}^{2}-x_{837}^{2}-x_{838}-x_{838}^{2}\\[3pt]
T_{473}={}&x_{344}-x_{839}^{2}-x_{840}^{2}-x_{841}-x_{841}^{2}\\[3pt]
T_{474}={}&-1-x_{2}+x_{345}-x_{842}^{2}-x_{843}^{2}-x_{844}-x_{844}^{2}\\[3pt]
T_{475}={}&x_{350}-x_{845}^{2}-x_{846}^{2}-x_{847}-x_{847}^{2}\\[3pt]
T_{476}={}&x_{351}-x_{848}^{2}-x_{849}^{2}-x_{850}-x_{850}^{2}\\[3pt]
T_{477}={}&x_{352}-x_{851}^{2}-x_{852}^{2}-x_{853}-x_{853}^{2}\\[3pt]
T_{478}={}&x_{353}-x_{854}^{2}-x_{855}^{2}-x_{856}-x_{856}^{2}\\[3pt]
T_{479}={}&-1-x_{34}+x_{354}-x_{857}^{2}-x_{858}^{2}-x_{859}-x_{859}^{2}\\[3pt]
T_{480}={}&x_{359}-x_{860}^{2}-x_{861}^{2}-x_{862}-x_{862}^{2}\\[3pt]
T_{481}={}&x_{360}-x_{863}^{2}-x_{864}^{2}-x_{865}-x_{865}^{2}\\[3pt]
T_{482}={}&x_{361}-x_{866}^{2}-x_{867}^{2}-x_{868}-x_{868}^{2}\\[3pt]
T_{483}={}&x_{362}-x_{869}^{2}-x_{870}^{2}-x_{871}-x_{871}^{2}\\[3pt]
T_{484}={}&-1-x_{35}+x_{363}-x_{872}^{2}-x_{873}^{2}-x_{874}-x_{874}^{2}\\[3pt]
T_{485}={}&x_{368}-x_{875}^{2}-x_{876}^{2}-x_{877}-x_{877}^{2}\\[3pt]
T_{486}={}&x_{369}-x_{878}^{2}-x_{879}^{2}-x_{880}-x_{880}^{2}\\[3pt]
T_{487}={}&x_{370}-x_{881}^{2}-x_{882}^{2}-x_{883}-x_{883}^{2}\\[3pt]
T_{488}={}&x_{371}-x_{884}^{2}-x_{885}^{2}-x_{886}-x_{886}^{2}\\[3pt]
T_{489}={}&-1-x_{1}+x_{372}-x_{887}^{2}-x_{888}^{2}-x_{889}-x_{889}^{2}\\[3pt]
T_{490}={}&x_{377}-x_{890}^{2}-x_{891}^{2}-x_{892}-x_{892}^{2}\\[3pt]
T_{491}={}&x_{378}-x_{893}^{2}-x_{894}^{2}-x_{895}-x_{895}^{2}\\[3pt]
T_{492}={}&x_{379}-x_{896}^{2}-x_{897}^{2}-x_{898}-x_{898}^{2}\\[3pt]
T_{493}={}&x_{380}-x_{899}^{2}-x_{900}^{2}-x_{901}-x_{901}^{2}\\[3pt]
T_{494}={}&-1-x_{36}+x_{381}-x_{902}^{2}-x_{903}^{2}-x_{904}-x_{904}^{2}\\[3pt]
T_{495}={}&x_{386}-x_{905}^{2}-x_{906}^{2}-x_{907}-x_{907}^{2}\\[3pt]
T_{496}={}&x_{387}-x_{908}^{2}-x_{909}^{2}-x_{910}-x_{910}^{2}\\[3pt]
T_{497}={}&x_{388}-x_{911}^{2}-x_{912}^{2}-x_{913}-x_{913}^{2}\\[3pt]
T_{498}={}&x_{403}-x_{914}^{2}-x_{915}^{2}-x_{916}-x_{916}^{2}\\[3pt]
\end{align*}
\endgroup

\end{document}
